\documentclass[lualatex,a4paper,10pt]{jlreq}
\usepackage{../../../templates/survey/jgaisurvey}
\usepackage{needspace}
\usepackage{booktabs}
\usepackage{tabularx}
\usepackage{array}
\usepackage{multicol}
\usepackage{xurl}
\addbibresource{references.bib}

\setlength{\columnsep}{1.45em}
\setlength{\multicolsep}{0.65em}

\surveysetup
  {SP-beyond-text-2026}
  {Japanese Generative AI Technical Survey Special}
  {Beyond Text――画像・音声・音楽・映像生成AIの技術史 / Thematic survey as of 2026-09-23}
  {Open history, evidence observed through 2026-09-23 UTC}
\surveyeditiondescriptor{Thematic Longform Technical Survey}

\surveycoverstory
  {Beyond Text}
  {画素・波形・フレーム・動き・同期音を生成可能にしたのは、表現・目的・条件づけ・制御・編集・時間構造・実行・評価の選択の進化である。有名製品の目録ではなく、何が詰まり、何が変わり、何が残ったかを追う技術史。}
  {Representation \quad / \quad Paradigms \quad / \quad Conditioning \quad / \quad Control \quad / \quad Editing \quad / \quad Speech \quad / \quad Music \quad / \quad Video \quad / \quad Temporal \quad / \quad Runtime \quad / \quad Evaluation \quad / \quad Convergence \quad / \quad Capstones \quad / \quad Reception}

\begin{document}
\surveycover
\clearpage

\section*{本巻の問いと読み方}
本巻の問いは、\textbf{高次元連続メディア（画素・波形・フレーム・動き・同期音）が、表現・生成目的・条件づけ・制御と編集・時間構造・サンプリングと実行・評価の選択の進化によって、いかに生成可能になったか}である。対象は画像・音声・音楽・映像の四つの系譜であり、各モダリティで「何が隘路で、何が変わり、何が改善し、何が悪化し、何が継承されたか」を追う。本巻は製品の順位付けでも、特定モデルの推奨でも、条件の異なる数値の横断比較でもない。\textbf{短縮（表現）が何を可能にし、目的と骨格とサンプリングのどれが何を決めたか}を区別して読む技術史である。

重要なのは、ここで三つの混同を避けることである。第一に、\textbf{骨格と目的とサンプリング}を混同しない。U-NetかTransformerかは骨格、自己回帰か敵対か拡散かflowかは目的、DDIMやsolverや蒸留は推論手続きであり、別の層の選択である。第二に、\textbf{生成と編集}を混同しない。新規生成が白紙から試料を引くのに対し、編集は保存対象と改変領域の定義を要する別問題である。第三に、\textbf{標本の鮮明さと長時間の能力}を混同しない。局所忠実度、短距離整合、同一性永続、同期、継続、編集保存、多回合振る舞いは別の軸であり、一つの得点では語れない。第1節で結ぶ表現の契約と第2節の三層分離が、本巻全体の共通言語になる。

\subsection*{確認範囲と主張境界}
本巻の技術的主張は、受理された証拠の範囲に縛られる。139件の典拠のうち126件は本文まで読み込めたが、8件は要旨・断片水準（PARTIAL）、5件は本体取得が阻まれ保留（NEEDS\_MORE/HOLD）である。概要水準の主張は概要水準として扱い、本文で裏付けられた言語では記述しない。提供元による速度・価格・精度の数値は、測定条件付きの主張として扱い、独立検証のない数値は事実として引用しない。異なる条件で測られた性能値を直接並べて優劣を論じることはしない。閉鎖製品は能力・workflow・可用性・lifecycleの文脈証拠としてのみ扱い、振る舞いからの機構推論は行わない。Xの27件の観察は一つの有界な受容記録であり、受容・配備・失敗の証拠に限定し、技術的事実の根拠にはしない。

\subsection*{本巻の限界の要約}
本巻の限界をあらかじめ要約する。EDM設計空間のablationは要旨限定であり本文言語では記述しない\autocite{btd022}。VALL-E 2の機構は消費できたが正確な同等余白や高速化の数値は表欠落のため不確定である\autocite{btd059}。FID原典の全文は未消費であり要旨と相互記述の範囲に留まる\autocite{btd083}。LibriSpeechの分割・整列・基準誤り率は断片水準のみである\autocite{btd089}。C2PA仕様本体は案内面のみで表明・断言・署名・結合の詳細は未消費である\autocite{btd098}。Imagen所在は製品ハブへ転送され能力記述を持たないためlifecycleの事実のみを記録する\autocite{btd106}。Moshiは枠組みのみを用い評価章の数値は確定しない\autocite{btd134}。Movie Genは頁要旨の段階のみを用い手法・評価部は未確立である\autocite{btd076}。以下は本体が阻まれ保留（HOLD）とし、機構の主張は兄弟典拠に委ねる。Stable Diffusion公開実装の版バインドは未解決である\autocite{btd024}。Phenakiは壁の向こうにありmasked-prior比較は未解決である\autocite{btd072}。ITU-R BS.1534勧告文は取得できず手順の細部は未確立である\autocite{btd091}。Kling 3.0 IR頁は取得できず報道で代替しない\autocite{btd120}。Wanハブ頁は未取得であり開放線の裏付けはリポジトリ記録に委ねる\autocite{btd125}。これらを隠して記事を整えないことが、本巻の編集方針である。

\subsection*{この巻で比較するもの}
本巻が横断的に比較するのは、モデルの絶対順位ではなく、\textbf{どの設計問題が第一級の競争軸になったか}である。具体的には、連続潜在・離散符号・残差量子化・意味音響階層・時空間圧縮という表現の選択、自己回帰・敵対・拡散・flowという目的の分岐、文章条件と案内の分離、空間・参照・同一性の制御、新規生成と編集の分離、音声・音楽・映像の固有の時間構造、段数・遅延・記憶・公開性の実行経済、分布・構成・選好・多面という評価の物差し、そして単一統合と専門家連携の収束問題を追う。2025--2026年の現行群は順位付けではなく能力・workflow・lifecycleの証拠として読み、最終節では座標そのものを点検する。

\subsection*{座標としての通史}
本巻の通史を座標で読む。起点は短縮である。画素・波形・フレームを扱える系列に短縮できなければ、目的も条件づけも始まらない。次に目的と骨格とサンプリングの分離が来る。同じ拡散でも骨格が違えば振る舞いが変わり、同じ骨格でも目的が違えば学ぶものが変わり、同じ学習でもサンプリングが違えば配備の可否が変わる。条件づけと案内、制御と参照、生成と編集がその上に積み、音声・音楽・映像の各系譜が固有の時間と実行と評価を持つ。層は断絶ではなく積み重ねであり、後の層は前の層を置き換えるのではなく、前提として使う。表現の短縮は実行の設計に、目的の分岐は評価の物差しに、制御と編集の分離は現場のworkflowに接続する。座標は、層の重なりを設計の言葉で言い直したものである。

\begin{center}
\small
\begin{tabularx}{\textwidth}{@{}p{0.20\textwidth}p{0.32\textwidth}X@{}}
\toprule
層 & 代表的手法 & 中心の問い \\
\midrule
表現の短縮 & VAE・VQ・RVQ符号・意味音響階層・時空間圧縮 & 何をどれだけ短くし何を失うか \\
目的と骨格とサンプリング & AR・GAN・拡散・flow・DiT・DDIM・蒸留 & どれが何を決めるか \\
条件と制御と編集 & CFG・ControlNet・参照・指示編集 & 何を保ち何を変えるか \\
時間と実行と評価 & 同期・物理・段数・遅延・多面評価 & 続くか動くか測れるか \\
収束と現場 & 統合・連携・capstone・受容 & どこへ向かい何が残るか \\
\bottomrule
\end{tabularx}
\end{center}

\subsection*{記号と表記の約束}
本巻では、生成の三層を骨格（U-Net・Transformer・DiT等）、目的（自己回帰・敵対・拡散・score・flow等）、推論手続き（DDIM・solver・蒸留・一致性等）と呼ぶ。潜在は符号化器で短縮された表現、符号帳は離散量子化の表、残差量子化は層を重ねる音声符号、意味符号は言語的継続を担う符号、音響符号は波形再構成を担う符号を指す。案内は推論時の外挿操作、適合は分布の近さ、整合は条件への一致、忠実度は保存対象の保たれ方を指す。提供元による数値は提供元主張、独立検証のない数値は未確認として扱い、本文で断定的に引用しない。技術的事実を確定する典拠は一次資料・公式文書・実装記録・公開手順書を基本とし、受容記録やコミュニティ資料を引用する場合は配置依存の観察・再現の手がかりに限定し、技術的事実の根拠には用いない。

\subsection*{引用の約束}
本巻の引用は、主張の出所を示すためのものであり、権威の飾りではない。一次資料の引用は、本文まで読んだ範囲でのみ構造と条件を語り、要旨だけ読んだ範囲では見出しの事実に留める。公式文書と手順書の引用は、告知と配置の範囲に留め、他構成への転用をしない。実装記録の引用は、存在と手順の証拠に留め、正確性の裏書きにはしない。提供元の数値の引用は、測定条件付きの主張として隔離し、事実としては引用しない。受容記録の引用は、配置依存の観察として扱い、技術的事実の根拠にはしない。引用の後に残る未解決は、境界として明示する。引用は主張を閉じるためではなく、読み手が確かめに戻れるようにするためにある。

\subsection*{読者への道案内}
時間の限られた読者は、本節の層表と各節末の境界記述だけを追っても、本巻の主張の骨格をつかめる。各節の用語は本節の三層分離に準拠する。骨格と目的とサンプリング、生成と編集、標本鮮明さと長時間能力の三つの分離が、各節で繰り返し使われる合言葉である。表現と目的の理論に関心のある読者は、第1--2節の系譜表を中心に読むとよい。条件と制御と編集の設計に関心のある読者は、第3--5節の対照表が近道になる。モダリティの系譜に関心のある読者は、第6--8節の lineage 表、運用と測定の実務に関心のある読者は、第9--11節の作法表、現場と収束に関心のある読者は、第12--14節の対照表が近道になる。結びは層の要点と指針と診断の実践を圧縮したものであり、本文の詳述に戻る際の索引として使える。

\begin{center}
\small
\begin{tabularx}{\textwidth}{@{}p{0.14\textwidth}p{0.34\textwidth}X@{}}
\toprule
節 & 問い & 要点 \\
\midrule
第1節 & 何を短くするのか & 表現の短縮と忠実度・編集可能性の交換 \\
第2節 & どれが何を決めるのか & 目的・骨格・サンプリングの三層分離 \\
第3節 & 何に従うのか & 条件づけと案内の分離 \\
第4節 & 何を保つのか & 空間・参照・同一性の制御 \\
第5節 & 何を変え何を残すのか & 生成と編集の分離 \\
第6節 & 声はどう作られるか & 音声・声の固有系譜 \\
第7節 & 音楽はどう続くか & コーデック言語モデルと潜在拡散の二系統 \\
第8節 & 映像はどう動くか & 時間的一貫性の系譜 \\
第9節 & 続くことと鮮明さの違い & 長時間・同期・編集の評価軸 \\
第10節 & 動かせるか & 段数・遅延・記憶・公開性 \\
第11節 & 数字をどう読むのか & 指標の非交換性 \\
第12節 & どこへ向かうのか & 統合か連携かの開かれた問い \\
第13節 & 2026年の到達点は何か & 能力・workflow・lifecycleの証拠 \\
第14節 & 現場は何を見たか & 27件の受容・配備・失敗の記録 \\
結び & 何を持ち帰るのか & 層の要点と指針と診断の実践 \\
\bottomrule
\end{tabularx}
\end{center}

\setcounter{tocdepth}{2}
\tableofcontents
\clearpage

\section{表現と圧縮：生成可能にする短縮の歴史}
\sectionkicker{REPRESENTATION}
\label{sec:representation}

\begin{multicols}{2}
\raggedcolumns
画像の画素、音声の波形、映像のフレーム列をそのまま系列として扱おうとすれば、長さはたちまち扱いきれない領域へ膨らむ。二百五十六画素四方の画像だけでも六万五千超の画素値が並び、秒間数万標本の音声は一秒で数万点の波形となり、十六フレームの映像はその十数倍の画素を抱える。生成模型がこれらを直接の予測対象にすれば、記憶と計算の両面で系列長が律速となり、長距離の依存関係を捉える前に学習が滞る。本節がたどるのは、この長大系列の問題に対して、短縮そのものを第一級の設計対象に据えた系譜である。連続潜在への写像、離散符号帳への量子化、残差量子化による加法的精緻化、意味と音響の階層分離、時空間方向の映像符号化という五つの転換が、圧縮と忠実度と編集可能性の三つの軸の上で何を削り何を残したかを整理する\autocite{btd001,btd002}。

本節の読み方として三つの分離を固定しておきたい。第一に、アーキテクチャと目的関数と推論時サンプリングは別の欄である。符号器と復号器の形が変わったのか、学習の損失が変わったのか、生成時の抽出手順が変わったのかを混同しない。第二に、圧縮率やトークン速度の主張は、報告された構成に結びつけて読む。同じ符号器でも層数や帯域や解像度が変われば短縮の度合いは変わるため、条件を離れた一般論にはしない\autocite{btd006,btd008}。第三に、意味の一貫性と波形や画素の忠実度は別の物差しである。言葉として筋が通ることと、音や絵として精細であることは、異なる符号が担う。本節の後半で見る意味トークンと音響トークンの分離は、この区別を表現の水準で制度化したものにほかならない\autocite{btd009,btd010}。

時代の流れをあらかじめ見渡す。起点は連続潜在変数を持つ有向生成模型であり、推論の償却化が学習可能領域を広げた\autocite{btd001}。続いて離散符号帳への量子化が導入され、音響構造を保つ離散表現の可能性が示された\autocite{btd002}。画像では階層化と知覚的符号化が二つの道を開き、高解像度への到達と短縮の両立が進んだ\autocite{btd003,btd004}。大規模な文条件づけとの結合は、画像と文を単一ストリームに載せる離散化によって先駆けられた\autocite{btd005}。音声では残差量子化が低速度と端末遅延の両立を担い\autocite{btd006,btd007,btd008}、意味と音響の階層と映像の時空間符号が本節の到達点となる\autocite{btd009,btd010,btd012}。各時代の改善は条件付きの報告であり、代償と境界を伴う。以下の各項では、隘路と機構と改善と代償と継承の五点で各転換を記述する。
\end{multicols}

\subsection{連続潜在から離散符号へ：推論の償却化と量子化}
連続潜在変数を持つ有向生成模型では、データ点ごとに扱いにくい事後分布の推論が必要となり、大規模データでの学習が滞っていたことが出発点である。変分オートエンコーダはこの隘路に対し、推論を個別最適化から償却化された認識模型へ置き換える転換を示した\autocite{btd001}。アーキテクチャとしては多層パーセプトロンによる確率的符号器と復号器を置き、標準正規分布を事前分布とする構成が採られた。目的関数は証拠下界であり、解析的な発散項と再パラメータ化による勾配推定が機構の中核にある。推論時サンプリングは認識模型を経由した祖先抽出で実行され、マルコフ連鎖を必要としない点が推論側の特徴である\autocite{btd001}。報告された改善は下界の向上と二次元潜在多様体の可視化に結びつき、連続表現が学習可能領域を広げたことが示された。一方で、消費された本文が示すのは平均場近似を用いないガウス事後の一例に限られ、離散潜在や大域的母数への変分は扱われないという境界がある。この限定は弱さの告白ではなく、読みの作法である。概要水準の主張を本文水準として語らないことが、本特集全体の約束だからである。

連続潜在の継承課題は、強力な復号器のもとで潜在変数が無視される事後崩壊であった。離散符号帳による量子化は、この隘路への継承策として位置づく。ベクトル量子化変分オートエンコーダは符号帳への最近傍探索とストレートスルー勾配を用い、一様カテゴリ事前分布のもとで学習される\autocite{btd002}。アーキテクチャと目的関数と推論時サンプリングの分離に注意すれば、ここでの転換は表現の離散化であり、事前分布の同時学習ではない。表現学習の段階では一様事前分布が置かれ、事前分布の学習は後段の課題として残された。音声では約二十五ヘルツの符号列で四九・三パーセントの音素対応が報告され、離散表現が音響構造を保つことが示された\autocite{btd002}。ただし画素損失に伴うぼけと符号帳崩壊の懸念が残り、短縮は得たものの忠実度の上限と編集の粒度は後継の階層化と知覚的符号化へ引き継がれる。圧縮と忠実度と編集可能性の三軸で見れば、第一世代の離散化は圧縮の軸を前進させ、忠実度の上限を次の時代の課題として残したことになる。

\subsection{階層と知覚：高解像度への二つの道}
単一スケールの離散符号だけでは大きな画像の高忠実度生成に届かないことが次の隘路となった。表現の側では二層の離散写像が導入され、二百五十六画素では上位三十二掛ける三十二と下位六十四掛ける六十四を条件づけて組み合わせ、高解像ではさらに層を重ねて千画素級へ伸ばす設計が採られた\autocite{btd003}。別の系統では知覚損失とパッチ識別器を備えた畳み込み符号帳が導入され、二百五十六画素を十六掛ける十六の指標列へ圧縮する短縮が実現した\autocite{btd004}。二つの道は対立ではなく、階層化による扱い域の拡大と、知覚的符号化による短縮の徹底という異なる方向への前進である。アーキテクチャとしては、前者でフィードフォワード符号器と復号器に自己回帰 prior を重ね、指数移動平均で符号帳を更新する構成、後者で符号器復号器と潜在Transformerを分離し滑動窓で扱う構成が選ばれた\autocite{btd003,btd004}。

目的関数の水準では、二段階構成への整理が要点である。再構成と符号帳およびコミットメント損失に、凍結符号上の prior の負対数尤度を加える形へ整理され、敵対項の適応重みが後者の特徴として加わる\autocite{btd004}。推論時サンプリングは上位から下位への条件付き生成と単一順伝播復号、あるいは指標列の自己回帰予測と温度付き抽出によって行われる\autocite{btd003}。改善としては棄却後の分類精度と人手忠実度の向上、画素空間に対する十倍超の高速化と分布距離の低下が条件付きで報告された\autocite{btd003,btd004}。ここでの条件付きという留保は重要である。潜在空間を跨ぐ尤度の比較はできないこと、棄却が多様性を削ること、窓の空間不変性が必要で文脈が固定されることが残るからである。符号器の質が生成上限を決めるという知見は、後の潜在拡散の第一段階へ継承される。生成の良し悪しを論じる前に符号器の忠実度を問う習慣は、この時代に定まった。

\begin{center}
\small
\begin{tabularx}{\textwidth}{@{}p{0.22\textwidth}p{0.36\textwidth}X@{}}
\toprule
転換 & 短縮の形 & 残された上限 \\
\midrule
連続潜在の償却化 & 推論の個別最適化を認識模型へ & 離散潜在と大域母数は範囲外 \\
単一符号帳の量子化 & 音響構造を保つ離散列 & 画素損失のぼけと符号帳崩壊の懸念 \\
階層離散写像 & 上位と下位の条件づけで千画素級へ & 層数と prior 学習の重さ \\
知覚的符号帳 & 二百五十六画素を十六掛ける十六指標へ & 復号上限が生成上限を決める \\
\bottomrule
\end{tabularx}
\end{center}

\subsection{文と画像の単一列：大規模条件づけの前の短縮}
離散表現の学習が無条件中心であったため、大規模な文条件づけと結びつかないことが隘路となった。表現の側では二百五十六画素を三十二掛ける三十二の千二十四トークンへ離散化し、語彙八千超の画像符号と二百五十六トークン以内のバイト対符号化文を単一ストリームとして連結する短縮が採られた\autocite{btd005}。この連結は、文と画像を別々の条件器でなく同一系列の要素として扱う発想であり、後の大規模文条件づけの先駆けとなった。アーキテクチャとしては畳み込み符号器と復号器にガンベル緩和の隘路を設け、六十層級の疎な復号器専用Transformerで行列と畳み込み注意を束ね、分散学習で扱う構成が選ばれた\autocite{btd005}。

目的関数は結合証拠下界であり、事前分布の交差エントロピーに文と画像で異なる重みを与え、分布外の文と画像の対を jointly に学ぶ点が特徴である\autocite{btd005}。生成手順は自己回帰生成ののち、対照模型による五百超の候補からの再順位付けで締める手順となり、三十数件を超えると利得が頭打ちになることが報告された。改善としては、条件付きで人手評価の写実性と文整合が高い割合を示し、学習重複を除いても結果が変わらないことが確認された\autocite{btd005}。代償として、強い圧縮は高周波細部を落とし、特定分野で専門模型との開きが残り、再順位付けへの依存や可変束縛と文字描画の不安定さが残る。離散トークンの書換え粒度と再構成誤差の分離も残り、潜在局所編集への継承課題として記録される。なお、離散符号帳そのものの基礎については前項の量子化の知見が前提となっており、符号帳学習と prior 学習の分離という考えはここでも保たれている\autocite{btd002}。

\subsection{音の残差量子化：低速度と遅延の両立}
単一符号器単一速度の設計では、音声と音楽と一般音響を低ビット速度で等しく扱い、端末遅延に収めることが難しいことが隘路となった。表現の側では因果畳み込み埋め込みを三百二十倍に間引き、残差ベクトル量子化で層数を増減できる加法的精緻化へ進んだ\autocite{btd006}。後継では複数ストリームの残差量子化と帯域別識別器と任意帯域対応が加わり\autocite{btd007}、さらに後継では蛇行活性化と分解正規化量子化と層脱落で約九十一倍の圧縮と可変速度が報告された\autocite{btd008}。三世代の継承関係は、層を足し引きできる加法的設計という一本の考えで貫かれている。速度を変えたいときに符号器を作り直すのではなく、使う層の数を選ぶことで目標速度を得る。この可変速度の発想が、音響符号を持続可能な生成部品にした。

アーキテクチャとしては因果符号器復号器に雑音除去条件や回帰型層、帯域別多解像度識別器や複素スペクトル識別器が組み合わされた\autocite{btd006,btd007}。目的関数は時間領域損失とメル再構成に敵対損失と特徴一致と符号帳損失を束ね、勾配均衡や層脱落で可変速度に対応する点が共通する\autocite{btd007,btd008}。推論は因果ストリーミングで行われ、層数選択で目標速度を得て、初期遅延十ミリ秒台と実時間比一桁が条件付きで報告された\autocite{btd006}。改善としては低速度で既存符号器を上回る聴取評価や、エントロピー符号化による二割超の削減、全帯域での既存高速度符号器に対する優位が条件付きで示された\autocite{btd007,btd008}。代償として、聴取評価は厳密方式ではなく群衆方式であること、一定速度動作であること、音楽が最も難しいことが残る。高帯域と符号化の両立では実時間を割る点と参照水準との開きも記録される。条件別の聴取得点と客観指標の対応、帯域別の識別器寄与、層数と速度の交換が整理され、残差量子化の考えは言語模型型音響生成へ継承される。速度可変の符号設計は音楽生成の制御適合へも接続する。

\subsection{意味と音響の分離、映像の時空間短縮}
音響符号だけでは意味の一貫性が崩れて呟きとなり、意味符号だけでは波形に戻せないことが隘路となった。表現の側では自己教師あり音響符号器の中間層を千語彙で量子化した意味トークンと、残差量子化音響トークンの粗密二層を組み合わせる階層が採られた\autocite{btd009}。意味の一貫性と言語としての筋道を担う層と、波形の質感を担う層を分けたことが要点である。映像では三次元量子化で十六フレームを千二十四トークンへ短縮し、二次元符号の膨張初期化で学習する設計が選ばれ、多課題マスク学習で補完と予測を束ねた\autocite{btd010}。さらに高圧縮映像変分符号器では時間高さ幅方向に六十四倍の短縮とパッチ化を重ね、消費者用図形処理装置での高精細生成へ接続した\autocite{btd012}。

アーキテクチャとしては、音響で三段階の復号器専用Transformer、映像で課題 prompt と類別を伴う双方向Transformer、公開系で高雑音と低雑音の二専門家混合が置かれた\autocite{btd009,btd010}。目的関数は段階別次トークン尤度や多課題交差エントロピーであり、公開系の学習損失は文書開示にないという境界がある\autocite{btd012}。推論時サンプリングは意味拡張から音響拡張への段階継続、あるいは十数段階の非自己回帰復号で行われ、秒間数十フレームの実効速度が条件付きで報告された\autocite{btd009,btd010}。改善としては意味保持と話者保持の両立や映像の分布距離の大幅低下が示された\autocite{btd009,btd010}。代償として固有名詞や雑音での誤りと固定長学習が残る。符号忠実度が生成上限となる点と文書開示の不足も記録される。意味と音響の分離は音楽の段階生成へ継承され、映像短縮は可変長生成の検証課題として残る。編集可能性の検証が後継へ引き継がれる。なお高圧縮映像符号の効率表の一部数値は画像条件に結びつくものであり、映像条件一般への転用はできないという留保が伴う\autocite{btd012}。

\subsection{三軸のまとめ：何を削り何を残したか}
五つの転換を圧縮と忠実度と編集可能性の三軸で振り返る。連続潜在の償却化は学習可能領域を広げたが、短縮そのものはもたらさず、系列長の問題は残した\autocite{btd001}。単一符号帳の量子化は短縮をもたらし、音響構造の保持を示したが、画素損失のぼけと符号帳崩壊の懸念を残した\autocite{btd002}。階層離散写像は扱える解像度を広げて千画素級へ届かせ、知覚的符号帳は十六掛ける十六指標への徹底した短縮を示したが、いずれも復号上限が生成上限を決めるという制約を残した\autocite{btd003,btd004}。単一列の連結は文条件づけへの道を開いたが、高周波細部の脱落と再順位付け依存を残した\autocite{btd005}。残差量子化は層数選択による可変速度で低速度と遅延を両立させたが、聴取評価の方式と音楽の難しさの留保を伴った\autocite{btd006,btd007,btd008}。意味と音響の分離は一貫性と質感の分担を制度化し、映像の時空間符号は十六フレームの千二十四トークン化と六十四倍短縮で生成の土台を築いたが、固有名詞と雑音での誤り、固定長学習、文書開示の不足を残した\autocite{btd009,btd010,btd012}。

編集可能性の軸は、五つの転換を通じて一貫して未成熟であった。離散トークンの書換え粒度と再構成誤差の分離は記録されるに留まり、潜在局所編集の検証は後継へ引き継がれる\autocite{btd005}。意味と音響の階層も編集の検証は後継の課題として残した\autocite{btd009}。映像短縮は可変長生成の検証課題を残した\autocite{btd010}。この残し方が次章以降への橋渡しとなる。表現が短縮を与え、目的関数と推論時サンプリングがその短縮の上で競い、条件づけと制御が短縮された列に働きかける。読者は本節の符号器を、後の潜在拡散の第一段階\autocite{btd004}、音響言語模型の符号\autocite{btd006}、映像生成の時空間土台\autocite{btd012}として再会することになる。短縮は終わりではなく、生成の始まりである。

\subsection{本節の読解上の境界}
本節の主張は報告条件に結びつく。圧縮率やトークン速度の数値は、述べた構成のもとでのものであり、構成を離れて比べない\autocite{btd006,btd008}。変分オートエンコーダの記述は平均場近似を用いないガウス事後の一例に限られ、離散潜在や大域的母数への変分は範囲外である\autocite{btd001}。下流の圧縮と忠実度の主張は第一段階符号器の裏付けを要し、単独論文だけでは閉じない\autocite{btd004}。表現学習時の事前分布は一様であり、事前分布との同時学習は行われていない。画素損失のぼけは本文中で記録されている\autocite{btd002}。符号帳崩壊と使用率の前線は後継資料の検証を要する。音声の聴取評価は厳密方式ではなく群衆方式であり、一定速度動作と音楽の難しさが伴う\autocite{btd006}。速度と品質の前線は整合条件での比較を要する\autocite{btd007}。識別器の抑え込みは確率的更新で扱われ、高帯域と符号化の併用は実時間を割り、立体音響は識別器調整による対応に留まる\autocite{btd008}。意味と音響の分離の評価は混合型資料の裏付けを要し\autocite{btd009}、音響のみの模型化の呟き化と意味のみの符号化の不可逆性、固有名詞と背景雑音での誤り、鍵盤継続の限定は本文記録の境界である\autocite{btd009}。長い結構の評価は後継の音楽資料に委ねられる。映像符号の忠実度が生成上限を決め、十六フレーム固定学習と素朴な覆い外しの失敗が動機づけの前提となる\autocite{btd010}。覆い prior と拡散 prior の比較は後継の対象であり、符号器を跨ぐ潜在尤度の比較は成立せず、分布距離と分類得点は符号帳ぼけの影響を受け、 prior 学習の重さも記録される\autocite{btd003}。階層符号の費用と適用範囲の比較は後継比較資料を要し、滑動窓生成は空間不変性と条件づけを前提とし、Transformer文脈は十六掛ける十六に固定される\autocite{btd004}。強い圧縮の高周波細部の脱落と特定分野の開き、再順位付け依存、可変束縛と文字描画の不安定さは本文記録の境界である\autocite{btd005}。離散変分符号器の上限と後継比較は後継資料の対象であり、参照水準との開きと環境音や特定音色での弱さ、均衡ある全帯域抽出の要請が残る\autocite{btd007}。生成下流での利得は採用側の証拠を要し、符号器指標だけでは閉じない。公開系の開示は文書水準に留まり、変分再構成の数値や学習損失の内訳はなく、効率表の一部は画像条件に結びつく\autocite{btd012}。公開系列の線引きは二・二までであり、それ以降は応用接続のみである。

\section{生成パラダイムと目的関数：自己回帰・敵対・拡散・フロー}
\sectionkicker{PARADIGMS}
\label{sec:paradigms}

\begin{multicols}{2}
\raggedcolumns
何をどう学ぶのかという問いは、模型の形と学習の損失と生成の手順を分けて初めて見通せる。本節がたどるのは、敵対的生成と自己回帰画素から出発し、雑音除去拡散と暗黙的短縮、得点の連続時間 framing、潜在拡散、Transformer denoiser、Flow Matching（フローマッチング）と整流とConsistency Model（整合性モデル）への分岐に至る系譜である\autocite{btd013,btd019,btd023}。通読の約束として三層の区別を固定する。第一に、U網かTransformerかはアーキテクチャの欄、拡散か流れかは目的関数の欄、解法や蒸留は推論時サンプリングの欄である。第二に、条件の異なる数値の横断順位付けはしない。解像度や予算や評価手順が違えば、分布距離の大小はそのまま優劣にならない。第三に、要旨のみの報告は本文未検証として限定する。とくに設計空間整理の一件は要旨消費に留まり、前処理や日程や解法の内訳と backbone 対目的の分離や他 modality への主張は本文で確かめられていない\autocite{btd022}。

時代の流れをあらかじめ見渡す。起点は尤度に基づく学習が扱いやすい密度や連鎖を要求した時代であり、逆伝播に親和的な暗黙的生成の不在が隘路となった。敵対的生成は生成器の標本だけからなる暗黙分布を置き、識別器との競合で分布一致を目指す転換を示した\autocite{btd013}。畳み込みの安定化指針が加わり\autocite{btd014}、様式生成器が潜在の分離と品質改訂を進めた\autocite{btd015,btd016}。並行して画素再帰網とImage Transformerが完全自己回帰分解で同時分布の捕捉を示した\autocite{btd017,btd018}。拡散模型が良質標本を示せず品質主導が敵対側にあった時代に、雑音除去拡散と大規模化と暗黙的短縮が転換をもたらし\autocite{btd019,btd020,btd030}、連続時間の確率微分方程式 framing が両者の統一を与えた\autocite{btd021}。画素拡散の費用が次の隘路となり、知覚学習済み符号器で一度圧縮して潜在で拡散する転換が置かれた\autocite{btd023}。U網必須の仮定が崩れ、Diffusion Transformer (DiT)とScalable Interpolant Transformer (SiT)が backbone 対目的の分離を示し\autocite{btd025,btd026}、Flow Matching（フローマッチング）とRectified Flow（整流フロー）と整合性モデルが高速抽出への分岐を開いた\autocite{btd027,btd028,btd029}。各転換の改善は条件付きの報告であり、代償と境界を伴う。各転換は表現と骨格と目的と確率経路と抽出と短縮の六欄で読み分け、骨格の交代を目的の交代と取り違えない。
\end{multicols}

\subsection{六軸分離の読み方：表現・骨格・目的・経路・抽出・短縮}
表現は分布を扱う空間であり、画素そのものか圧縮潜在か雑音との補間かが問われる\autocite{btd017,btd023,btd026}。骨格は模型の形であり、多層パーセプトロンか畳み込み指針か時刻条件U網かTransformerかが問われる\autocite{btd013,btd014,btd019,btd025}。目的は最小化する損失であり、競合損失か離散負対数尤度か雑音予測か速度回帰か整合性かが問われる\autocite{btd013,btd017,btd019,btd027,btd029}。確率経路は前向き過程と逆過程の定式化であり、固定ガウス前向きと学習逆過程か、周辺保存の非マルコフ前向きか、連続時間確率微分方程式か、線形補間かが問われる\autocite{btd019,btd020,btd021,btd028}。抽出は生成に要する段階であり、単一順伝播か画素逐次か千段階祖先抽出か適応積分か暗黙的軌道かが問われる\autocite{btd013,btd017,btd019,btd021,btd023}。短縮は軌道をどう間引くかであり、部分軌道か十から二百段階の暗黙的軌道か直線軌道の少数段階か自己整合性かが問われる\autocite{btd020,btd023,btd027,btd029}。

\begin{center}
\small
\begin{tabularx}{\textwidth}{@{}p{0.24\textwidth}p{0.34\textwidth}X@{}}
\toprule
欄 & 問うこと & 例 \\
\midrule
表現 & どの空間で分布を扱うか & 画素多値か圧縮潜在か補間か \\
骨格 & 模型の形は何か & 畳み込み指針か時刻条件U網かTransformerか \\
目的 & 何を最小化したか & 競合か雑音予測か速度回帰か整合性か \\
確率経路 & 前向きと逆過程の定式化は何か & 固定ガウスか非マルコフか連続時間か補間か \\
抽出と短縮 & 生成に何段階要しどう間引くか & 単一順伝播か画素逐次か千段階か直線化か \\
\bottomrule
\end{tabularx}
\end{center}

\subsection{敵対的生成の系譜：暗黙分布から様式分離へ}
尤度に基づく学習が扱いやすい密度や連鎖を要求したことが、逆伝播に親和的な暗黙的生成の不在という隘路となった。敵対的生成はこの隘路に対し、生成器の標本だけからなる暗黙分布を置き、識別器との競合で分布一致を目指す転換を示した\autocite{btd013}。アーキテクチャとしては多層パーセプトロンの生成器と識別器を交互に更新する構成が出発点であり、後に歩幅畳み込み識別器と分数歩幅生成器、集合同期や生成出力への正規化除外などの設計指針が加わった\autocite{btd014}。目的関数は競合損失であり、一致時に識別が偶然水準となることが理論上の帰着として述べられ、生成器の非飽和損失が実用上の工夫とされた\autocite{btd013}。サンプリングは潜在からの単一順伝播であり、連鎖を必要としない点が機構上の利点とされた\autocite{btd013}。改善としては分布類似の代理評価や、下流分類のプローブで同学構成の教師ありを上回ることが条件付きで報告された\autocite{btd013,btd014}。代償として、明示密度を持たないこと、識別器の同期が崩れると様式崩壊を招くこと、高次元での代理評価の分散が大きいことが残る。長期学習での濾波崩壊と対数尤度の不在も記録される。安定な畳み込み構成は後の音声 vocoder や符号器の敵対項へ継承され、単独の画像系譜としては後の拡散系へ座を譲る。現在の役割は単独の到達点ではなく、 vocoder や符号器や後学習の混合利用にあるという位置づけが境界として残る。骨格側は歩幅畳み込み識別器と分数歩幅生成器、集合同期や生成出力への正規化除外、生成器側の整流化線形と双曲線正接と識別器側の漏れ整流化線形、〇・〇〇〇二の更新幅であり、競合損失と非飽和損失の工夫とは欄が異なる\autocite{btd013,btd014}。分類器プローブ八二・八パーセント、千標本条件の誤り率二二・四八対同学構成教師あり二八・八七が条件付きで報告された\autocite{btd014}。

初期の敵対的生成では絡み合った潜在と水滴状や blob 状の作画上の乱れ、位置選好が隘路となった。様式生成器はこの隘路に対し、入力潜在を中間空間へ写像し、層ごとに様式と確率的雑音図で制御する分離を導入した\autocite{btd015}。アーキテクチャとしては八層写像網と学習済み定数からの合成網、層別適応正規化と尺度別混合が中核であり、改訂版では外的ばらつきと雑音の配置を見直し、標準偏差のみ変調と重み復調、跳躍生成器と残差識別器へ整理され、漸進成長を廃した点が異なる\autocite{btd016}。目的関数は二潜在の混合正則化に、遅延付き正則化と経路長正則化を組み合わせ、ヤコビの直交性を促す方向へ進んだ\autocite{btd015,btd016}。推論時サンプリングは尺度別様式混合と打切り、雑音の再抽出による確率的細部制御で行われ、帰属のための射影最適化が伴った\autocite{btd015}。改善としては分布距離と経路長の同時低下が条件付きで報告され、線形分離性の向上が伴った\autocite{btd015,btd016}。代償として、混合正則化が経路長をわずかに歪めること、入力空間が絡んだままであることが残る。分布距離と知覚経路の質感偏りや非構造集合での相反も記録される。中間空間の制御概念は参照適合の先駆けとして位置づき、この制御と品質改訂の系譜は後の参照適合や個人化の制御へ継承される。品質上限自体は後の拡散系の比較対象となる。八層写像網と学習済み定数からの合成網と層別適応正規化と尺度別混合が骨格側、混合正則化が目的側であり、改訂版では雑音配置の見直しと標準偏差のみ変調と重み復調と跳躍生成器と残差識別器への整理が骨格側、遅延付き正則化と経路長正則化の追加が目的側である\autocite{btd015,btd016}。千二十四画素条件で八・〇四から四・四〇へ、経路長四百台から二百台前後への低下と、改訂版で四・四〇から二・八四へ、二百台から百四十台への低下が条件付きで報告された\autocite{btd015,btd016}。

\subsection{完全自己回帰：画素の同時分布を捉える}
潜在推論の扱いにくさを避けて完全な同時分布を扱い、画素間の長距離依存を捉えることが隘路となった。画素再帰網は行優先と色分解による完全自己回帰分解を置き、画素値を多値分類として扱う表現を選んだ\autocite{btd017}。Image Transformerはこの系統をTransformerへ置き換え、画素埋め込みと二次元座標で系列を扱い、遮蔽付き局所多頭注意で計算を抑えた\autocite{btd018}。アーキテクチャとしては行方向や対角方向の再帰層と遮蔽畳み込みの多重解像度版、前者に対して復号器専用Transformerと超解像用符号器復号器が対応し、無条件とクラス条件と超解像の三形態が用意された\autocite{btd017,btd018}。目的関数は離散負対数尤度であり、ビット毎次元の尺度で比較され、位置和の最尤として整理される\autocite{btd017}。学習時は並列に尤度を評価できる一方、生成時は画素を一つずつサンプリングするため費用は画素数に比例して積み上がり、十二層までの再帰や遮蔽畳み込みの多重解像度版もこの逐次性を解消しない\autocite{btd017}。

推論時サンプリングは学習時は並列だが生成時は画素逐次であり、条件付きサンプリングによる塗りつぶしや、温度付き逐次抽出と低解像からの超解像復号が採られた\autocite{btd017,btd018}。改善としては桁規模での尤度向上と超解像の騙し率が条件付きで報告され、画素単位の忠実な分布捕捉が示された\autocite{btd017,btd018}。代償として逐次生成の費用が大きく単一尺度では大域構造が弱いことが残る。六十四画素の単一尺度では大域構造が弱く、規模は図形処理装置の記憶と時間に制約される\autocite{btd017}。局所注意の文脈切詰めと小解像への限定も記録される\autocite{btd018}。条件別の層数と尤度の対応、注意範囲と大域構造の交換、逐次費用と並列学習の対比が整理され、この自己回帰 prior の考えは符号帳 prior や音響言語模型へ継承される。完全同時分布の捕捉は符号帳系列の事前分布設計へ引き継がれ、超解像接続の考えも後継へ残る。規模と記憶時間の律速という費用の物差しは、後の潜在系列の設計にも通じる。七層対角双方向再帰で七九・二〇ナット、三・〇〇ビット毎次元、Image Transformerで二・九〇ビット毎次元が条件付きで報告された\autocite{btd017,btd018}。

\begin{center}
\small
\begin{tabularx}{\textwidth}{@{}p{0.24\textwidth}p{0.34\textwidth}X@{}}
\toprule
欄 & 問うこと & 例 \\
\midrule
アーキテクチャ & 模型の形は何か & U網かTransformerか、畳み込み指針か \\
目的関数 & 何を最小化したか & 競合損失か雑音予測かFlow Matching（フローマッチング）か \\
推論時サンプリング & 生成に何段階要するか & 単一順伝播か千段階抽出か短縮軌道か \\
\bottomrule
\end{tabularx}
\end{center}

\subsection{雑音除去の転換：拡散と短縮抽出}
拡散模型が良質標本を示せず、画像の品質主導が敵対的生成にあったことが隘路となった。雑音除去拡散は固定ガウス前向き過程と学習逆過程を置き、雑音予測の簡易目的が変分下界の小時刻項を抑える機構を示した\autocite{btd019}。別系統の大規模化は百二十八基底路と多解像注意、残差組立てと適応正規化による時刻と類別の条件づけで品質を押し上げ、分類器誘導を組み合わせた\autocite{btd030}。別系統の大規模化と分類器誘導ではImageNet百二十八画素条件で分布距離二・九七対六・〇二と敵対的生成を下回ったことが条件付きで報告された\autocite{btd030}。アーキテクチャの改良と目的関数の簡易化は区別されるべきであり、ここでは画素Transformer類似のU網が時刻埋め込みと共有される点が前者、重みなし雑音予測と混合変分項と学習分散補間が後者である\autocite{btd019,btd030}。前向きは千時刻の線形日程で十万分の一から百分の二までの分散を固定し、逆過程は学習ガウス遷移を置く\autocite{btd019}。雑音予測の形は多雑音水準の雑音除去得点整合との等価性を露わにし、簡易目的は小時刻項を抑える重みなしの形である\autocite{btd019}。Transformer式正弦時刻埋め込みと十六画素単位の自己注意のU網共有が骨格側であり、大規模化の系統の骨格側と混合変分項と学習分散補間の目的側の分離も同様である\autocite{btd019,btd030}。

推論時サンプリングは千段階の祖先抽出が出発点であり、二百五十段階や二十五段階への短縮と、後に非マルコフ前向き過程で周辺を保ったまま軌道を間引く暗黙的手法が導入された\autocite{btd030,btd020}。改善としては簡易目的での識別得点と分布距離の向上が条件付きで報告され\autocite{btd019}、CIFAR-10条件では識別得点九・四六、分布距離三・一七が簡易目的で示された\autocite{btd019}。暗黙的短縮の側でも、CIFAR-10条件で確定的復号の十段階一三・三六から百段階四・一六へと軌道間引きのもとでの分布距離の維持が条件付きで示され、二十から百段階で千段階級の品質に迫る十倍から五十倍の高速化が報告された\autocite{btd020}。代償として尤度は競争的でなく情報の多くが知覚困難な細部に割かれることが残る\autocite{btd019}。千段階抽出の日単位費用と短軌道化での品質低下も記録される。条件別の段階数と分布距離の対応、誘導尺度と適合再現の交換が整理され、この簡易目的と間引き抽出の考えは潜在拡散と蒸留へ継承される。周辺保存の間引き概念は後継の高速抽出へ引き継がれ、誘導併用の設計も残る。なお軌道を短くするほど分布距離が低下する関係や、実装由来の分散の振る舞い、常微分方程式と確率流の開きは少数段階の境界として記録される。暗黙的短縮は学習目的を変えない抽出側の転換である\autocite{btd020}。周辺を保ったまま非マルコフ前向きで軌道を間引き、確率変数で確率性を切り替え、決定論的側で常微分方程式の見方と潜在の符号化を開く\autocite{btd020}。十段階から百段階の部分軌道で千段階級の品質に迫り、十倍から五十倍の短縮が条件付きで示された\autocite{btd020}。

\subsection{条件誘導の分岐：分類器誘導と分類器なし誘導}
品質主導を拡散側へ移した大規模化の系統では、雑音付き分類器が祖先抽出や暗黙的抽出を操り、尺度が適合と再現の交換を担った\autocite{btd030}。この分類器誘導は学習時の分類器追加と推論時の参照を伴い、分類器と大規模級上げの計算は重い\autocite{btd030}。分類器なし誘導はこの隘路に対し、単一網を条件付きと無条件で共同学習し、無条件化率一割から二割のもとで推論時に条件付き予測と無条件予測の外挿で誘導する転換を示した\autocite{btd032}。学習時の無条件化と推論時の外挿に寄せる点が、分類器誘導との違いである\autocite{btd030,btd032}。六十四画素条件で分布距離一・五五が弱い誘導で分類得点二六〇・二が強い誘導で、百二十八画素条件で二・四三と四二一が条件付きで報告された\autocite{btd032}。代償として二回の順伝播を要し、強い誘導では彩度が飽和し多様性が落ち、クラス条件を超える文条件への転移は確立していない\autocite{btd032}。この推論時外挿の考えはTransformer denoiserのクラスドロップアウトや潜在拡散の誘導尺度へ継承される\autocite{btd023,btd025,btd032}。

\subsection{連続時間への統一と潜在への移行}
離散時刻の拡散と得点整合が別定式化であり、新しい解法や厳密尤度や制御生成の拡張が阻まれていたことが隘路となった。連続時間の確率微分方程式 framing は両者を統一し、時間依存得点網と予測子修正子の解法、確率流常微分方程式による適応積分と厳密尤度、無条件得点からの条件生成を導入した\autocite{btd021}。アーキテクチャとしては過程別の得点網と修正子、目的関数としては連続雑音除去得点整合、推論時サンプリングとしては逆過程解法と適応積分が区別される\autocite{btd021}。解法の費用は依然として高く、千から二千回の評価を要し、過程の変種と修正子段階の調整が残ることも記録される。離散の雑音予測と連続の得点はこの framing のもとで対応し、解法の選択肢が拡散側にも得点側にも開かれた\autocite{btd021}。

設計空間整理の一件は、前処理と日程と解法の部品化で三十五回評価の抽出と分布距離の低下を示すが、ここは要旨のみの消費であり、前処理や日程や解法の内訳と backbone 対目的の分離や他 modality への主張は本文未検証として扱う限定が必要である\autocite{btd022}。要旨水準の数値を本文水準として語らないことが、本特集の約束である。画像域での焼き直しや音声映像への転移は、本文と modality 資料の裏付けを要する\autocite{btd022}。この限定を保ったまま、統一 framing の概念はFlow Matching（フローマッチング）の比較軸へ引き継がれる。

画素拡散の費用が次の隘路となり、知覚学習済み符号器で一度圧縮して潜在で拡散し一度だけ復号する転換が置かれた\autocite{btd023}。アーキテクチャは知覚と敵対項の符号器、時刻条件U網、文や配置の交叉注意であり、目的関数は潜在雑音整合である\autocite{btd023}。推論時サンプリングは十から二百段階の暗黙的軌道で行われ、畳み込み拡張で千画素級へ伸ばす。改善として二百五十六画素無条件条件でCelebA-HQ五・一一、FFHQ四・九八等の分布距離が報告された\autocite{btd023}。細部の第一段階律速と過圧縮の停滞、逐次費用も記録される。交叉注意の条件づけの役割や符号器対 prior の誤差帰属は焼き直し検証を要する。要旨限定の主張は後継の本文検証へ引き継がれる。本文検証の完了が後継課題として残る。この移行は表現欄の転換であり、骨格や目的の転換と取り違えない\autocite{btd023}。符号器は知覚損失と敵対項で学ばれ、正則化は標準正規への近似か復号側量子化のいずれかで、文や配置の条件は領域符号器と交叉注意で入る\autocite{btd023}。誘導尺度一・五から十で二百から二百五十段階を回す\autocite{btd023}。公開実装庫の版や重みの主張は本文未検証であり、本節の機構の重みは潜在拡散の報告に置く\autocite{btd023}。

\subsection{Transformer denoiserと流れ・整合性モデルへの分岐}
U網の帰納バイアスが必須という仮定と、目的関数の比較が backbone と絡んでいたことが隘路となった。Diffusion Transformer (DiT)は雑音除去器をTransformerへ置き換え、凍結変分符号潜在のパッチ化と適応正規化による零初期化で規模則を示した\autocite{btd025}。Scalable Interpolant Transformer (SiT)は同一 backbone で前向き結合を補間へ置き換え、速度目的と得点導出を整え、拡散と補間の差を分離し、確率流と逆過程の二解法を比べた\autocite{btd026}。アーキテクチャの交代と目的関数の交代はここで明確に区別される。同一 backbone での目的関数比較の方法、補間係数の特異点と正則化の対応が整理され、 backbone 対目的の分離概念は後継の評価設計へ引き継がれる\autocite{btd026}。装置依存の処理量数値や評価手順への鋭敏さ、画素空間での未検証は境界として残る\autocite{btd025}。Transformerへの置換えは骨格欄の転換であり、拡散の置換えではない\autocite{btd025}。凍結変分符号潜在のパッチ化と適応正規化による零初期化が骨格側、簡易雑音目的の継承とクラスドロップアウトによる誘導対応が目的・条件側である\autocite{btd025}。抽出計算を積んでも模型計算を補えないことが条件付きで示された\autocite{btd025}。Scalable Interpolant Transformer (SiT)は同一骨格で前向き結合を補間へ置き換え、速度目的と得点導出を整える\autocite{btd026}。四十万段十七・二対十九・五、七百万段八・三対九・六、誘導付き二百五十六画素二・〇六対二・二七が同一骨格・同一計算量のもとで条件付きで報告された\autocite{btd026}。

Flow Matching（フローマッチング）は常微分方程式模擬を要する最尤の費用に対し、条件付き場への回帰でsimulation-free学習（シミュレーション不要の学習）を可能とし、最適輸送版で軌道を直線化し、訓練費一定で解法費を下げた\autocite{btd027}。Rectified Flow（整流フロー）は線形補間の方向への最小二乗回帰で結合を直し、再流で輸送費と直線性を下げ、一段階蒸留へ接続した\autocite{btd028}。推論時サンプリングは適応解法やオイラー段階で実行され、直線軌道ほど少数段階で足りる\autocite{btd027,btd028}。整合性モデルは軌道上の任意点を起点へ写す自己整合性関数（consistency function）を置き、境界条件と跳躍変数化で既存拡散網を再利用し、蒸留と単独学習の二経路を示した\autocite{btd029}。改善として規模別の分布距離低下や評価数削減、一段階での距離低下が条件付きで報告された\autocite{btd025,btd027,btd029}。装置依存と端点特異性、再流誤差の蓄積も記録される\autocite{btd027,btd028}。条件別の規模と距離の対応、直線性と段階数の交換が整理され、映像編集側の比較へ引き継ぐ。少数段階の品質と後継蒸留との比較や誘導との相互作用は報告表の範囲を超えて確立していない\autocite{btd029}。端点特異性への正則化や離散化の発見的要素、再流回数と誤差蓄積の交換が整理され、 backbone 対目的の分離概念は後継の評価設計へ引き継がれる\autocite{btd026,btd027,btd028}。

\subsection{流れへの再定式化：simulation-free回帰（シミュレーション不要の回帰）と直線化}
連続時間の最尤が常微分方程式模擬の費用を要したことが隘路となった。Flow Matching（フローマッチング）はこの隘路に対し、周辺場の代わりに標本ごとの条件付き場へ回帰するsimulation-free学習（シミュレーション不要の学習）を置き、勾配の等しさを要点とする転換を示した\autocite{btd027}。ガウス条件付き経路と最適輸送版で軌道を直線化した\autocite{btd027}。分布距離六・三五対七・四八、評価数百四十二対二百七十四が条件付きで報告された\autocite{btd027}。代償として条件付きの最適輸送でも周辺場は最適輸送に至らず、適応解法の評価数は百超に留まる\autocite{btd027}。Rectified Flow（整流フロー）は線形補間の方向への最小二乗回帰で結合を直し、決定論的な結び直しで周辺を保ちつつ輸送費を下げ、再流でさらに直線化し、一段階蒸留は最後にのみ接続した\autocite{btd028}。改善として一段階で分布距離四・八五と再現率〇・五一が条件付きで報告され、一回整流は二段階以上で良く、二回整流は一段階で直線に近づく\autocite{btd028}。代償として厳密理論は厳密速度場と一意常微分方程式を要し、推定誤差は再流で蓄積するため一回の再流が推奨され、振る舞いの悪い場には平滑化を要する\autocite{btd028}。

\subsection{短縮と蒸留から実行費用への接続}
千段階級の抽出費用が隘路として残った。整合性モデルはこの隘路に対し、軌道上の任意点を起点へ写す自己整合性関数を置き、境界条件と跳躍変数化で既存拡散網を再利用し、教師解法からの蒸留と単独学習の二経路を示す転換を置いた\autocite{btd029}。小規模画像条件で一段階三・五五と二段階二・九三、六十四画素条件で一段階六・二〇と二段階四・七〇が条件付きで報告され、単独学習は教師なしで漸進蒸留に並ぶことが条件付きで示された\autocite{btd029}。代償として定理の誤差率は正則条件と高次解法を要し、連続時間は順方向微分を要し、貪欲探索は単峰仮定に立つ\autocite{btd029}。半減蒸留、consistency distillation（整合性蒸留）と単独学習、潜在での整合性蒸留、敵対的少数段階、携帯端末一段階などの方向が後継として整理され\autocite{btd078,btd079,btd080,btd081,btd082}、詳細な遅延や記憶量や量子化の数値は実行章に譲る。

\subsection{本節の読解上の境界}
本節の区別を保つ。U網対Transformerはアーキテクチャ、拡散対流れは目的関数、解法と蒸留は推論であり、条件を跨ぐ順位付けはしない。敵対的生成は明示密度を持たず、識別器の同期が崩れれば様式崩壊を招き、高次元での代理評価は分散が大きい\autocite{btd013}。現在の役割は vocoder や符号器や後学習の混合利用であり、単独の到達点ではない\autocite{btd013,btd014}。混合正則化は経路長をわずかに歪め、入力空間は絡んだままである\autocite{btd015}。分布距離と知覚経路の指標は質感偏りを伴い、非構造集合での相反と規模への留保が伴う\autocite{btd016}。敵対と拡散の比較は解像度と予算の規律を要する。拡散の尤度は他尤度模型に競争的でなく、情報の多くが知覚困難な細部に割かれ、千段階抽出は日単位の費用を要する\autocite{btd019}。段階費用の削減は暗黙的短縮や蒸留の後継を要する\autocite{btd020}。軌道短縮での距離低下や実装由来の分散の振る舞い、常微分方程式と確率流の開きは少数段階の境界である\autocite{btd020}。少数段階の到達点と後継蒸留との比較はここでは行わない。解法費用は依然として高く、過程の変種と修正子段階の調整を要し、アーキテクチャ利得と解法利得は絡む\autocite{btd021}。解法の最適性と後継比較、画像外への転移は未確立である。画素級の細部は第一段階に律速され、過圧縮は停滞し、逐次費用と誤用偏りの注意が伴う\autocite{btd023}。交叉注意の条件づけの役割と誤差帰属は焼き直し検証を要する。Transformer denoiserの処理量数値は装置に結びつき、評価手順に鋭敏であり、画素空間は未検証である\autocite{btd025}。規模則と後継や文条件転移の比較は後継資料を要する。最適輸送版でも周辺場は最適輸送に至らず、適応解法の評価数は百超に留まり、画素空間の比較を要する\autocite{btd027}。軌道直線性と再流の比較や高解像文条件転移は未確立である。整流の厳密理論は厳密速度場と一意常微分方程式を要し、推定誤差は再流で蓄積し、振る舞いの悪い場には平滑化を要する\autocite{btd028}。再流費用と後継比較、一次元外の結合の最適性は未確立である。整合性モデルの定理の誤差率は正則条件と高次解法を要し、連続時間は順方向微分を要し、貪欲探索は単峰仮定に立つ\autocite{btd029}。少数段階品質と後継や誘導との相互作用は報告表の範囲を超えて確立していない。大規模化の抽出は敵対側より遅く、分類器と大規模級上げの計算は重く、分布距離と分類得点と適合再現は不完全な代理である\autocite{btd020}。誘導機構と後継分類器なし誘導の比較、類別を超える文条件は未確立である。長期学習での濾波崩壊の可能性と対数尤度の不在、集約プールの安定速度の交換は記録される\autocite{btd014}。安定性解析は様式生成器時代の工学に継承される。逐次生成は高価であり、六十四画素単一尺度の大域構造は弱く、規模は図形処理装置の記憶と時間に制約される\autocite{btd017}。費用の主張と後継Transformerとの対比が残る。局所注意は全域文脈を切り詰め、最大生成は三十二画素に超解像を伴うのみである\autocite{btd018}。全域注意の大規模比較を要する。設計空間整理の一件は要旨のみの消費であり、前処理と日程と解法の内訳と backbone 対目的の分離と他 modality への主張は本文未確立である\autocite{btd022}。画像域の焼き直しと音声映像転移は本文と modality 資料の裏付けを要する。Scalable Interpolant Transformer (SiT)の端点特異性への正則化や離散化の発見的要素、拡散対流れの開きと再流比較は境界として残る\autocite{btd026}。画素再帰の費用主張と後継の対比、逐次生成の重さと並列学習の対比は整理の対象である\autocite{btd017,btd018}。転換ごとの読み分けを再確認する。雑音除去拡散は前向き固定と逆過程学習と簡易目的の機構であり、大規模化は骨格と誘導の寄与である\autocite{btd019,btd030}。暗黙的短縮は学習目的を変えない抽出側の転換である\autocite{btd020}。連続時間 framing は離散と得点の統一であり、設計空間整理は要旨限定である\autocite{btd021,btd022}。潜在移行は表現側であり、公開実装庫の主張に機構の重みを置かない\autocite{btd023}。Transformer置換えは骨格側、補間は目的・経路側であり、混同しない\autocite{btd025,btd026}。Flow Matching（フローマッチング）と整流はsimulation-free回帰（シミュレーション不要の回帰）と直線化の再定式化であり、少数段階の到達と後継比較は確立していない\autocite{btd027,btd028}。誘導は尺度と適合再現の交換を伴い、強い誘導では彩度飽和と多様性低下と二回の順伝播の費用が残り、クラス条件を超える文条件への転移は確立していない\autocite{btd030,btd032}。

\section{条件づけとアライメント}
\sectionkicker{CONDITIONING}
\label{sec:conditioning}

\begin{multicols}{2}
\raggedcolumns
文から絵や音を生むとき、条件づけはどこで行われるのか。学習の段階で条件器ごと育てるのか、推論の段階で誘導として外挿するのか、凍結した本体に適合器として足すのか、文脈参照として並べるのか。本節がたどるのは、凍結符号器の文条件づけから、分類器なし誘導と段階専門家集成、音響の言語整合、空間制御の適合までであり、学習時条件と推論時誘導と適合器と文脈参照の四区分で整理する\autocite{btd011,btd032}。通読の約束として、文整合と知覚品質は交換可能ではない。 prompt への従順さが上がったことと、絵や音としての良さが上がったことと、分布としての近さが増したことは別の物差しである\autocite{btd011}。符号器に関する知見は時代に結びつく。当時の符号器で得られた順序が、後の符号器でも保たれるとは限らない。

時代の流れをあらかじめ見渡す。起点は文から画像への生成が、条件器自体を jointly に育てずに深い言語理解を要した時代であり、対照言語画像事前学習が四億対の結合埋め込み空間を置き、雛形 prompt と集成でゼロショット分類が既存視覚 backbone に並ぶことが報告された\autocite{btd031}。この joint 空間を条件づけに用いる発想が、凍結大規模文符号器による拡散条件づけへ継承された\autocite{btd011}。忠実度と多様性の調整が雑音付きの別学習分類器を要した時代に、分類器なし誘導が条件づけ脱落で条件付きと無条件を jointly に学び、推論時のみ外挿する転換を示した\autocite{btd032}。文誘導拡散の別系統はこの誘導を雑音付き対照勾配の誘導と比較し、塗りつぶしと部分的編集へ拡張した\autocite{btd033}。共有雑音除去器が抽出段階を跨いで隘路となった時代に、専門家集成が共有 denoiser を事前学習したのち雑音水準で二分木分割して段階専門家へ微調整する転換を示した\autocite{btd034}。音響では対照言語音響事前学習が十二万超の対で千二十四次元 joint 空間を置き、余弦で雛形 prompt と照合する評価推論を示した\autocite{btd035}。凍結された大型拡散に空間制御を加える際の忘却と費用の隘路に対し、空間制御網が固定U網に学習複写の符号器と中間層を零初期化畳み込みで接続し、零から育てることで初期の破壊を避ける転換を示した\autocite{btd036}。各転換の改善は条件付きの報告であり、代償と境界を伴う。
\end{multicols}

\subsection{凍結符号器による文条件づけ：対照事前学習の継承}
文から画像への生成が、条件器自体を jointly に育てずに深い言語理解を要したことが隘路となった。対照言語画像事前学習は四億対の結合埋め込み空間を置き、雛形 prompt と集成でゼロショット分類が既存視覚 backbone に並ぶことが報告された\autocite{btd031}。アーキテクチャとしては画像と文の二符号器と学習温度が前者であり、目的関数は対称交差エントロピーである\autocite{btd031}。推論手順では雛形 prompt と埋め込み集成によるゼロショット合成が採られた\autocite{btd031}。この joint 空間を条件づけに用いる発想が、凍結大規模文符号器による拡散条件づけへ継承された\autocite{btd011}。

継承先のアーキテクチャとしては、凍結文符号器の系列と集約埋め込みを交叉注意で受ける多段拡散と文条件超解像が置かれた\autocite{btd011}。目的関数は文脱落つき拡散学習に区別され、一割の脱落で条件付きと無条件を jointly に学ぶ\autocite{btd011}。推論時サンプリングでは静的動的しきい値付き誘導と雑音添加 cascade が採られた\autocite{btd011}。改善としてはゼロショットでの分布距離の低下と、人手の整合が参照に迫り、描画課題での選好が条件付きで報告された\autocite{btd011,btd031}。代償として文整合は知覚品質や分布品質と別物であることが残る。人物写実の弱さと符号器知見の時代限定、内外資料の制約も記録される\autocite{btd011}。数え上げの失敗や細分類での開き、偏りと監視用途の懸念も本文記録の境界である\autocite{btd031}。条件別のゼロショット精度と分布距離の対応、符号器規模と整合の関係が整理され、後継の交叉注意へ接続する。学習時条件と推論時誘導の区別がここに定まる。条件別の符号器規模とゼロショット精度の対応、 cascade 段数としきい値の交換、資料構成と人物写実の限界が整理され、凍結符号器の利用概念は後継の多段条件づけへ引き継がれる\autocite{btd011}。知覚の歴史は本特集の範囲外であり、生成との関連でのみ読むという境界も保たれる\autocite{btd031}。対照事前学習の成果を生成の条件づけとして読むとき、ゼロショット分類の強さがそのまま文条件生成の強さにはならないことに注意する。分類の物差しと生成の物差しは別であり、 joint 空間の利用概念は後継の多段条件づけへ引き継がれるものの、そのままの転用は検証を要する\autocite{btd011,btd031}。

\subsection{推論時誘導への転換：分類器なし誘導と文誘導拡散}
忠実度と多様性の調整が、雑音付きの別学習分類器を要したことが隘路となった。分類器なし誘導は条件づけ脱落で条件付きと無条件を jointly に学び、推論時のみ外挿する転換を示した\autocite{btd032}。アーキテクチャではなく学習と推論の変更である点が重要であり、視覚と言語の大型拡散と級上げ器の構成は維持される\autocite{btd032}。目的関数は無条件化率一割から二割での jointly 学習であり、別系統では空文割合での微調整と塗りつぶし微調整が加わり、文と画像の結合を保つ\autocite{btd032,btd033}。推論時サンプリングは尺度掃引による外挿であり、二回の順伝播を要し、尺度が小さいほど分布類似、大きいほど識別得点に寄る交換が報告された\autocite{btd032}。

改善としては小解像での分布距離の低下や、人手の写実と文整合の優位、ゼロショットでの分布距離が条件付きで報告された\autocite{btd032,btd033}。代償として高い尺度での彩度飽和と二回計算、多様性低下が残る\autocite{btd032}。クラス条件のみの検証という限定も記録される。文誘導拡散の別系統はこの誘導を雑音付き対照勾配の誘導と比較し、塗りつぶしと部分的編集へ拡張した\autocite{btd033}。異常対象の失敗と選別資料の偏り保持も記録される\autocite{btd033}。条件別の尺度と距離の対応、誘導方式の比較が整理され、潜在拡散へ継承される。推論時誘導の考えがここに定まる。条件別の脱落率と尺度感度の対応、二回計算の費用と多様性の交換、分類器併用の要否の対比が整理され、 jointly 学習の概念は後継の誘導日程へ引き継がれる\autocite{btd032,btd033}。塗りつぶし拡張の設計も残る。文条件づけへの転移や脱落率の最適値を超える主張は、検証済みの値を離れて確立していないという留保が伴う。別系統の微調整で加わった空文割合や塗りつぶし微調整の寄与は、報告された構成のもとでのものであり、他の backbone や他の解像度への転用は検証を要する\autocite{btd033}。誘導尺度の掃引で得られた分布類似と識別得点の交換も、報告された小解像の条件に結びつく\autocite{btd032}。条件を離れた一般論にはしないことが、本節の読みの作法である。

\begin{center}
\small
\begin{tabularx}{\textwidth}{@{}p{0.24\textwidth}p{0.34\textwidth}X@{}}
\toprule
区分 & 変更の位置 & 例 \\
\midrule
学習時条件 & 学習で条件器ごと育てる & 凍結文符号器の交叉注意条件づけ \\
推論時誘導 & 推論で外挿して調整する & 分類器なし誘導の尺度掃引 \\
適合器と文脈参照 & 本体を凍結して足す & 空間制御網と段階専門家 \\
\bottomrule
\end{tabularx}
\end{center}

\subsection{段階専門家への分岐：共有 denoiser の分割}
共有雑音除去器が抽出段階を跨いで隘路となり、単一符号器の条件づけにも限りがあったことが次の課題となった。専門家集成は共有 denoiser を事前学習したのち、雑音水準で二分木分割して段階専門家へ微調整する転換を示した\autocite{btd034}。アーキテクチャとしては対数正規雑音の共有学習と三専門家構成、条件づけとしては大規模文符号器と対照文と対照画像の三符号器を独立脱落で束ねる点が中核である\autocite{btd034}。目的関数は共有事前学習ののち専門家微調整という二相であり、単一目的の微調整とは区別され、二値木の分割基準が記録された\autocite{btd034}。

推論時サンプリングは段階ごとに一専門家を用いるため一段あたり費用は変わらず、多段超解像と交叉注意偏向による学習なしの語描画が組み合わされた\autocite{btd034}。改善としてはゼロショットの分布距離で先行構成を下回り、六十万段階で八十万段階の共有 baseline を上回ることが条件付きで報告され、分布と整合の曲線全体での優位が伴った\autocite{btd034}。凍結符号器の系譜と並列に立つ位置づけであり、適合器条件づけと併存し後継の評価結合へ接続する。代償として専門家数が学習費を増やすことが残る\autocite{btd034}。中間区間追加の小ささと対照のみの構成性不足、前景の弱さも記録される。条件別の専門家数と距離の対応、符号器組合せの寄与が整理され、対照のみでは構成性に限りがあり、文符号器のみでは前景に弱さが残るという交換が示された\autocite{btd034}。条件別の分割位置と段階費用の対応、符号器組合せと構成性の交換、学習段階数と曲線全体の優位の関係が整理され、段階専門家の概念は後継の効率化設計へ引き継がれる。先行構成との対照はゼロショットの条件でのものであり、条件を離れた順位付けにはしない。

\subsection{音響の言語整合と空間制御の適合}
音響が柔軟な自然言語の教師を欠き、予測と条件づけの接続が乏しかったことが隘路となった。対照言語音響事前学習は十二万超の対で千二十四次元 joint 空間を置き、余弦で雛形 prompt と照合する評価推論を示した\autocite{btd035}。アーキテクチャとしては畳み込み音響符号器と文符号器に線形射影と学習温度を重ね、百二十八対の束で学習する構成であり、目的関数は対称対照損失である\autocite{btd035}。推論手順は雛形集合との類似度計算であり、拡散を伴わず、ゼロショットの分類として整理される\autocite{btd035}。改善としては環境音のゼロショットで人手水準に迫り、音楽対話話の完全分離や、監視設定で複数課題の高い結果が条件付きで報告された\autocite{btd035}。代償として、話や感情や鍵語では偶然水準に近いこと、文品質が隘路となること、特定増強が害となることが残る\autocite{btd035}。生成側の音響拡散での条件利用や人手整合の妥当性は未確立として扱う。条件別のゼロショット精度と文品質の対応、増強の功罪が整理され、潜在拡散の対照条件や音楽の指標検証へ接続する。 joint 空間の概念は生成側の条件利用の検証軸へ引き継がれる\autocite{btd035}。音響整合の限界も記録される。

凍結された大型拡散に空間制御を加える際、忘却と費用が隘路となった。空間制御網は固定U網に学習複写の符号器と中間層を零初期化畳み込みで接続し、零から育てることで初期の破壊を避ける転換を示した\autocite{btd036}。アーキテクチャとしては固定側と学習側の分離と微小条件符号器が中核であり、十二層と一中間層の複写規模が記録された\autocite{btd036}。目的関数は prompt 脱落つき条件学習であり、五割の脱落で無条件生成を保つ点が特徴である\autocite{btd036}。推論時サンプリングでは加算可能な複合制御と解像度重みが採られ、曖昧入力への prompt 補助が伴った\autocite{btd036}。改善としては素描の品質と忠実度、意味分割での重なりと分布距離が条件付きで報告され、万段階未満での急な収束が伴った\autocite{btd036}。代償として、曖昧入力は prompt を要すること、複合重みが手動であることが残る\autocite{btd036}。単一制御の framing と条件別資料規模の外側は未検証として記録される。制御信号の忠実度と標本品質は別物として読むべきことがここに定まる。条件別の制御強度と忠実度の対応、複写規模と収束の関係が整理され、訓練時条件づけと推論時合成の区別が保たれる。軽量適合器や参照適合器へ継承される。条件別の複写規模と収束速度の対応、脱落率と無条件保持の交換、解像度重みと複合制御の関係が整理され、零初期化の概念は後継の適合器設計へ引き継がれる\autocite{btd036}。

\subsection{四区分の使い分け：条件と誘導と適合と参照}
本節の四転換を四区分で振り返る。凍結文符号器の交叉注意条件づけは学習時条件の代表であり、条件器を育て直さず深い言語理解を借りる道を示した\autocite{btd011,btd031}。分類器なし誘導は推論時誘導の代表であり、学習時の脱落と推論時の外挿で忠実度と多様性のつまみを手に入れた\autocite{btd032,btd033}。段階専門家は学習時条件の発展形であり、共有学習ののち段階分割で denoiser を専門化させた\autocite{btd034}。音響の対照空間は評価推論としての joint 空間であり、生成条件としての利用は検証課題として残した\autocite{btd035}。空間制御網は適合器の代表であり、本体を凍結したまま零から育てることで忘却と費用を避けた\autocite{btd036}。

使い分けの指針を具体化する。文の理解を深めたいときは符号器の欄を見る。どの符号器を凍結し、系列と集約のどちらを交叉注意で受けるかが問われる\autocite{btd011}。生成の振る舞いを推論時に振りたいときは誘導の欄を見る。脱落率と尺度の掃引と二回計算の費用が問われる\autocite{btd032}。段階ごとに異なる仕事をさせたいときは専門家の欄を見る。分割位置と専門家数と学習費の交換が問われる\autocite{btd034}。新しい modality に言葉の教師を与えたいときは対照空間の欄を見る。対規模と文品質と増強の功罪が問われる\autocite{btd035}。凍結本体に空間や参照を足したいときは適合器の欄を見る。複写規模と零初期化と脱落率が問われる\autocite{btd036}。欄を取り違えれば、改善の帰属を誤る。符号器の利得を誘導の利得と混同せず、専門家の利得を解法の利得と混同しないことが、本節の契約の実践である。

次節への橋渡しとして、空間制御網の位置づけを確認しておく。空間制御網は条件づけの側から見れば適合器であり、次節の制御の側から見れば配置前史の継承先である\autocite{btd036}。軽量適合器や参照適合器へ継承されるという記述は、この両義性を指す。零初期化で固定側を壊さず育てる考えは、画像参照の分離注意や低階数因子の零初期化にも通じる。読者は本節の四区分を、次節の適合器の系譜を読む物差しとして携えてほしい。条件づけと制御は別の節に分かれるが、凍結本体に触れず足すという発想は共通である。

\subsection{符号器規模と誘導尺度：報告された交換の整理}
条件別の対応をさらに具体化する。凍結符号器の側では、符号器規模とゼロショット精度の対応、 cascade 段数としきい値の交換が整理され、対照事前学習の側では雛形 prompt と集成の寄与が整理された\autocite{btd011,btd031}。誘導の側では、脱落率と尺度感度の対応、二回計算の費用と多様性の交換が整理され、分類器併用の要否の対比が整理された\autocite{btd032,btd033}。専門家の側では、分割位置と段階費用の対応、符号器組合せと構成性の交換が整理された\autocite{btd034}。音響の側では、対規模とゼロショット精度の対応、束規模と温度の対応が整理された\autocite{btd035}。空間制御の側では、複写規模と収束速度の対応、脱落率と無条件保持の交換、解像度重みと複合制御の関係が整理された\autocite{btd036}。いずれも条件付きの対応であり、条件を離れた一般論にはしない。束ねる量や文脈長や装置の別をそろえたうえでなければ、数値の比較は意味を持たないという本特集全体の作法が、ここでも適用される。とくにゼロショットの条件での対照は、ゼロショットの条件でのものとして読む。学習段階数の差や符号器組合せの差がもたらす曲線全体の優位も、報告された評価手順のもとでのものである\autocite{btd034}。人手の写実や文整合の優位も、報告された課題と手順のもとでのものである\autocite{btd032,btd033}。素描の品質と忠実度や意味分割での重なりも、報告された制御と資料のもとでのものである\autocite{btd036}。

\subsection{本節の読解上の境界}
文への従順さは知覚や分布の品質と別物であり、符号器の知見は時代に結びつく\autocite{btd011}。人物写実の弱さと内外資料の制約、符号器知見の時代限定は本文記録の境界である\autocite{btd011}。超解像の文有用性の焼き直しは報告掃引の範囲を超えて確立していない。高い尺度での彩度飽和と二回計算と多様性低下、クラス条件のみの検証は誘導の境界である\autocite{btd032}。文条件づけへの転移や脱落率の最適値を超える主張は検証済みの値を離れて確立していない。異常対象と異常場面の失敗、選別済み公開資料の偏り保持は誘導併用の境界である\autocite{btd033}。曖昧入力の prompt 要請と複合重みの手動性、単一制御の framing は空間制御の境界である\autocite{btd036}。条件別資料規模の外側は報告曲線の範囲を超えて確立していない。数え上げの失敗や細分類での開き、偏りと監視用途の懸念は対照事前学習の境界である\autocite{btd031}。知覚の歴史は別巻の対象であり、生成条件づけの品質はここでは確立しない。専門家数の増加は学習費を増やし、対照のみでは構成性に限りがあり、文符号器のみでは前景に弱さが残る\autocite{btd034}。話や感情や鍵語の偶然水準への近さ、文品質の隘路、特定増強の害は音響対照の境界である\autocite{btd035}。生成側音響拡散での条件利用や人手整合の妥当性は未確立として扱う。

\section{制御と参照：空間・参照・主体性の保存}
\sectionkicker{CONTROL}
\label{sec:control}

\begin{multicols}{2}
\raggedcolumns
新しく生むことと、保ちながら変えることは別の技術問題である。本節がたどるのは、凍結模型に触れず空間と参照と主体を足す適合の系譜であり、配置前史から軽量適合、画像参照、少数主体、低階数、映像と音楽の時間信号まで接続する\autocite{btd037,btd039,btd042}。通読の約束として、制御信号の忠実度と標本品質は別物として読む。配置に従ったことと、絵や映像として良いことは別の物差しである。主体に関する主張は述べた手順に結びつく。枚数も prompt の形も評価手順も違えば、保存の度合いは比べられない。制御信号は生成の制御としてのみ読み、知覚の歴史としては語らない。

時代の流れをあらかじめ見渡す。起点は意味覆いから写実画像への制御が、正規化で配置を洗い流す時代であり、空間適応正規化が覆いから空間変化する変調母数を二層畳み込みで求め、配置を保ったまま活性を変える転換を示した\autocite{btd037}。この前史が凍結拡散への軽量適合器へ継承され、七十メガ級から数メガ級まで三段階の適合器が凍結符号器の四尺度へ特徴を注入し、配置保持を拡散期へ持ち込んだ\autocite{btd041}。文のみ条件では画像 prompt の能力がなく、文能力を壊さず参照を加えることが隘路となった時代に、画像 prompt 適合器が凍結拡散と凍結視覚符号器を保ち、大域埋め込みを四トークンへ射影して層別に加える転換を示した\autocite{btd038}。参照類似を超えて主体を一貫させ、かつ母数効率よく切替えることが隘路となった時代に、少数画像個人化が三から五枚に稀少識別子と類名詞の文を添え、全層微調整と自動生成類損失で主体を焼き付けつつ類を保つ転換を示し\autocite{btd039}、低階数適合が基底を凍結し問合せと値への低階数因子だけを学び、推論時に統合して遅延を増やさない転換を示した\autocite{btd040}。映像で撮影と対象の動きを分け、見た目なしに姿勢と軌跡だけで制御することが隘路となった時代に、運動制御が凍結映像 backbone に軽量撮影モジュールと対象モジュール二つを段階学習し、実景の姿勢注釈で学ぶ転換を示した\autocite{btd042}。音楽で旋律と強弱と律動を時間変化する信号として与えることが隘路となった時代に、音楽制御網がメル図の畳み込みU網を札条件で事前学習し、制御別多層知覚器と零畳み込みの適合器で微調整する転換を示した\autocite{btd043}。各転換の改善は条件付きの報告であり、代償と境界を伴う。
\end{multicols}

\subsection{配置の前史と軽量適合：洗い流さない工夫}
意味覆いから写実画像への制御が、正規化で配置を洗い流すことが隘路となった。空間適応正規化は覆いから空間変化する変調母数を二層畳み込みで求め、配置を保ったまま活性を変える転換を示した\autocite{btd037}。アーキテクチャとしては符号器なし残差生成器と多重解像識別器、変分様式符号器の組合せであり、敵対期の構成として整理される\autocite{btd037}。目的関数は対覆い画像学習であり、後の拡散制御とは学習様式が異なる点が重要である。生成手順では配置と様式の同時制御が行われ、人手選好の比較で評価された\autocite{btd037}。改善としては意味重なりと分布距離と人手選好が条件付きで報告され、従来の直接連結を上回ることが示された\autocite{btd037}。代償として、対資料を要すること、検索併用が分布距離では優れるが整合では劣ることが残る。拡散期への転移は継承としてのみ読み、前史の地位に留めるという境界が伴う。

この前史が凍結拡散への軽量適合器へ継承された。適合器は七十メガ級から数メガ級まで三段階があり、凍結符号器の四尺度へ特徴を注入し、配置保持を拡散期へ持ち込む\autocite{btd041}。立方時刻抽出で初期段階を厚く学ぶ目的関数が特徴である\autocite{btd041}。生成手順は重み付き和で合成でき、再学習を要せず、派生系へ汎化する\autocite{btd041}。改善として分布距離の低下が条件付きで報告された。条件別の適合規模と距離の対応、注入位置の寄与が整理される。多適合の手動合成と粗い素描の分散も記録され、参照適合へ接続する。条件別の適合規模と分布距離の対応、時刻抽出の偏りと初期段階の寄与、派生系への汎化範囲が整理され、軽量注入の概念は後継の参照適合へ引き継がれる\autocite{btd041}。配置保持の考えが残る。直接の容量比較は同時期引用の範囲を超えて確立していないという留保が伴う。

\begin{center}
\small
\begin{tabularx}{\textwidth}{@{}p{0.24\textwidth}p{0.34\textwidth}X@{}}
\toprule
適合の形 & 凍結するもの & 足すもの \\
\midrule
軽量配置適合 & 拡散本体と文符号器 & 四尺度への注入特徴 \\
画像参照適合 & 拡散と視覚符号器 & 層別四トークンの画像枝 \\
主体焼付けと低階数 & 基底行列と類知識 & 識別子文脈と低階数因子 \\
\bottomrule
\end{tabularx}
\end{center}

\subsection{画像参照の分離注意：文能力を壊さない足し方}
文のみ条件では画像 prompt の能力がなく、文能力を壊さず参照を加えることが隘路となった。画像 prompt 適合器は凍結拡散と凍結視覚符号器を保ち、大域埋め込みを四トークンへ射影して層別に加える転換を示した\autocite{btd038}。アーキテクチャとしては分離交叉注意の画像枝を層ごとに足し、問合せ共有で文枝と並列に扱い、学習母数を二千万級に抑えた点が中核である\autocite{btd038}。目的関数は画像脱落つき学習であり、文脱落の考えと対応し、推論時の重みで参照度を調整する\autocite{btd038}。

推論時サンプリングでは画像注意の重みで参照度を調整し、再利用と空間制御との併用が可能とされ、五十段階級の暗黙的軌道と誘導尺度が条件として記録された\autocite{btd038}。改善としては画像類似と文類似の両立が条件付きで報告され、千万級対での百万段階学習と開放語彙での汎化が添えられた\autocite{btd038}。代償として、内容と様式の類似に留まり、少数画像の主体一貫には届かないことが残る\autocite{btd038}。主体保存の指標比較の未確立と参照強度の交換も記録される。条件別の重みと類似度の対応、層別注入の寄与が整理され、文脈参照条件という区分がここに定まる。少数主体化と低階数適合へ接続し、映像の動作適合や音楽の制御適合と対をなす。条件別の参照重みと類似度の対応、層別注入と文枝並列の関係、対規模と汎化の交換が整理され、分離注意の概念は後継の主体保存へ引き継がれる\autocite{btd038}。類似限定の境界も記録される。条件別のトークン数と参照度の対応が整理され、後継の適合器設計へ引き継がれる。開放語彙での汎化は千万級対での百万段階学習という報告条件に結びつき、条件を離れた一般論にはしない\autocite{btd038}。空間制御との併用や再利用の可能性は、報告された五十段階級の暗黙的軌道と誘導尺度の条件のもとでのものであり、他の解法や他の尺度への転用は検証を要する\autocite{btd038}。

\subsection{主体の焼付けと低階数：少数画像と効率切替え}
参照類似を超えて主体を一貫させ、かつ母数効率よく切替えることが隘路となった。少数画像個人化は三から五枚に稀少識別子と類名詞の文を添え、全層微調整と自動生成類損失で主体を焼き付けつつ類を保つ転換を示した\autocite{btd039}。低階数適合は基底を凍結し、問合せと値への低階数因子だけを学び、推論時に統合して遅延を増やさない転換を示した\autocite{btd040}。アーキテクチャとしては、前者が凍結拡散の全面微調整と超解像の継承、後者が凍結行列と可換因子の分離と零初期化である\autocite{btd039,btd040}。目的関数は識別子微調整と因子学習に区別され、倍率と脱落で安定化され、事前保存項が前者に伴う\autocite{btd039,btd040}。

生成手順では識別子 prompt による再文脈化や視点合成と作風適用、因子差替えによる即時切替が行われる\autocite{btd039,btd040}。改善としては三十主体での主体類似と文類似の両立や人手の主体選択、巨大模型で万分の一の母数と速度向上が条件付きで報告された\autocite{btd039,btd040}。代償として稀少文脈での失敗や文脈外観の絡みと色ずれが残る\autocite{btd039}。訓練付近への過適合と複合バッチの困難、階数調整の限定も記録される\autocite{btd040}。条件別の枚数と類似度の対応、階数と精度の交換が整理され、拡散での低階数挙動は言語域の証拠に留まる\autocite{btd040}。映像動作適合への継承がここに定まる。条件別の画像枚数と主体類似度の対応、階数と切替え費用の交換、倍率と安定性の関係が整理され、識別子焼付けの概念は後継の参照制御へ引き継がれる\autocite{btd039}。過適合の範囲も記録される。条件別の脱落率と保存項の対応が整理され、後継の効率適合へ引き継がれる。主体の主張は述べた手順に結びつき、手順を離れた同一性一般には拡張しない。三十主体での主体類似と文類似の両立や人手の主体選択は、述べた枚数と prompt の形と評価手順のもとでの報告であり\autocite{btd039}、巨大模型での万分の一の母数と速度向上は言語域の証拠に留まる\autocite{btd040}。条件別の画像枚数と主体類似度の対応、階数と切替え費用の交換、倍率と安定性の関係が整理され、識別子焼付けの概念は後継の参照制御へ引き継がれる\autocite{btd039}。過適合の範囲も記録される。条件別の脱落率と保存項の対応が整理され、後継の効率適合へ引き継がれる。

\subsection{時間信号の分離制御：映像の動きと音楽の変化}
映像で撮影と対象の動きを分け、見た目なしに姿勢と軌跡だけで制御することが隘路となった。運動制御は凍結映像 backbone に軽量撮影モジュールと対象モジュール二つを段階学習し、実景の姿勢注釈で学ぶ転換を示した\autocite{btd042}。アーキテクチャとしては時間Transformerの第二自己注意への姿勢注入と、対象軌跡の疎な条件符号化が中核であり、外観層を凍結する点が特徴である\autocite{btd042}。目的関数は撮影ののち対象という二相であり、同時学習とは区別され、順序が撮影先行であることが記録された\autocite{btd042}。推論時サンプリングでは外観なしの姿勢軌跡条件を単独または jointly に与え、十六フレーム級の短尺で評価された\autocite{btd042}。改善としては基本と複雑の撮影誤差や対象誤差で先行法を下回り、文整合が維持されることが条件付きで報告され、粒子運動の評価と対照類似が添えられた\autocite{btd042}。代償として複雑撮影と複雑軌跡の同時制御の成功率が低いことが残る\autocite{btd042}。順序の固定と場面多様性の限定、参照数十件を要する適用も記録される。空間制御の考えの映像継承がここに定まる。音楽の時間変化制御と対をなし、編集保存の評価へ接続する。条件別の軌跡疎度と誤差の対応、二相学習の順序効果、参照件数と適用範囲の交換が整理され、分離制御の概念は後継の時間信号設計へ引き継がれる\autocite{btd042}。

音楽で旋律と強弱と律動を時間変化する信号として与えることが隘路となった。音楽制御網はメル図の畳み込みU網を札条件で事前学習し、制御別多層知覚器と零畳み込みの適合器で微調整する転換を示した\autocite{btd043}。アーキテクチャとしては事前学習網と制御網の分離、表現としては抽出された旋律と力学と律動の信号、時間幅六秒の窓が中核である\autocite{btd043}。目的関数は札事前学習ののち制御微調整であり、全脱落と区間覆いと任意部分集合の即興化が併用される\autocite{btd043}。推論時サンプリングは百段階級の暗黙的軌道と大域様式誘導で行われ、部分区間は即興で補い、拍弦予測の文駆動が伴った\autocite{btd043}。改善としては全制御での旋律精度や力学と拍の両立が条件付きで報告され、抽出制御での優位と人手の音楽性の向上が添えられた\autocite{btd043}。代償として学習窓が六秒であり外挿で雑音が増え分布距離が悪化することが残る\autocite{btd043}。創作制御の単調化と札のみ文で器楽のみの範囲、歌声対象外も記録される。制御信号の忠実度と標本品質の分離が音楽でも保たれる。適合器制御の多様式展開として位置づき、編集保存の評価へ接続する。条件別の信号組合せと精度の対応、窓内外の劣化率、即興補完の範囲が整理され、部分時間条件の概念は後継の音楽編集へ引き継がれる\autocite{btd043}。器楽限定の範囲も記録される。

\subsection{適合器の系譜：凍結本体に触れず足す}
本節の五転換を適合の視点で振り返る。共通するのは、本体を凍結し、足す部分だけを学ぶという発想である。軽量配置適合は拡散本体と文符号器を凍結し、四尺度への注入特徴を学んだ\autocite{btd041}。画像参照適合は拡散と視覚符号器を凍結し、層別四トークンの画像枝を学び、学習母数を二千万級に抑えた\autocite{btd038}。主体焼付けは拡散の全面微調整という例外を含むが、類知識の保存項で凍結時の知識を保つ工夫を伴った\autocite{btd039}。低階数適合は基底行列を凍結し、低階数因子だけを学び、推論時に統合して遅延を増やさない\autocite{btd040}。運動制御は映像 backbone の外観層を凍結し、撮影と対象の二モジュールを段階学習した\autocite{btd042}。音楽制御網は事前学習網を凍結し、制御別多層知覚器と零畳み込みの適合器で微調整した\autocite{btd043}。零から育てるという空間制御網の考え\autocite{btd036}は直接本節の範囲外だが、その継承先として軽量適合器と参照適合器が位置づく。

使い分けの指針を具体化する。配置を保ちたいときは空間適合の欄を見る。注入位置と適合規模と時刻抽出の偏りが問われる\autocite{btd037,btd041}。文能力を壊さず参照を加えたいときは分離注意の欄を見る。参照重みと層別注入と対規模が問われる\autocite{btd038}。少数の画像から主体を一貫させたいときは焼付けの欄を見る。画像枚数と識別子文脈と保存項が問われる\autocite{btd039}。切替えの効率を上げたいときは低階数の欄を見る。階数と倍率と脱落が問われる\autocite{btd040}。撮影と対象を分けたいときは分離制御の欄を見る。姿勢誤差と軌跡誤差と二相の順序が問われる\autocite{btd042}。旋律と強弱と律動を時間信号で与えたいときは音楽制御の欄を見る。信号組合せと窓内外の劣化が問われる\autocite{btd043}。欄を取り違えれば、類似の改善と一貫の改善を混同する。内容と様式の類似に留まる参照適合\autocite{btd038}と、主体を焼き付ける個人化\autocite{btd039}は別の欄である。

\subsection{保存と編集への橋渡し：制御忠実度と標本品質の分離}
本節の最後に、次章の編集への橋渡しとして分離の指針を固定する。制御信号の忠実度と標本品質は別物である。配置に従ったことと絵として良いことは別の物差しであり\autocite{btd041}、軌跡誤差の小ささと映像として自然なことは別の物差しであり\autocite{btd042}、旋律精度の高さと音楽として良いことは別の物差しである\autocite{btd043}。この分離を保つことが、編集の評価へ接続する条件となる。保存を伴う改変では、何を保つのかを制御信号の忠実度で測り、改変の良さを標本品質で測る。二つを一つの点数に丸めれば、保ったのか変えたのかが分からなくなる。

主体の側でも同様の分離が要る。主体類似と文類似の両立は二つの物差しの両立であり\autocite{btd039}、画像類似と文類似の両立も二つの物差しの両立である\autocite{btd038}。片方だけを見て適合器の良し悪しを論じないことが、本節の契約の実践である。参照強度のつまみは両者の交換を動かす\autocite{btd038}。識別子の焼付けは主体側へ寄せ、保存項は類側へ引き戻す\autocite{btd039}。低階数の階数は効率と精度の交換を動かす\autocite{btd040}。つまみの位置を記録せずに数値だけを比べれば、比較は意味を失う。条件と手順を添えて読むことが、適合器の系譜を読む作法である。運動制御の二相の順序も手順の一部であり、撮影先行という順序を変えれば結果は変わり得る\autocite{btd042}。音楽制御の信号組合せも手順の一部であり、全制御での報告と部分制御での報告は分けて読む\autocite{btd043}。抽出制御での優位と人手の音楽性の向上も、報告された窓と手順のもとでのものである\autocite{btd043}。

\subsection{本節の読解上の境界}
制御信号は生成の制御としてのみ読み、知覚の歴史としては語らない。主体の主張は述べた手順に結びつく。画像参照は内容と様式の類似に留まり、少数画像の主体一貫には届かない\autocite{btd038}。主体保存の指標比較は未確立であり、参照強度の交換が伴う。稀少文脈での失敗や文脈外観の絡みと色ずれ、訓練付近への過適合は少数画像個人化の境界である\autocite{btd039}。複雑撮影と複雑軌跡の同時制御の成功率の低さ、順序の撮影先行、場面多様性の限定は運動制御の境界である\autocite{btd042}。多適合の手動合成と粗い素描の分散は軽量適合の境界である\autocite{btd041}。直接の容量比較は同時期引用の範囲を超えて確立していない。複合バッチの困難と多層知覚器や正規化や偏向の除外、階数調整の要請は低階数適合の境界である\autocite{btd040}。拡散での低階数挙動は言語域の証拠に留まる。対資料の要請、検索併用の分布距離の優位と整合の劣位は前史の境界である\autocite{btd037}。拡散期への転移は継承としてのみ読む。六秒の学習窓と外挿での雑音増加と分布距離悪化、創作制御の単調化、札のみ文と器楽のみと歌声対象外は音楽制御の境界である\autocite{btd043}。

\section{生成から編集へ：保存を伴う改変の系譜}\sectionkicker{EDITING}\label{sec:editing}

\begin{multicols}{2}
\raggedcolumns
新規生成の品質を上げることと、既存画像の一部を残したまま改変することは、別の技術課題である。新規生成では画面全体を作り直せばよいが、編集では既知領域の保存と未知領域の作り替えを同時に満たさなければならない。本節は、案内画像の部分拡散とマスク合成に始まり、潜在空間での局所合成、生成過程の内部表現を操作する注意写像と最適化編集、日常語の指示で操作する指示追従学習、そして一段階の動画編集へ至る系譜をたどる\autocite{btd044,btd045,btd046}。各段階の共通の問いは、何を保存対象とし、どの範囲を局所的に変更し、どのような条件で失敗するかである。

保存の主張はつねにマスクと既知領域の定義に束縛される。新規生成の品質指標をそのまま編集に流用することはできない。たとえば知覚類似度は、正当で多様な補完に対して罰を与えうる\autocite{btd045}。また画像ごとの最適化を伴う手法では、実行時間と映像メモリ使用量が実用上の制約になる\autocite{btd046,btd048}。本節では各手法の報告値を条件付きの記述として扱い、条件をそろえない横断順位づけは行わない。報告された選好率や類似度の数値は、比較対象と課題設定とマスク種別に依存する\autocite{btd044,btd045}。

系譜の見取り図は次のとおりである。第一に、事前学習済み拡散模型を課題別再学習なしに使う部分拡散とマスク合成の手続きが確立した\autocite{btd044,btd045}。第二に、画素空間での合成の重さと継ぎ目を受けて、潜在拡散の内部で前景と背景を分けて扱う局所合成へ移行した\autocite{btd046}。第三に、語の差し替えが画像全体を作り替えてしまう問題に対し、交差注意写像の注入と対象埋め込みの最適化が内部表現を編集界面にした\autocite{btd047,btd048}。第四に、マスクや注意操作の手続きを露出させず日常語の指示だけで編集する指示追従学習が現れた\autocite{btd049}。第五に、画像拡散を一段だけ動画化する一段階動画編集が、単一動画からの時間的一貫した改変を示した\autocite{btd050}。本節で用いる三つの言葉を定めておく。保存対象とは、改変後も残すことが求められる領域または構造であり、マスク外既知領域、背景潜在、参照生成の注意構造、既存画像の保存領域、単一動画の外観と運動が段階ごとに相当する\autocite{btd044,btd045,btd046,btd047,btd048,btd049,btd050}。局所性とは、変更を限定する機構であり、時刻$t_0$、マスク合成、潜在混合、写像注入の打ち切り時刻と重み、二重ガイダンスの重み、疎参照範囲が相当する\autocite{btd044,btd045,btd046,btd047,btd049,btd050}。失敗条件とは、主張が成り立たない範囲であり、各手法の代償節に記録する\autocite{btd044,btd045,btd046,btd047,btd048,btd049,btd050}。新規生成との対照では、効率化の動機が異なる点に注意する。潜在化の高速化は、新規生成の効率化ではなく保存領域の再計算回避にある\autocite{btd046}。単一動画微調整は、生成器の記憶と運動事前分布の分離を前提とする\autocite{btd050}。以下ではこの順序で詳述する。
\end{multicols}

\subsection{部分拡散とマスク合成：再学習なしの保存手続き}
事前学習済みの拡散模型を任意の案内画像や任意マスクに対応させるには、課題別の再学習なしで既知領域を保存しつつ未知領域だけを作り替える手続きが必要だった\autocite{btd044,btd045}。変化は二つの方向で現れた。第一は、案内画像を時刻$t_0$まで雑音化してから学習済みスコア模型で逆拡散する部分拡散である。$t_0$を$0.3$から$0.6$の範囲で選ぶことで、案内への忠実さと実在感の交換を調整する\autocite{btd044}。$t_0$が小さいほど案内の構造が残り、大きいほど模型の事前分布に寄った実在感のある出力になる。この交換は利用者が明示的に選ぶ設計であり、唯一の正解は存在しない。

第二は、既知領域は入力を順方向に雑音化したものから取り出し、未知領域だけを模型予測から取り出してマスク合成し、往復の再拡散で境界をなじませる手順である\autocite{btd045}。ジャンプ長$10$と再標本$10$回の条件で、知覚類似度は$0.068$まで低下し、比較対象への選好は$56.25\%$と報告された\autocite{btd045}。報告された改善は条件付きである。人物描画入力での$L_2$距離は$32.55$まで低下し、実在感選好は$80$から$98\%$、満足度は$75$から$91\%$の範囲で示され、広域・狭域・細線・太線の四種マスクでは$95\%$の選好が報告された\autocite{btd045}。一方で部分拡散の側の条件も記録すべきである。妥当な案内と$t_0$探索が不可欠であり、白画素主体の案内では大きな$t_0$が必要になり忠実さが失われる\autocite{btd044}。

代償は両者に共通して明確である。反復再拡散は$250 \times 10$回相当の評価を要し低速である\autocite{btd045}。知覚類似度は正当な多様な補完に罰を与え、細線構造では破綻が残り、失敗事例は補遺に記録されている\autocite{btd045}。保存の主張はマスクと既知領域の定義に束縛され、新規生成の品質指標をそのまま流用できない点も編集固有の制約である\autocite{btd044,btd045}。ここで確立した部分拡散と保存領域合成の考え方は、潜在空間での局所編集へ継承される。すなわち、作り直す範囲と残す範囲を分離し、境界の整合を反復手続きで確保するという枠組みが、以後の編集手法の共通言語になった\autocite{btd044,btd045,btd046}。

\subsection{潜在局所合成：前景と背景の分離}
画素空間での合成が高品質でも、画像ごとの最適化が重く境界の継ぎ目や色ずれが残ることが次の隘路になった\autocite{btd046}。変化は潜在拡散の内部で前景と背景を分けて扱うことである。テキスト条件の前景潜在を各段階で生成しつつ、背景潜在は雑音化した入力から取り出して合成し、$8$分の$1$に圧縮された潜在解像度でダウンサンプリングされたマスクにより混合する\autocite{btd046}。必要に応じて画像ごとの復号器微調整で再構成を補い、複数の予測から言語画像類似度で順位付けして選択する。潜在化による高速化は新規生成の効率化とは動機が異なり、保存領域の再計算を避けて改変箇所だけを生成し直すことにある\autocite{btd046}。

報告された改善は速度と精度の両面である。画素空間の先行手法が約$25$分かかるのに対し桁違いに高速であり、比較対象の三手法に対する精度向上が報告されている\autocite{btd046}。背景保存と前景変更の分離により、生成し直しではなく保存を伴う改変という編集固有の課題設定が明確になった\autocite{btd046}。ただし代償は潜在表現の損失性に由来する。変分自己符号化器の再構成限界により継ぎ目やポアソン合成時の色ずれが生じ、潜在最適化は過度の平滑化を招く\autocite{btd046}。薄いマスクでは細部再現が依然として失敗する。保存の主張はマスク定義に束縛され、実行時間や映像メモリ使用量の定量化は消費本文では確立していない\autocite{btd046}。

自然さの軸での改善は、可知度や話者類似度とは区別して読む必要がある。ここでの自然さは画像編集の文脈での見た目の整合であり、音声系譜の評価軸とは無関係である\autocite{btd046}。この局所合成の枠組みは注意写像の操作へ継承される。潜在空間で保存と変更を分離する考え方は、以後、注意機構の内部表現を直接操作する手法の前提になった\autocite{btd046,btd047}。なお画像ごとの復号器微調整は任意の補正であり、つねに必要というわけではない。微調整なしでも合成は成立するが、再構成誤差が残る場合の選択肢として位置づけられる\autocite{btd046}。

\subsection{注意写像と最適化編集：内部表現を界面にする}
小さな指示語の差し替えが画像全体を作り替えてしまうことは、実用上の大きな隘路だった\autocite{btd047,btd048}。変化は生成過程の内部表現を編集界面にすることである。第一の方法は、同一乱数種で生成した参照側の交差注意写像を対象側生成の$64 \times 64$段階へ注入し、語入替では打ち切り時刻、語追加では重み付け、強弱調整では係数操作で局所変更を実現する\autocite{btd047}。反転が歪む場合には注意から導出したマスクで代替する。第二の方法は、対象埋め込みを雑音除去損失で最適化した後に生成器本体を超解像も含めて微調整し、補間係数$\eta$で忠実さと変更度の交換を制御し、八つの乱数種から選択する\autocite{btd048}。$\eta$は$0.6$から$0.8$の範囲で平均的な交換が報告されている\autocite{btd048}。

報告された改善は条件付きである。前者はケーキや自転車から自動車への変換、様式や仕様の変更、強弱スライダー操作を定性的に示した\autocite{btd047}。後者は$100$組の編集課題で$9000$超の二者択一回答により三種の先行手法に$70\%$超の選好を得た\autocite{btd048}。定量的選好と定性的作例は別の証拠であり、両者を対応づけて一般化することはできない。作例は可能性を示し、選好率は課題分布のもとでの傾向を示す\autocite{btd047,btd048}。

代償も共有される。反転は歪みを生みやすく、参照指示文が必要で、低解像注意が精度上限となり物体の空間移動はできない\autocite{btd047}。語彙外の対象や複雑な空間配置では弱さが残り、保存度と変更度の交換は係数と乱数種の選択に依存する\autocite{btd047,btd048}。顔や人物の再現は基盤模型の限界を継承し、編集層では解消されない\autocite{btd048}。画像ごとの最適化時間は実用上の制約であり、高速化は後継の指示学習に委ねられた\autocite{btd048}。参照指示文の必要性は編集の再現性を支える一方で、指示なし編集への一般化を残す\autocite{btd047}。注意由来マスクの代替は反転歪みの回避策として位置づけられる\autocite{btd047}。保存度と変更度の交換は、係数探索の手続きとともに記録すべきである\autocite{btd047,btd048}。

\begin{center}
\small
\begin{tabularx}{\textwidth}{@{}p{0.22\textwidth}p{0.24\textwidth}p{0.24\textwidth}X@{}}
\toprule
手法 & 保存対象 & 局所性の機構 & 失敗条件 \\
\midrule
部分拡散\autocite{btd044} & 案内画像の構造（$t_0$依存） & 時刻$t_0$までの雑音化と逆拡散 | $t_0 \in [0.3, 0.6]$ & 白画素主体の案内では忠実さが失われる \\
マスク合成\autocite{btd045} & マスク外の既知領域 | 順方向雑音化の既知領域と模型予測の未知領域の合成、再拡散 & 低速（$250 \times 10$回相当）、細線構造の破綻 \\
潜在局所合成\autocite{btd046} & 背景潜在（入力由来） & 潜在解像度での前景・背景混合、復号器微調整は任意 & 継ぎ目・色ずれ・平滑化、薄いマスクの細部失敗 \\
注意注入\autocite{btd047} & 参照生成の注意構造 | $64 \times 64$段階への写像注入、打ち切り時刻・重み・係数 & 反転歪み、空間移動は不可、低解像上限 \\
最適化編集\autocite{btd048} & 対象埋め込みの忠実さ（$\eta$依存） | 埋め込み最適化と生成器微調整、$\eta \in [0.6, 0.8]$、八種選択 & 低速、$\eta$と乱数種への依存、顔・人物は基盤継承 \\
\bottomrule
\end{tabularx}
\end{center}

\subsection{指示追従学習と一段階動画編集：手続きの隠蔽と時間への拡張}
マスクや注意操作の手続きを露出させず、日常語の指示だけで編集したいという要求が次の隘路になった\autocite{btd049}。変化は指示追従を学習課題に変えることである。$700$組の三つ組みから微調整した言語模型で$45$万超の指示文と説明文の組を生成し、安定拡散に画像条件チャネルを追加して指示組で学習し、方向性フィルタで品質を保つ\autocite{btd049}。推論時は画像とテキストの二重ガイダンスで編集強度を調整し、多段階の対話的連鎖も可能にする。報告された改善は交換曲線上の位置である。同一の方向類似度では画像類似度が高く、先行する案内画像手法を上回ることが示され、テキスト誘導$7.5$の条件での比較が報告され、全量と選別済みデータの除去検証では全量学習が他の選別条件を上回った\autocite{btd049}。生成し直しではなく既存画像の保存領域を残す点で、新規生成とは別の技術課題として位置づけられる\autocite{btd049}。

代償は新規生成の限界と保存課題の交差に現れる。視点変更や拡大、過度な変更、物体の分離や入替、個数や空間推論は失敗し、基盤拡散模型と言語模型の偏見を継承する\autocite{btd049}。微細すぎる変更や撮影角度のずれも弱点として記録されている。二重ガイダンスの重みは編集強度と保存度の交換を担い、多段階連鎖では誤差蓄積が懸念される\autocite{btd049}。視点や個数や空間推論の失敗は基盤模型の構成性限界の継承であり、編集層では解消されない。全量学習の優位は選別基準の再検討を示唆する。保存の主張は領域定義に束縛される\autocite{btd049}。この指示編集の系譜は動画への一段階適用へ接続する\autocite{btd049,btd050}。

一枚の動画だけから時間的一貫性を保ったまま対象動作や様式を変えることは、画像編集の延長では解けない隘路だった\autocite{btd050}。変化は画像拡散を一段だけ動画化することである。$3 \times 3$畳み込みを$1 \times 3 \times 3$へ膨張させ、最初と直前のフレームだけを参照する疎な時空間注意を挿入し、問い合わせ行列を中心に単一動画で約$500$段階微調整する\autocite{btd050}。推論時は無条件側を反転で求めてから編集指示文でサンプリングする。除去検証では疎注意なしで内容漂流、反転なしで動作停滞、微調整なしでちらつきが生じることが示され、各要素の役割が分離された\autocite{btd050}。

報告された改善は短区間の整合性である。$42$本の動画と$140$指示の条件で、枠間整合性は$92.40$となり比較対象の$90.64$と$88.89$を上回り、整合性選好は$87.86\%$と$62.14\%$、意味整合も同等以上だった\autocite{btd050}。単一動画微調整は、生成器の記憶と運動事前分布の分離を前提とする。疎注意の参照範囲は計算量と整合性の交換であり、全参照への拡張は検証されていない\autocite{btd050}。長時間地平の物語構造や同一性保持は確立しておらず、短区間整合の主張を超えて読んではならない。複数物体の重なりでは混合が生じ、対象画像模型の物体限界も継承する\autocite{btd050}。実行メモリや長時間整合の数値は報告されていない。編集保存の評価は枠間整合と意味整合の二軸で行い、単一指標への縮約は避けるべきである\autocite{btd050}。

\begin{center}
\small
\begin{tabularx}{\textwidth}{@{}p{0.22\textwidth}p{0.24\textwidth}p{0.24\textwidth}X@{}}
\toprule
手法 & 保存対象 | 交換の係数 & 局所性・時間の機構 & 失敗条件 \\
\midrule
指示追従学習\autocite{btd049} & 既存画像の保存領域 | 二重ガイダンスの重み | 画像条件チャネル追加の学習、対話的連鎖は誤差蓄積 & 視点・拡大・分離・入替・個数・空間推論の失敗、基盤偏見の継承 \\
一段階動画編集\autocite{btd050} & 単一動画の外観と運動 | 参照範囲（最初・直前のみ）と微調整約$500$段階 | 膨張畳み込みと疎時空間注意、反転による無条件側の取得 & 重なり混合、長時間同一性・物語構造は未確立、長区間への外挿は不可 \\
\bottomrule
\end{tabularx}
\end{center}

\paragraph{生成と編集の対照}
新規生成は画面全体の品質を問うが、編集は保存と変更の同時成立を問う\autocite{btd044,btd045,btd046}。時刻$t_0$、補間係数$\eta$、二重ガイダンスの重み、疎参照範囲はいずれも保存度と変更度または実在感の交換係数であり、唯一の正解ではなく利用者の選択として記録すべきである\autocite{btd044,btd048,btd049,btd050}。参照指示文の必要性は再現性を支える一方で指示なし編集への一般化を残し、八つの乱数種からの選択は交換の手続きの一部である\autocite{btd047,btd048}。多段階の対話的連鎖は利便性を上げるが誤差蓄積が懸念される\autocite{btd049}。

\paragraph{係数と手続きの記録作法}
再現に必要な手続き情報は、$t_0$範囲、ジャンプ長と再標本回数、潜在混合の解像度と復号器微調整の有無、注入段階と打ち切り時刻・重み・係数、$\eta$範囲と乱数種数、二重ガイダンス条件と指示組の規模、膨張形式と参照範囲と微調整段階数である\autocite{btd044,btd045,btd046,btd047,btd048,btd049,btd050}。選好率や類似度の数値は比較対象と課題設定とマスク種別に依存し、条件をそろえない横断順位づけは行わない\autocite{btd044,btd045,btd047,btd048,btd050}。実行時間・映像メモリ・量子化の多くは消費本文で報告されていない\autocite{btd044,btd045,btd046,btd047,btd048,btd049,btd050}。

\paragraph{短区間主張の範囲}
一段階動画編集の枠間整合$92.40$は、$42$本の動画と$140$指示という短区間条件の記述であり、長時間地平の物語構造や同一性保持の証拠にはならない\autocite{btd050}。疎注意の参照範囲を最初と直前のフレームに限る設計は、計算量と整合性の交換の選択であり、全参照への拡張は検証されていない\autocite{btd050}。編集保存の評価は枠間整合と意味整合の二軸で行い、単一指標への縮約は避けるべきである\autocite{btd050}。

\subsection*{読解上の境界}
保存の主張はマスクと領域定義に束縛され、画像ごとの最適化費用を添えて読む\autocite{btd044,btd045,btd046,btd047,btd048}。案内画像手法は妥当な案内と$t_0$探索を要し、白画素主体の案内では大きな$t_0$が必要になり忠実さが失われる\autocite{btd044}。マスク合成の反復再拡散は低速であり、知覚類似度は正当な多様補完に罰を与え、細線構造の破綻と補遺の失敗事例が残る\autocite{btd045}。潜在局所合成は損失性変分符号化に由来する継ぎ目と色ずれと平滑化を伴い、薄いマスクの細部は失敗し、実行時間と映像メモリ使用量の定量化は消費本文では確立していない\autocite{btd046}。注意注入は反転歪みと参照指示文の必要と低解像上限を伴い物体の空間移動はできない\autocite{btd047}。最適化編集は微細すぎる変更や撮影角度ずれに弱く$\eta$と乱数種に依存し、顔や人物は基盤限界を継承し画像ごとの最適化は低速である\autocite{btd048}。指示追従学習は視点・拡大・過度変更・分離・入替・個数・空間推論に失敗し基盤と言語模型の偏見を継承する\autocite{btd049}。一段階動画編集は重なり混合を伴い、疎注意・反転・微調整の各要素は除去検証で役割分離されたが、対象画像模型の物体限界を継承し、実行メモリや長時間整合の数値は報告されていない\autocite{btd050}。壁時計時間・映像メモリ・量子化・実時間係数の多くは消費本文で報告されていない\autocite{btd044,btd045,btd046,btd047,btd048,btd049,btd050}。新規生成の品質指標を編集の保存主張にそのまま適用してはならない\autocite{btd044,btd045,btd046}。

\section{音声・声の系譜：波形からEnd-to-End、ニューラルコーデック言語モデル、全二重へ}\sectionkicker{SPEECH}\label{sec:speech}

\begin{multicols}{2}
\raggedcolumns
音声合成の系譜は、画像生成の付随物ではない。波形の標本列を直接作る課題から出発し、系列変換とニューラルボコーダの分離、並列合成とHiFi-GANボコーダ、端末間統合、話者条件づけとニューラルコーデック言語モデル、流れ補完による非自己回帰化と実時間化、そして全二重対話へ至る独立の lineage である\autocite{btd051,btd052,btd053}。本節を通読する際の約束は、誤り率・可知度と自然さと話者類似度と韻律と遅延・実時間性を混同しないことである。各指標は別の軸であり、ある軸の改善を別軸の改善として読んではならない\autocite{btd051,btd057,btd058}。提供元作成の評価は隔離し、独立再現なしに順位づけを行わない\autocite{btd061,btd062}。

系譜の見取り図は次のとおりである。第一に、膨張因果畳み込みの自己回帰波形網と系列変換注意機構と復元網付き復号器が、素片結合なしの波形生成を切り開いた\autocite{btd051,btd052,btd053}。第二に、持続長予測による並列化とHiFi-GANボコーダの分離と条件付き変分自己符号化器による端末間統合が、遅さと注意破綻の隘路を解いた\autocite{btd054,btd055,btd056}。第三に、話者条件づけの端末間統合から残差符号の言語模型化への移行が、数秒の指示音声による未知話者のゼロショット話者複製を可能にした\autocite{btd057,btd058}。第四に、最適輸送条件付きFlow Matching（フローマッチング）による非自己回帰化と前後文脈の補充条件が、生成と編集の統一と実時間化を進めた\autocite{btd060,btd132,btd133}。第五に、多能高忠実模型と検証済み多言語翻訳系と実時間対話枠組みが、実環境の人間並み合成と対話遅延の両立を課題化した\autocite{btd061,btd062,btd134}。評価軸の定義を定めておく。誤り率は可知度の軸、平均意見値は自然さの軸、符号器類似度と類似度意見値は話者性の軸、実時間係数と段階数は遅延の軸である\autocite{btd051,btd054,btd057,btd058,btd132,btd133}。同一の数値が別軸を語ることはない。たとえば意見値$4.43$の報告は自然さの記述であり、可知度や話者性の証拠にはならない\autocite{btd056}。ボコーダの分離は忠実度と制御性の分離でもあり、韻律や話者性の制御はボコーダ層では扱わない\autocite{btd055}。音素前処理の残存や持続長誤差の伝播は、可知度の制約として読む\autocite{btd054,btd056}。模型の位置づけの軸で読むと、言語特徴から音響中間、離散符号、意味符号へと、何を確率変数とし何で条件づけるかが移っている\autocite{btd051,btd054,btd058,btd133}。この移行は第1節が扱う表現の選択の議論の音声側の現れであり、言葉による対応づけとして読む（第1節の本文には立ち入らない）。以下ではこの順序で詳述する。

証拠の深さに関する留保を先に記す。群化符号と言語模型改良の後継は結果表が消費本文で欠落しているため、本稿では群化と反復回避サンプリングという機構の変更のみを扱い、等価性や高速化の数値は確立済みとしない\autocite{btd059}。実時間対話枠組みは評価節が未消費のため、本稿では枠組みのみを扱い、品質順位や数値比較は行わない\autocite{btd134}。いずれも PARTIAL 消費であり、本文数値としての主張は行わない\autocite{btd059,btd134}。
\end{multicols}

\subsection{波形自己回帰と系列変換：素片結合なしの出発点}
波形を直接生成する段階の隘路は、一秒あたり一万六千点の標本列を素片結合なしに作ることだった\autocite{btd051}。変化は三段階で現れた。第一に、膨張因果畳み込みによる自己回帰波形網で、言語特徴や対数基本周波数の条件づけにより自然さを底上げした\autocite{btd051}。英語と中国語での平均意見値は$4.21$と$4.08$となり、従来の系列網や結合型を上回り、自然音声との隔たりを英語で$51\%$縮めたことが報告された\autocite{btd051}。ここでの平均意見値は自然さの軸であり、可知度や話者類似度とは別の評価である\autocite{btd051}。第二に、文字から音響への系列変換注意機構により、帯域メル目標への直接予測と後段処理網で可変長抑揚を扱えるようにした\autocite{btd052}。第三に、位置鋭敏注意と復元網付き復号器に改良型WaveNetボコーダを結合し、平均意見値$4.526$と自然音声$4.582$に迫る水準を示した\autocite{btd053}。

構成の細部は次のとおりである。波形網は対数圧縮二百五十六値のsoftmaxとゲート付き素子、残差と跳躍接続を備える\autocite{btd051}。系列変換側は注意の再帰復号器と帯域メル八十目標、後段処理網からなり、波形化は五十回反復の近似合成に頼る\autocite{btd052}。復元網付き復号は位置鋭敏注意と復号器内復元網に改良型WaveNetボコーダを結合し、三十層・三膨張循環・十要素混合・十六ビットのボコーダでメル（五十ミリ秒枠・十二・五ミリ秒ずれ）から波形化する\autocite{btd053}。

代償も共通する。受容野は$240$から$300$ミリ秒程度で長域抑揚を捉えきれず、逐次サンプリングは低速で、楽曲の長域整合は欠ける\autocite{btd051}。系列変換では反復や飛び越しが起き、近似波形合成由来の人工感が上限となる\autocite{btd052}。発音誤りや領域外固有名詞の弱さも残る\autocite{btd053}。サンプリングの逐次性は遅延の制約となり、並列化への動機となった\autocite{btd051,btd054}。可知度と自然さと話者性の三軸は、この段階ですでに分離して評価すべきである\autocite{btd051,btd052,btd053}。受容野の限界は韻律の模型的上限であり、文脈窓の拡張は後継の課題として残された。音楽への転用では長域整合の欠如が露呈した\autocite{btd051}。言語横断の隔たり縮小（英語五一％、中国語六九％）の報告は、四十四時間・百九話者の条件の枠内である\autocite{btd051}。これらは並列化とニューラルボコーダの系譜へ継承される\autocite{btd051,btd052,btd053,btd054}。

\subsection{並列化とニューラルボコーダと端末間統合：遅さと注意破綻への応答}
系列変換の逐次生成は遅く、注意の破綻が誤りなく読み上げる用途の隘路になった\autocite{btd054}。変化は持続長予測による並列化とHiFi-GANボコーダの分離、そして端末間統合である\autocite{btd054,btd055,btd056}。第一に、順伝播Transformerに持続長予測器と長さ調整器を備え、教師整列と系列蒸留で学習し、メル生成を$0.025$秒、全体を$0.180$秒まで短縮した\autocite{btd054}。難文$50$文での誤りは零となり、比較対象の$7$や$15$や$17$を下回った\autocite{btd054}。第二に、転置畳み込み生成器と多受容野融合、複数周期識別器からなるボコーダで、平均意見値$4.36$と自然音声$4.45$に迫り、画像処理装置で$3701$キロヘルツ、軽量版は中央処理装置でも実時間超えを達成した\autocite{btd055}。第三に、条件付き変分自己符号化器に量保存流れと単調整列探索を組み込み、HiFi-GANボコーダで直接波形化する端末間統合で、平均意見値$4.43$と真値$4.46$に並び、$67$倍速の実時間性を示した\autocite{btd056}。

ボコーダの分業は次のとおりである。ボコーダは音響から波形への変換層である\autocite{btd055}。多周期識別器の除去で意見値が$2.28$まで落ち、識別器の必需を示す\autocite{btd055}。端末間統合が統合したものは、事後分布と文章事前分布と流れと整列探索と持続長予測とボコーダによる復号であり、流れなしでは意見値$2.98$まで落ちる\autocite{btd056}。

代償として、流れ事前分布は必須で、単一話者では発音素変換の前処理が残り、持続長誤差は伝播する\autocite{btd056}。音素なしでは誤読が増え、性別や言語の不均衡も残る\autocite{btd056}。ボコーダの分離は忠実度と制御性の分離でもあり、韻律や話者性の制御はボコーダ層では扱わない\autocite{btd055}。速度と自然さと話者性の三軸は区別して読む必要がある\autocite{btd054,btd055,btd056}。実行資源や量子化の数値は報告されていない\autocite{btd054,btd055}。並列化の代償として教師整列への依存が生じ、持続長予測の誤差は音響全体に伝播する\autocite{btd054}。端末間統合は二段階の逐次学習を解消したが、流れ事前分布の必要性や音素前処理の残存が新たな制約となった\autocite{btd056}。未知話者へのゼロショット話者複製の評価は後継に委ねられ、単一話者の枠内の記述として読む\autocite{btd056}。速度と自然さと話者性の三軸は区別して読む必要がある\autocite{btd054,btd055,btd056}。

\begin{center}
\small
\begin{tabularx}{\textwidth}{@{}p{0.22\textwidth}p{0.26\textwidth}p{0.24\textwidth}X@{}}
\toprule
段階 & 機構 | 報告（条件付き） & 分離すべき評価軸 \\
\midrule
波形自己回帰\autocite{btd051} & 膨張因果畳み込み、言語特徴・対数基本周波数条件 | 意見値$4.21$・$4.08$、隔たり$51\%$縮小（英語） & 自然さ。可知度・話者類似度とは別 \\
系列変換\autocite{btd052}・復元網復号\autocite{btd053} & 文字→帯域メル注意、位置鋭敏注意と復元網＋WaveNetボコーダ | 意見値$4.526$対真値$4.582$ & 自然さ。反復・飛び越し・人工感は別途記録 \\
並列化\autocite{btd054} | 持続長予測器＋長さ調整器、教師整列＋系列蒸留 | メル$0.025$秒・全体$0.180$秒、難文誤り零 & 速度と可知度。持続長誤差の伝播に注意 \\
HiFi-GANボコーダ\autocite{btd055} | 転置畳み込み＋多受容野融合＋複数周期識別器 | 意見値$4.36$対真値$4.45$、$3701$キロヘルツ & 自然さと速度。韻律・話者制御はボコーダ外 \\
端末間統合\autocite{btd056} | 条件付き変分自己符号化器＋量保存流れ＋単調整列探索 | 意見値$4.43$対真値$4.46$、$67$倍速 | 自然さと速度。流れ事前分布は必須 \\
\bottomrule
\end{tabularx}
\end{center}

\subsection{何を何で表すか：言語・意味・音響・符号・波形の分担}
系譜を模型の位置づけの軸で読み直す。波形網は標本値の分布を言語特徴と対数基本周波数で条件づける\autocite{btd051}。系列変換は文字から帯域メルへの写像を注意で整列させ、復元網付き復号はメルからWaveNetボコーダへの二段構えとする\autocite{btd052,btd053}。並列Transformerは持続長で長さを決めてメルを一括生成し、ボコーダはメルから波形への変換に専念する\autocite{btd054,btd055}。端末間統合は文字の事前分布と線形スペクトルの事後分布を流れで結び、波形まで一貫して学習する\autocite{btd056}。ニューラルコーデック言語モデルは音響を離散符号に替え、第一量子化器を自己回帰、残りを非自己回帰で予測する\autocite{btd058}。流れ補完は対数メル上のベクトル場と持続長模型に分け、前後文脈の補充条件で生成と編集を統一する\autocite{btd060}。実時間二系統は意味符号化と言語模型と因果流れ、あるいはDiffusion Transformer (DiT)と文章精緻化層に分ける\autocite{btd132,btd133}。この移行は第1節が扱う表現の選択の議論の音声側の現れであり、言葉による対応づけとして読む（第1節の本文には立ち入らない）。

意味符号は内容の保持を担い、有限スカラ量子化の教師あり意味符号化器や第一層に意味を蒸留する分割残差符号化がこれに当たる\autocite{btd133,btd134}。音響符号は声質や韻律の保持を担い、残差符号化器の複数量子化器がこれに当たる\autocite{btd058}。神経符号は波形の離散近似そのものであり、言語模型が次符号予測の対象とする\autocite{btd058,btd061}。意味と音響の分離は内容忠実度と話者類似度の軸分離に対応する\autocite{btd057,btd058}。

\begin{center}
\small
\begin{tabularx}{\textwidth}{@{}p{0.20\textwidth}p{0.28\textwidth}p{0.26\textwidth}X@{}}
\toprule
段階 & 何を模型化するか & 条件・入力 & 境界・分離軸 \\
\midrule
波形直接\autocite{btd051} & 標本値の分布（対数圧縮二百五十六値） & 言語特徴・対数基本周波数 & 自然さ。受容野約$240$から$300$ミリ秒、逐次サンプリングは低速 \\
文字から音響\autocite{btd052,btd053} & 帯域メルの系列（注意による整列） & 文字、位置鋭敏注意と復元網 & 自然さ。反復・飛び越し・人工感は別途記録 \\
並列とボコーダ\autocite{btd054,btd055} & 持続長とメルの一括生成、メルから波形 & 教師整列と系列蒸留、転置畳み込みと複数周期識別 & 速度と可知度。韻律・話者制御はボコーダ外 \\
端末間\autocite{btd056,btd057} & 文字事前分布とスペクトル事後分布の結合 & 量保存流れと単調整列探索、話者同一性と言語埋め込み & 自然さと速度。流れ事前分布は必須 \\
ニューラルコーデック言語モデル\autocite{btd058} & 残差符号列（第一量子化器のみ自己回帰） & 音素と三秒音響指示 & 可知度と話者性を分離。三秒指示依存、非実時間 \\
流れと実時間\autocite{btd060,btd132,btd133} & 対数メル上のベクトル場と持続長、意味符号と因果流れ & 前後文脈の補充、区画因果と文章精緻化 & 遅延は構成依存。系列間の比較は未確立 \\
\bottomrule
\end{tabularx}
\end{center}

\subsection{ゼロショット話者複製とニューラルコーデック言語モデル：未知話者への拡張}
単一話者の枠を超えて未知話者の声を数秒の指示音声で複製することが次の隘路になった\autocite{btd057}。変化は話者条件づけの端末間統合からニューラルコーデック言語モデルへの移行である\autocite{btd057,btd058}。第一に、端末間統合の骨格に話者同一性条件や言語埋め込みを加え、一致損失を選択的に併用することで、未知話者での符号器類似度$0.864$、平均意見値$4.21$、類似度意見値$4.16$と真値並みを示し、声質変換でも同等水準を達成した\autocite{btd057}。第二に、残差符号化器の符号列を言語模型で扱い、第一量子化器を自己回帰、残りを非自己回帰で予測し、三秒の音響指示でゼロショット話者複製を実現した\autocite{btd058}。誤り率$5.9$と類似度$0.580$で従来の$7.7$と$0.337$を上回り、主観類似度も$4.38$と$3.45$を上回った\autocite{btd058}。ここでの誤り率は可知度の軸、類似度は話者性の軸、意見値は自然さの軸であり、混同してはならない\autocite{btd057,btd058}。

第三の後継は群化符号と言語模型の改良だが、結果表が消費本文で欠落しているため、本稿では群化と反復回避サンプリングという機構の変更のみを扱い、等価性や高速化の数値は確立済みとしない\autocite{btd059}。この PARTIAL 扱いは証拠の深さの境界であり、機構の継承関係のみを記録する\autocite{btd059}。群化の機構は群寸法一・二・四・八での群埋め込みと群予測、残り符号の非自己回帰化であり、サンプリングは反復率の閾値で核サンプリングから無作為抽出へ切り替える反復回避方式である\autocite{btd059}。群化後継の数値的主張は表の欠落により本稿では扱わず、機構の継承関係のみを記録する\autocite{btd059}。

代償として、サンプリングの不安定や無限反復、三秒指示への依存、雑音付き擬似音素標識の問題、非実時間性が残る\autocite{btd058}。誤用や詐称の危険も本文で指摘されている\autocite{btd058,btd059}。これらは流れ補完の系譜へ継承される\autocite{btd058,btd060}。可知度と話者性と自然さの三軸は、条件ごとに分離して読まなければならない\autocite{btd057,btd058}。指示音声の品質や雑音への依存は、実環境適用の制約として残る\autocite{btd057,btd058}。誤用防止の観点も、評価軸とは別に記録すべきである\autocite{btd058}。

\subsection{ニューラルコーデック言語モデルの生成方式とゼロショットの条件づけ}
生成方式を定式化しておく。復号のみ型で第一量子化器を自己回帰、第二から第八を非自己回帰で予測し、適応層正規化の段階づけと共有埋め込み・予測層を備える\autocite{btd058}。入力は音素と三秒の音響指示であり、サンプリング復号ののち符号化器復号器で波形化する\autocite{btd058}。ゼロショットは話者埋め込みの事前学習に代えて指示音声そのものを音響指示として与える方式であり、三秒指示への依存と雑音付き擬似音素標識の問題を伴う\autocite{btd058}。話者条件づけの端末間統合と比べ、話者情報を埋め込みではなく指示音響の符号列として受け取る点が違う\autocite{btd057,btd058}。

類似度と内容忠実度の交換関係として、一致損失は類似度と引き換えに自然さを損なう方向に働く\autocite{btd057}。強化学習の仕上げは報酬不正の指摘と添えて読む\autocite{btd061}。

\subsection{流れ補完と実時間化：非自己回帰と前後文脈の統一}
過去のみを見る自己回帰は遅く、標識付き整列を要することが実時間用途の隘路になった\autocite{btd060}。変化は流れ補完による非自己回帰と、実時間化の二系統である\autocite{btd060,btd132,btd133}。第一に、最適輸送条件付きFlow Matching（フローマッチング）で学習したTransformerベクトル場に持続長模型とボコーダによる復号を組み合わせ、前後文脈の補充条件で生成と編集を統一した\autocite{btd060}。ゼロショットでは誤り率$5.9\%$から$1.9\%$、類似度$0.580$から$0.681$へ進み、最大$20$倍の高速化が報告され、言語横断でも誤り率$10.9\%$から$5.2\%$へ進んだ\autocite{btd060}。第二に、Diffusion Transformer (DiT)と文章精緻化層からなる開放型実時間模型は、$32$段階で誤り率$2.42$、実時間係数$0.31$、$16$段階で$2.53$と$0.15$を示したが、非実時間の枠内である\autocite{btd132}。第三に、有限スカラ量子化の意味符号化器と言語模型と区画因果Flow Matching（フローマッチング）を組み合わせた実時間模型は、誤り率$2.47$と人間並みに迫り、実時間版でもほぼ無劣化と報告された\autocite{btd133}。

流れ補完の構成は八十次元対数メルと枠整列音素上のTransformerベクトル場と持続長模型からなり、常微分方程式解法の評価回数を十回未満に抑え、外部のHiFi-GANクラスのボコーダに渡す\autocite{btd060}。雑音除去の誤り$-8.8\%$・類似$+0.450$・意見$+0.80$の報告は、課題条件の枠内である\autocite{btd060}。一方は二十二層十六頭のDiffusion Transformer (DiT)と四層の文章精緻化層を組み合わせ、揺動サンプリングとオイラー解法で推論する\autocite{btd132}。他方は二十五ヘルツ意味符号化器と言語模型と区画因果Flow Matching（フローマッチング）から五十ヘルツメルを経てボコーダに渡す\autocite{btd133}。

ただし遅延は構成依存で、英語は資料不均衡で弱く、話者類似度の指標は話者検証依存である\autocite{btd132,btd133}。供給者系列の比較可能性は確立していない\autocite{btd132,btd133}。可知度と遅延と話者性の軸は区別して読む\autocite{btd060,btd132,btd133}。三十二段階と十六段階の報告（誤り率$2.42$・係数$0.31$と誤り率$2.53$・係数$0.15$）は、同一装置・短片条件の枠内の記述である\autocite{btd132}。開放型の実時間報告は構成依存であり、装置や段階数の拘束なしに比較してはならない\autocite{btd132}。英語の弱さや指標の検証依存は、供給者系列の比較可能性の限界を示す\autocite{btd132,btd133}。流れ補完の編集統一は、生成と改変の境界を接続する成果であり、可制御実時間化の系譜へ継承される\autocite{btd060}。実時間係数や段階数の報告は、同一構成内の記述に留める\autocite{btd132,btd133}。

\subsection{因果性と遅延と全二重：速さと同時性の分離}
非自己回帰化は過去のみを見る逐次性をやめ、前後文脈の補充条件で欠落部を埋める方式に替える\autocite{btd060}。他方は区画因果Flow Matching（フローマッチング）により部分入力の流れ化を可能にし、遅延模型を言語模型・Flow Matching（フローマッチング）・ボコーダの遅延の積み上げ$L=M\cdot(d_{lm}+d_{fm}+d_{voc})$として明示する\autocite{btd133}。遅延は構成依存であり、装置や段階数や区画長の拘束なしに比較してはならない\autocite{btd132,btd133}。速さの改善と同時性の実現を混同してはならない\autocite{btd062,btd134}。

対話の同時性は速さとは別の制約である。全二重は発話の交互ではなく重なりの扱いであり、割り込みと相づちと同時発話の検出と継続判断を要する\autocite{btd134}。枠組みは大規模言語骨格と実時間符号化器と階層深度Transformerからなり、二系列を音響遅延つきで並列に模型化し、時刻整列した文章接頭辞の内言で内容を支える\autocite{btd134}。実時間符号化器は二十四キロヘルツを十二・五ヘルツ・八十ミリ秒枠に落とし、第一層に意味を蒸留する分割残差符号化を用いる\autocite{btd134}。理論百六十ミリ秒・実践約二百ミリ秒の遅延は枠組み条件の記述であり、品質順位の証拠ではない\autocite{btd134}。翻訳系の流れ化は別系統の制約であり、単調注意の段階的方策と遅延・分散正則化と放出閾値からなり、平均遅延と長さ適応遅延と終了ずれで評価する\autocite{btd062}。

\subsection{Seed-TTS（多用途・高忠実音声生成）・翻訳・全二重対話：実環境と遅延の両立}
実環境での人間並み合成と対話遅延の両立が最後の隘路になった\autocite{btd061,btd062,btd134}。変化は三方向である。第一に、多能高忠実模型は音響符号を言語模型で扱い、指示微調整や教師あり微調整や強化学習で仕上げ、文脈内学習で誤り率$2.249$、類似度$0.762$と人間に迫り、微調整で$+0.37$、強化学習で$+0.14$の比較意見値を得て、難文誤り率を$7.585$から$6.423$へ下げた\autocite{btd061}。ただし難発音指示では弱く、遅く澄んだ発話への報酬不正が指摘され、評価は供給者作成で独立再現は確立していない\autocite{btd061}。第二に、検証済みの多言語翻訳系は非自己回帰の階層的単位復号と韻律単位拡張や表情声質変換を備え、翻訳課題の多目標で学習する\autocite{btd062}。表現力と実時間性の両立を翻訳課題の枠内で示したものであり、対話遅延の主張とは区別する\autocite{btd062}。

多能高忠実模型の管轄は符号化器から符号と言語模型とDiffusion Transformer (DiT)とボコーダまでであり、自己蒸留分解と音素持続長部なしの非自己回帰変種を伴う\autocite{btd061}。配備遅延〇・〇二八倍・〇・一三二倍の報告は、提供元（ByteDance）作成の条件付き主張として隔離する\autocite{btd061}。翻訳系の復号は部分語・文字・単位の階層的伸長からなり、部分入力流れ化のための長さ分離で三倍速を示す\autocite{btd062}。翻訳品質の報告（対英二六・六対二四・一、音声間二九・七対二五・八）と流れ化の報告（閾値〇・五で対英二〇・〇・遅延一・六八）は、Fleurs・二・三B規模の提供元（Meta）条件の枠内である\autocite{btd062}。短形式一般領域の研究段階で長形式は劣化し認定なしという境界も添えて読む\autocite{btd062}。

第三の実時間対話枠組みは、大規模言語骨格と実時間符号化器と階層深度Transformerからなり、理論$160$ミリ秒、実践約$200$ミリ秒の全二重遅延を掲げるが、評価節が未消費のため、本稿では枠組みのみを扱い、品質順位や数値比較は行わない\autocite{btd134}。実時間対話の枠組み的主張は、遅延の理論値と実践値の記述に留め、品質比較は評価節の消費を待つ\autocite{btd134}。対話の多巡回測定は区切評価の依存性を残す\autocite{btd134}。供給者評価は隔離し、可知度と自然さと話者性と遅延を区別する\autocite{btd061,btd062,btd134}。供給者作成の評価は隔離し、独立再現なしに順位づけを行わない\autocite{btd061}。多巡回測定の区切依存は、方法論の開放課題として残る\autocite{btd134}。これが現行の音声前線の到達点である\autocite{btd061,btd062,btd134}。

\begin{center}
\small
\begin{tabularx}{\textwidth}{@{}p{0.22\textwidth}p{0.26\textwidth}p{0.24\textwidth}X@{}}
\toprule
段階 & 機構 | 報告（条件付き） | 境界 \\
\midrule
話者条件端末間\autocite{btd057} & 話者同一性条件・言語埋め込み、一致損失の選択併用 | 類似度$0.864$、意見値$4.21$、類似意見$4.16$ | 可知度・自然さ・話者性を分離 \\
ニューラルコーデック言語モデルゼロショット\autocite{btd058} | 残差符号の言語模型化、第一量子化器のみ自己回帰 | 誤り$5.9$対$7.7$、類似$0.580$対$0.337$ | 三秒指示依存、非実時間、誤用危険 \\
群化後継\autocite{btd059} | 群化符号＋反復回避サンプリング（機構のみ） | 数値は未確立（表欠落、PARTIAL） | 機構の継承のみ記録 \\
流れ補完\autocite{btd060} | 最適輸送条件付きFlow Matching（フローマッチング）＋補充条件 | 誤り$5.9\%\to1.9\%$、最大$20$倍速 | 生成と編集の統一、軸分離は継続 \\
実時間二系統\autocite{btd132,btd133} | Diffusion Transformer (DiT)＋精緻化層 / 有限スカラ量子化＋区画因果流れ | 誤り$2.42$・係数$0.31$ / 誤り$2.47$ | 遅延は構成依存、系列比較は未確立 \\
多能・翻訳・対話\autocite{btd061,btd062,btd134} | 指示・教師あり・強化学習 / 階層単位復号 / 全二重枠組み | 誤り$2.249$・類似$0.762$ / 翻訳枠内 / 遅延$160$・約$200$ミリ秒 & 供給者評価は隔離、対話枠組みは枠組みのみ \\
\bottomrule
\end{tabularx}
\end{center}

\paragraph{供給者評価の隔離作法}
多能高忠実模型の評価は供給者作成で独立再現は確立しておらず、供給者作成の数値は測定条件付きの主張として隔離する\autocite{btd061}。微調整による$+0.37$や強化学習による$+0.14$の比較意見値は、その作成条件の枠内で読む\autocite{btd061}。難発音指示の弱さと遅く澄んだ発話への報酬不正の指摘も添えて読む\autocite{btd061}。提供元（ByteDance・Meta）の別を明示して読む\autocite{btd061,btd062}。翻訳系の主張は翻訳課題の枠内の記述であり、対話遅延の主張とは区別する\autocite{btd062}。独立再現なしに順位づけを行わない\autocite{btd061,btd062}。

\paragraph{PARTIAL項目の扱い}
結果表の欠落により群化と反復回避サンプリングという機構の変更のみを扱い、等価性や高速化の数値は確立済みとしない\autocite{btd059}。評価節未消費のため、大規模言語骨格と実時間符号化器と階層深度Transformerからなる枠組みと遅延の理論値・実践値のみを扱い、品質順位や数値比較は行わない\autocite{btd134}。いずれも証拠の深さの境界であり、消費の拡張を待つ\autocite{btd059,btd134}。対話の多巡回測定の区切依存は、方法論の開放課題として残る\autocite{btd134}。

\paragraph{持続長と整列の制約}
並列化は教師整列と系列蒸留に依存し、持続長予測の誤差は音響全体に伝播する\autocite{btd054}。端末間統合は単調整列探索を内蔵するが、流れ事前分布は必須であり、単一話者では発音素変換の前処理が残る\autocite{btd056}。音素なしでは誤読が増え、持続長誤差と発音素誤りの区別は可知度の診断に必要である\autocite{btd054,btd056}。流れ補完の補充条件は標識付き整列の要求を緩和するが、遅延は構成依存であり、装置や段階数の拘束なしに比較してはならない\autocite{btd060,btd132}。性別や言語の不均衡、英語の資料不均衡、話者類似度指標の話者検証依存は、条件ごとに分離して記録する\autocite{btd056,btd132,btd133}。符号の役割分離（意味・音響・神経符号）は内容忠実度と話者類似度の軸分離に対応する\autocite{btd057,btd058,btd133}。因果性の制約（区画因果・放出閾値・音響遅延）と同時性の制約（割り込み・重なり・終了ずれ）は、速度の軸とは別に記録する\autocite{btd062,btd133,btd134}。

\subsection*{読解上の境界}
誤り率・可知度、自然さ、話者類似度、韻律、遅延・実時間性は区別し、提供元評価は隔離する\autocite{btd051,btd057,btd058,btd061,btd062}。受容野は約$240$から$300$ミリ秒で長域抑揚を捉えきれず、逐次サンプリングは低速で、楽曲の長域整合は欠ける\autocite{btd051}。系列変換の反復・飛び越しと近似合成の人工感、領域外固有名詞の弱さが残る\autocite{btd052,btd053}。並列化は教師整列と系列蒸留に依存し持続長誤差が伝播し、流れ事前分布は必須で単一話者の発音素前処理が残る\autocite{btd054,btd056}。ゼロショット話者複製は三秒指示に依存しサンプリングの不安定と無限反復と雑音付き擬似標識と非実時間性を伴い、誤用・詐称の危険が本文で指摘されている\autocite{btd058}。詳細表が消費本文で欠落しており機構のみを扱い、等価性・高速化の数値は確立済みとしない（PARTIAL）\autocite{btd059}。流れ補完の遅延は構成依存で英語は資料不均衡で弱く、話者類似度指標は話者検証依存で供給者系列の比較可能性は確立していない\autocite{btd132,btd133}。Seed-TTS（多用途・高忠実音声生成）は難発音指示に弱く報酬不正が指摘され評価は供給者作成で独立再現は確立していない\autocite{btd061}。枠組みのみを扱い品質順位・数値比較は行わない（PARTIAL）。対話の多巡回測定は区切評価に依存する\autocite{btd134}。実行資源・量子化・並列推論遅延・多話者ゼロショット類似度の細部は報告されていない\autocite{btd054,btd055,btd056}。

\section{音楽・一般音響の系譜：コーデック言語モデルと潜在拡散}\sectionkicker{MUSIC}\label{sec:music}

\begin{multicols}{2}
\raggedcolumns
音楽・一般音響の生成は、音声合成の延長でも画像生成の付録でもない。分単位の楽曲の長域整合、文章と音楽の整合、旋律や拍や和音の可制御性、歌詞の扱い、継続と補完という独立の課題設定を持つ\autocite{btd063,btd065}。本節は、階層的自己回帰から意味と音響の段階化、一段階可制御符号音楽、波形拡散と潜在拡散、時間条件と楽理制御、大規模人間選好検証へ至る系譜を、音響忠実度と文章音楽整合と長域形式と歌詞整合と継続補完の分離によりたどる\autocite{btd063,btd065,btd066}。各段階の報告値は同一条件内の記述であり、持続時間や手順の互換性のない横断順位づけは行わない\autocite{btd066,btd067,btd068}。

本節の評価に関する約束を先に記す。距離や乖離や整合の自動指標は代理水準に留め、閉鎖系の順位づけは人間選好検証を経なければならない\autocite{btd136}。距離指標と対照整合指標には既知の限界があり、条件の異なる数値を直接並べて優劣を論じることはしない\autocite{btd063,btd065,btd066,btd067,btd068,btd136}。具体的には、音楽 caps 十秒評価、音響 caps 検証条件、開放距離条件など、各報告はその条件の枠内に留まる\autocite{btd065,btd067,btd068}。相関の exact 値は図表束縛で本文数値としては確立しない\autocite{btd136}。

系譜の見取り図は次のとおりである。第一に、三種符号化器の階層的自己回帰と、凍結音響符号・自己教師意味符号・結合音楽文章符号の三段階変換が、分単位楽曲の生音響生成を切り開いた\autocite{btd063,btd065}。第二に、遅延配置の単一Transformer復号器による一段階自己回帰符号音楽が、多段階層の低速と弱い可制御性に応答した\autocite{btd066}。第三に、波形拡散と潜在拡散の二系統が、局所忠実度の上限と文章条件の一般音響化に応答した\autocite{btd064,btd067}。第四に、変分潜在と楽理交叉注意による開放重みでの時間条件と楽理属性の精密制御が現れた\autocite{btd068,btd069}。第五に、大規模人間選好検証が自動指標と人間の判断の対応を測定し、本系譜の評価境界を定めた\autocite{btd136}。条件枠の読み方を定めておく。音楽 caps 十秒評価、音響 caps 検証条件、声素材のボコーダによる波形生成条件、$10$秒断片条件、十秒器楽断片条件は、いずれもその枠内の記述であり、枠外への一般化は行わない\autocite{btd065,btd066,btd067,btd064,btd069,btd136}。持続時間の異なる報告の横断比較は行わない。$15$秒刻みの継続、$30$秒学習の延長、約$24$秒の上位文脈は、それぞれ別の条件の記述である\autocite{btd065,btd063}。開放重みと閉鎖系の比較は、資料条件が異なるため条件をそろえずに行わない\autocite{btd068}。距離指標は音響分布の隔たりの代理であり、対照整合指標は文章と音楽の対応の代理である。いずれも長域形式や歌詞整合や動機可制御性を測らない\autocite{btd063,btd065,btd136}。以下ではこの順序で詳述する。

本節では二つの分離を読み筋とする。局所忠実度と長域構造は別の軸であり\autocite{btd063,btd064}、文章整合と音楽的一貫性も別の軸である\autocite{btd065,btd067}。延長と維持、生成と継続と編集と制御も分け、評価は音質と整合と構造と選好に分割する\autocite{btd065,btd067,btd068,btd069,btd136}。
\end{multicols}

\subsection{階層的自己回帰と意味・音響段階化：分単位楽曲への第一歩}
分単位の楽曲を生音響のまま生成する隘路は、歌唱と伴奏の長域整合を保つことだった\autocite{btd063}。変化は階層的自己回帰から意味と音響の段階化である\autocite{btd063,btd065}。第一に、三種の符号化器で$8$や$32$や$128$の間引き率の符号列を作り、$50$億級の疎Transformerで上位事前分布を、二種の上位変換で下位を補い、窓化や呼び水の祖先サンプリングで延長した\autocite{btd063}。帯域収束は下位で$23.0$デシベルに達したが、整合は約$24$秒の上位文脈までで、その後は窓延長で漂流し、繰返し副歌や旋律記憶はなく、新規様式への一般化は弱い\autocite{btd063}。第二に、凍結音響符号と自己教師意味符号と結合音楽文章符号の三段階変換で、文章条件を音響条件に置換して推論し、$15$秒刻みの物語式継続や旋律条件拡張を可能にした\autocite{btd065}。音楽 caps 十秒評価で距離$4.0$と$9.6$や$13.4$、乖離$1.01$、整合$0.51$、人間勝利$312$と$158$や$97$、記憶再現は$0.2\%$未満と報告された\autocite{btd065}。

階層的自己回帰の因果を掘り下げる。上位は長い文脈を、下位は波形の質感を受け持つ\autocite{btd063}。$8192$符号の約$24$秒では整合が保たれるが、窓をずらすと漂流する\autocite{btd063}。符号帳比較の下位$-23.0$と$-15.9$と$-40.5$デシベルは局所再構成であり、副歌の証拠ではない\autocite{btd063}。

意味と音響の段階化の因果を掘り下げる。音響$50$ヘルツ残差$12$帳と意味$25$ヘルツと結合$128$次元$10$秒窓が分担する\autocite{btd065}。学習では意味を音響側で条件づけ、推論では文章側に置換する\autocite{btd065}。$4.3$億級三Transformerは温度$1.0$と$0.95$と$0.4$で分担し、意味除去で乖離$1.01$から$1.05$へ、整合$0.51$から$0.49$へ変化した\autocite{btd065}。$30$秒延長で節等の模型はなく失敗を継承する\autocite{btd065}。

代償として、結合符号の否定や時間順序の失敗を継承し、$30$秒学習の自己回帰延長で節や副歌の模型はなく、模型公開もない\autocite{btd065}。音響忠実度と文章音楽整合と長域形式は区別して読む\autocite{btd063,btd065}。窓延長による漂流は長域構造の未確立を示し、節や副歌の記憶は後継の課題として残る\autocite{btd063}。結合符号の失敗継承は条件符号化の限界であり、音楽理論属性の直接制御は一段階可制御系に委ねられた\autocite{btd065,btd066}。記憶再現の低さは学習資料の扱いの報告であり、品質の保証ではない\autocite{btd065}。これらは一段階可制御の系譜へ継承される\autocite{btd063,btd065,btd066}。

局所忠実度と長域構造は別の軸である。帯域収束の改善は再構成誤差を縮めるが、文脈切断の漂流を解消しない\autocite{btd063}。$15$秒刻みの継続は前断片を引き継ぐ推論手順であり、楽式記憶の証拠ではない\autocite{btd065}。

\subsection{一段階可制御符号音楽：多段階層の低速への応答}
多段階層の低速と弱い可制御性が次の隘路になった\autocite{btd066}。変化は一段階の自己回帰符号音楽である\autocite{btd066}。$32$キロヘルツ単一信号を$640$間隔で$50$ヘルツの四重残差符号に落とし、遅延配置の単一Transformer復号器で文章交叉注意と色度接頭辞を備え、誘導脱落$0.2$で学習し、上位サンプリングと誘導尺度$3.0$で生成し、立体化微調整では符号帳を倍化する\autocite{btd066}。音楽 caps 条件での報告は、距離$3.1$と$7.5$や$14.8$、乖離$1.28$、整合$0.31$、全体と関係性の人間評価は$78.43$と$81.11$から$84.81$へ、旋律色度類似は$0.66$と文章のみの$0.10$とされた\autocite{btd066}。これらは同一条件内の記述であり、他条件との順位づけには用いない\autocite{btd066}。

遅延配置と旋律条件を掘り下げる。各帳を一段ずつずらして単一復号器で予測し、多段階を一段階にまとめる\autocite{btd066}。類似$0.66$と$0.10$の差は同一条件内の旋律引き継ぎである\autocite{btd066}。色度は客観指標と交換し、平坦化は$6000$と$1500$段階の交換がある\autocite{btd066}。

代償として、微細制御は分類器なし誘導のみに依存し、平坦化が他の条件を上回るが$6000$と$1500$段階の交換があり、色度条件は客観指標を劣化させる\autocite{btd066}。旋律の精密制御には音響増強の研究が必要とされる\autocite{btd066}。音響忠実度と整合と可制御性は区別される\autocite{btd066}。潜在拡散の対抗系譜と並んで、閉鎖系の能力比較の文脈となる\autocite{btd066,btd067}。微細制御の誘導依存は可制御性の上限であり、楽理属性の直接操作は後継の順次交叉注意に委ねられた\autocite{btd066,btd069}。色度条件と客観指標の交換は、制御と品質の分離を示す\autocite{btd066}。立体化の符号帳倍化は、段階数を保ったままの拡張であり、遅延の増加を伴わない点が特徴である\autocite{btd066}。音響忠実度と整合と可制御性は区別される\autocite{btd066}。

文章整合と音楽的一貫性は別の軸である。整合$0.31$は対応の代理、類似$0.66$は引き継ぎの代理である\autocite{btd066}。直接操作は順次交叉注意に委ねられる\autocite{btd066,btd069}。

\subsection{波形拡散と潜在拡散：局所忠実度と一般音響化の二系統}
局所忠実度の上限と文章条件の一般音響化が次の隘路になった\autocite{btd064,btd067}。変化は波形拡散と潜在拡散の二系統である\autocite{btd064,btd067}。第一に、双方向膨張畳み込みの雑音予測網で非加重変分目的を最適化し、声素材$22.05$キロヘルツのボコーダによる波形生成で平均意見値$4.44$と波形網の$4.43$や真値の$4.52$に並び、軽量高速版でも$4.37$を保ち、浮動$32$で$2.1$倍速の実時間性を示した\autocite{btd064}。ただし局所忠実度の波形拡散に留まり、長域形式や歌詞整合や動機可制御性は確立していない\autocite{btd064}。第二に、メル変分符号化で$4$分の$1$に圧縮し、対照言語音響整合で整列した潜在拡散を音響埋め込みで学習し文章埋め込みでサンプリングし、音響のみ混合で増強した\autocite{btd067}。音響 caps 検証条件で距離$23.31$と$47.68$、指標$8.13$と$4.01$、乖離$1.59$と$2.52$、距離$1.96$と$7.75$、全体$65.91$と$45.0$と報告された\autocite{btd067}。

波形拡散の因果を掘り下げる。$30$層核$3$の膨張畳み込みで雑音を予測し、$20$から$200$段階を$6$段階に畳んで逐次律速を緩和し、ボコーダによる波形生成条件で意見値$4.44$と真値$4.52$の近接が報告された\autocite{btd064}。ただし局所忠実度の軸であり、長域や歌詞や動機の証拠ではない\autocite{btd064}。

潜在拡散の転移を掘り下げる。$4$分の$1$圧縮で拡散を軽くし、圧縮率$4$と$8$と$16$で$29.48$と$33.50$と$34.32$が報告された\autocite{btd067}。音響学習・文章サンプリングは対不足を補う転移で、音響のみ混合が$23.31$と$25.79$で上回った\autocite{btd067}。塗り足しや様式転移は編集の課題である\autocite{btd067}。

代償として、文章と音響の整合間隙が残り、抽象説明文は条件づけを損ない、比率$1$や$2$の学習は単一民生画像処理装置では不能である\autocite{btd067}。音楽固有の長域や動機や歌詞の構造は確立していない\autocite{btd067}。忠実度と整合の軸は区別される\autocite{btd064,btd067}。局所忠実度の主張は声素材のボコーダによる波形生成条件に束縛され、楽曲への一般化は確立していない\autocite{btd064}。文章音響の整合間隙は条件符号化の上限であり、抽象説明文の扱いは後継の課題である\autocite{btd067}。民生装置での学習不能条件は、再現性の拘束として記録すべきである\autocite{btd067}。時間条件や理論属性制御の系譜へ接続する\autocite{btd067,btd068}。

\begin{center}
\small
\begin{tabularx}{\textwidth}{@{}p{0.22\textwidth}p{0.26\textwidth}p{0.24\textwidth}X@{}}
\toprule
系統 & 機構 | 報告（同一条件内） | 境界 \\
\midrule
階層的自己回帰\autocite{btd063} & 三種符号化器、上位疎Transformer、窓化・呼び水延長 | 帯域収束$23.0$デシベル、整合は約$24$秒まで | 副歌・旋律記憶なし、新規様式への一般化は弱い \\
意味・音響段階化\autocite{btd065} & 凍結音響・意味・結合音楽文章の三段階、$15$秒刻み継続 | 距離$4.0$、乖離$1.01$、整合$0.51$（音楽 caps 十秒） | 否定・時間順序の失敗継承、模型未公開 \\
一段階符号音楽\autocite{btd066} & 遅延配置単一復号器、文章交叉注意＋色度接頭辞 | 距離$3.1$、乖離$1.28$、整合$0.31$ | 他条件との順位づけは不可、色度は客観指標と交換 \\
波形拡散\autocite{btd064} | 双方向膨張畳み込み、非加重変分目的 | 意見値$4.44$対真値$4.52$、$2.1$倍速 | 局所忠実度のみ、長域・歌詞・動機は未確立 \\
潜在拡散\autocite{btd067} | メル変分$1/4$圧縮、対照整合、音響学習・文章サンプリング | 距離$23.31$、乖離$1.59$ほか（音響 caps） | 整合間隙、抽象文は条件を損なう \\
\bottomrule
\end{tabularx}
\end{center}

\subsection{時間条件・楽理制御と人間選好検証：開放重みと評価境界}
開放重みでの時間条件と楽理属性の精密制御が次の隘路になった\autocite{btd068,btd069}。変化は変分潜在と楽理交叉注意である\autocite{btd068,btd069}。第一に、五畳み込み残差の連続潜在を$21.5$ヘルツ$64$次元とし、文章符号器と時刻接頭辞付きDiffusion Transformer (DiT)で可変長を無音充填で扱い、共有許諾資料のみで学習した\autocite{btd068}。音響 caps 条件で開放距離$78.24$と$103.66$や$101.11$、乖離$2.14$、整合$0.29$と報告された\autocite{btd068}。代償として、共有許諾資料が楽曲品質を制限し、接続詞や可知発話や歌唱に失敗し、英語のみである\autocite{btd068}。第二に、潜在拡散に対照音楽網を重ね、文章と拍と和音の順次交叉注意で拍優先の制御学習を行い、楽理増強資料で補った\autocite{btd069}。検証条件で距離$26.35$や$25.97$や$25.18$、乖離$1.21$や$1.12$や$1.16$、可制御拍色度は$11.64$から$17.99$と$6.61$や$3.97$、専門家音楽性は$5.76$から$6.10$と$2.75$と報告された\autocite{btd069}。ただし拍速度はほぼ同等で、$10$秒断片のみである\autocite{btd069}。和音や調や拍の精密可制御や局所忠実度の符号比較は条件外では確立していない\autocite{btd069}。継続と補完の軸は別に評価する\autocite{btd068,btd069}。

可変長の機構を掘り下げる。上限$47$秒までを無音充填で扱い後に切り詰める\autocite{btd068}。器楽条件の距離$96.51$や信号対歪み$-0.93$と$6.28$、重複監査$3693$と$856$件の非該当確認は同一条件内である\autocite{btd068}。共有許諾条件は再現性を支えるが歌唱に失敗する\autocite{btd068}。

楽理制御の機構を掘り下げる。文章と拍と和音の順次交叉注意を拍優先で学ぶ\autocite{btd069}。楽理増強$52768$組で補い、$10$秒での拍色度$11.64$から$17.99$の変化が報告された\autocite{btd069}。拍速度はほぼ同等で抽出雑音があり、長域の証拠ではない\autocite{btd069}。

共有許諾資料の制約は、開放重みの再現性と品質の交換であり、商用資料との比較は条件をそろえずに行わない\autocite{btd068}。拍速度の同等性は、拍抽出の雑音を考慮して読むべきである\autocite{btd069}。$10$秒断片の主張は、長域形式の証拠ではない\autocite{btd069}。和音や調や拍の精密可制御は、検証条件の枠内に留まる\autocite{btd069}。これらは人間選好検証の対象として継承される\autocite{btd069,btd136}。

自動指標が人間選好と一致するかが最後の隘路になった\autocite{btd136}。変化は大規模人間選好検証である\autocite{btd136}。標識三重の楽曲組合せで$12$模型掛ける$500$の十秒器楽断片、計$6000$曲に$2500$人超の評価者から$15600$の判断を集め、二値選好と整合判断から序列と Bradley Terry 係数で自動指標と対照した\autocite{btd136}。報告は条件付きである。特定商用$3.5$版が二軸で首位、開放第二版が第一版を$5$点上回り、品質相関は音楽音響学習の距離変種が相対的に高く、整合相関は音楽学習の対照変種が相対的に高いとされたが、相関の exact 値は図表束縛で本文数値としては確立しない\autocite{btd136}。対象は$3$や$3.5$世代で$5$や$6$世代ではなく、十秒器楽のみで、供給者薄い領域がある\autocite{btd136}。

人間選好の因果を掘り下げる。余弦$0.1382$以下$500$組で重複を抑え、$12$模型掛ける$500$断片で条件をそろえ、$7800$と$7800$の判断で序列と自動指標を対照する\autocite{btd136}。品質で距離変種、整合で対照変種が相対的に高いが、相関 exact 値は図表束縛で未確立である\autocite{btd136}。$3.5$版の首位や$5$点変化は十秒器楽の条件付き到達点である\autocite{btd136}。

継承すべき教訓は明確である。距離や整合の自動指標は代理水準に留め、閉鎖系の順位づけは人間選好検証を経なければならない\autocite{btd136}。互換性のない持続時間や手順の横断順位は行わない\autocite{btd136}。音響忠実度と文章音楽整合と長域形式と歌詞整合と継続補完は分離し、長編歌曲構造は開放課題として残す\autocite{btd136}。世代の限定は主張の適用範囲の拘束であり、新世代への外挿は行わない\autocite{btd136}。器楽断片の主張は、歌唱や長編の証拠ではない\autocite{btd136}。これが本系譜の評価境界である\autocite{btd136}。

\subsection{生成・継続・編集・制御の分離：課題設定の区別}
生成と継続と編集と制御は別の課題である\autocite{btd065,btd067,btd068,btd069}。生成は新規生成、継続は既存断片の引き継ぎ、編集は潜在の一部保存による塗り足しや様式転移、制御は旋律や拍や和音や時刻の追随である\autocite{btd065,btd067,btd068,btd069}。

継続の系譜をたどる。窓化や呼び水は約$24$秒を越えると漂流する\autocite{btd063}。$15$秒刻みの物語式継続や旋律拡張、$30$秒学習の延長は延長の記述であるが、節や副歌の証拠ではない\autocite{btd065}。継続は前断片との整合で評価し、生成の品質と混ぜない\autocite{btd063,btd065}。

編集の系譜をたどる。覆い潜在の塗り足しや超解像は覆い外を保って覆い内を起こし、浅い逆行の様式転移は中途から様式を変える\autocite{btd067}。長域や動機や歌詞の証拠ではなく、保存と変化で評価する\autocite{btd067}。

制御の系譜をたどる。色度接頭辞は追随を強めるが客観指標と交換になり、時刻接頭辞は時間指定だが楽式維持の証拠ではない\autocite{btd066,btd068}。拍和音の順次交叉注意は拍色度$11.64$から$17.99$の変化が報告されたが、$10$秒断片の枠内に留まる\autocite{btd069}。属性追随で評価し、忠実度や整合と混ぜない\autocite{btd066,btd069}。

\begin{center}
\small
\begin{tabularx}{\textwidth}{@{}p{0.20\textwidth}p{0.30\textwidth}p{0.28\textwidth}X@{}}
\toprule
課題 & 機構（同一条件内） & 報告の読み方 | 混ぜない軸 \\
\midrule
生成\autocite{btd065,btd066} & 文章条件の新規生成 | 十秒等の枠内 | 継続の整合と混ぜない \\
継続\autocite{btd063,btd065} & 窓化・呼び水、$15$秒刻み、旋律拡張 | 約$24$秒まで、以後漂流 | 延長と維持を分ける \\
編集\autocite{btd067} & 塗り足し・超解像、浅い逆行転移 | 保存と変化、一般音響の延長 | 長域等の証拠ではない \\
制御\autocite{btd066,btd068,btd069} & 色度・時刻接頭辞、拍和音交叉注意 | 追随の代理、$10$秒の枠内 | 忠実度・整合と混ぜない \\
\bottomrule
\end{tabularx}
\end{center}

\subsection{持続の延長と構造の維持：時間の二軸}
持続の延長と構造維持は別の軸である\autocite{btd063,btd065,btd068,btd069}。延長は音を続けられるかの記述、維持は節や副歌や動機や調が保たれるかの記述であり、後者は作曲者や genre や歌詞窓条件の範囲外では確立していない\autocite{btd063,btd065}。

時間報告を整理する。約$24$秒文脈と窓延長の漂流\autocite{btd063}、$15$秒刻み継続と$30$秒学習の延長\autocite{btd065}、上限$47$秒の無音充填可変長\autocite{btd068}、$10$秒断片の拍制御や器楽選好\autocite{btd069,btd136}は別の条件であり、横断比較は行わない\autocite{btd063,btd065,btd068,btd069,btd136}。いずれも節や副歌構造の証拠ではない\autocite{btd065,btd068,btd069}。

構造維持は未確立である。副歌反復や旋律記憶はなく\autocite{btd063}、節や副歌の模型はなく否定や時間順序の失敗を継承し\autocite{btd065}、拍制御は$10$秒断片と抽出雑音の拘束を伴う\autocite{btd069}。長編歌曲構造は開放課題として残す\autocite{btd136}。

\subsection{評価の四分割：音質・整合・構造・人間選好}
評価は四つに分割する。音質と整合と構造と選好は別の代理で測られ、自動指標は代理に留まり、閉鎖系の順位づけは選好検証を経る\autocite{btd063,btd065,btd066,btd067,btd068,btd069,btd136}。条件の異なる数値の優劣は論じない\autocite{btd063,btd065,btd066,btd067,btd068,btd136}。

第一に、音質の軸である。帯域収束$23.0$デシベルは再構成の記述に留まる\autocite{btd063}。意見値$4.44$と真値$4.52$はボコーダによる波形生成条件の局所到達点であり、楽曲への一般化はない\autocite{btd064}。距離$23.31$や全体$65.91$は一般音響の潜在忠実度であり、開放距離$78.24$は共有許諾条件付きの到達点である\autocite{btd067,btd068}。距離は分布の隔たりの代理であり、長域や歌詞や動機を測らない\autocite{btd063,btd065,btd136}。

第二に、整合の軸である。乖離$1.01$や整合$0.51$と勝利$312$と$158$や$97$は同一条件内の記述で、意味除去で$1.05$や$0.49$へ変化した\autocite{btd065}。$1.28$や$0.31$や$0.29$は各枠内で比較できない\autocite{btd066,btd068}。否定等の失敗を継承し\autocite{btd065}、抽象文は条件を損なう\autocite{btd067}。

第三に、構造と制御の軸である。類似$0.66$と$0.10$は引き継ぎの代理であり、精密制御には増強の研究が必要である\autocite{btd066}。拍色度$11.64$から$17.99$や音楽性$5.76$から$6.10$は$10$秒断片の代理であり、速度同等性は抽出雑音を考慮する\autocite{btd069}。長域指標や可知歌詞率や動機可制御性は範囲外では確立していない\autocite{btd063,btd065}。

第四に、人間選好の軸である。序列は代理水準の検証手順であり、品質で距離変種、整合で対照変種が相対的に高いとされたが、相関 exact 値は図表束縛で未確立である\autocite{btd136}。$3.5$版の二軸首位や$5$点変化は$3$や$3.5$世代と十秒器楽の条件付き到達点である\autocite{btd136}。閉鎖系の順位づけは検証を経なければならない\autocite{btd136}。

\begin{center}
\small
\begin{tabularx}{\textwidth}{@{}p{0.22\textwidth}p{0.26\textwidth}X@{}}
\toprule
指標 & 測る軸（代理水準） & 測らないもの \\
\midrule
距離（FAD系）\autocite{btd065,btd066} & 音響忠実度の代理 | 長域形式・歌詞整合・動機可制御性は測らない \\
乖離（KL系）\autocite{btd065,btd066} & 分布の隔たりの代理 & 文章音楽整合・可制御性は測らない \\
整合（CLAP系）\autocite{btd065,btd068} & 文章音楽対応の代理 | 否定・時間順序の失敗は継承\autocite{btd065}、抽象文は条件を損なう\autocite{btd067} \\
旋律色度類似\autocite{btd066} & 旋律条件の代理 | 精密制御には音響増強の研究が必要 \\
拍色度・専門家音楽性\autocite{btd069} & 拍優先制御の代理（$10$秒断片） | 拍速度の同等性は抽出雑音を考慮 \\
人間選好・Bradley Terry\autocite{btd136} & 二値選好と整合判断の序列 | 新世代・歌唱・長編への外挿は不可、相関 exact 値は未確立 \\
\bottomrule
\end{tabularx}
\end{center}

\paragraph{条件内比較の作法}
本節の数値はすべて同一条件内の記述であり、他条件との順位づけには用いない\autocite{btd066}。音楽 caps 条件での距離$3.1$と音響 caps 条件での距離$23.31$は、母集団も手順も異なるため直接比較できない\autocite{btd066,btd067}。開放距離$78.24$は共有許諾資料という条件付きの到達点であり、商用資料の条件とは切り分けて読む\autocite{btd068}。人間勝利$312$と$158$や$97$、人間評価$78.43$と$81.11$から$84.81$へ、可制御拍色度$11.64$から$17.99$、専門家音楽性$5.76$から$6.10$は、いずれもその検証条件の枠内の記述である\autocite{btd065,btd066,btd069}。条件をそろえない比較は、本節では成立しない比較として扱う\autocite{btd063,btd065,btd066,btd067,btd068,btd069,btd136}。互換性のない持続時間や手順の横断順位は行わない\autocite{btd136}。

\paragraph{開放重みと再現性の交換}
共有許諾資料のみの学習は再現性を支えるが楽曲品質を制限し、接続詞や可知発話や歌唱に失敗し英語のみである\autocite{btd068}。開放距離の報告は音響 caps 条件の枠内に留まる\autocite{btd068}。商用資料との比較は条件をそろえずに行わない\autocite{btd068}。民生装置での学習不能条件（比率$1$や$2$）は、再現性の拘束として記録すべきである\autocite{btd067}。開放線の主張は、共有許諾という資料条件付きの到達点である\autocite{btd068}。

\paragraph{長編歌曲の開放課題}
節や副歌の模型はなく、窓延長では漂流が生じる\autocite{btd063,btd065}。$30$秒学習の自己回帰延長は節・副歌構造の証拠ではない\autocite{btd065}。$10$秒断片の主張は長域形式の証拠ではなく\autocite{btd069}、器楽断片の主張は歌唱や長編の証拠ではない\autocite{btd136}。音響忠実度と文章音楽整合と長域形式と歌詞整合と継続補完は分離し、長編歌曲構造は開放課題として残す\autocite{btd136}。記憶再現の低さは学習資料の扱いの報告であり、品質の保証ではない\autocite{btd065}。

\paragraph{歌詞と継続の分離}
歌詞の扱いと継続・補完は、別の軸として評価する\autocite{btd068,btd069}。共有許諾資料の条件では歌唱に失敗し、英語のみという制約が残る\autocite{btd068}。$15$秒刻みの物語式継続や旋律条件拡張は、文章条件を音響条件に置換する推論の記述であり、節・副歌構造の証拠ではない\autocite{btd065}。継続と補完の軸は生成条件とは別に評価し、和音や調や拍の精密可制御は検証条件の枠内に留める\autocite{btd068,btd069}。歌詞整合の可知率や動機水準の可制御性は、作曲者・genre・歌詞窓条件の範囲外では確立していない\autocite{btd063,btd065}。長編歌曲構造は開放課題として残す\autocite{btd136}。

\subsection*{読解上の境界}
互換性のない持続時間や手順の横断順位は行わず、距離・対照整合指標の限界を保持する\autocite{btd063,btd065,btd066,btd067,btd068,btd136}。整合は約$24$秒の上位文脈までで窓延長では漂流し、繰返し副歌や旋律記憶はなく、新規様式への一般化は弱い\autocite{btd063}。長域形式指標・可知歌詞率・動機水準可制御性は、作曲者・ genre・歌詞窓条件の範囲外では確立していない\autocite{btd063,btd065}。結合符号の否定・時間順序の失敗を継承し、$30$秒学習の自己回帰延長で節・副歌模型はなく模型公開もない\autocite{btd065}。拍・和音・調の精密可制御と局所忠実度の符号比較は条件外では確立していない\autocite{btd069}。微細制御は分類器なし誘導のみに依存し、平坦化は$6000$対$1500$段階の交換を伴い、色度条件は客観指標を劣化させる\autocite{btd066}。歌詞歌唱や楽理制御は教師なし色度旋律類似の範囲を超えては確立していない\autocite{btd066,btd069}。共有許諾資料は楽曲品質を制限し、接続詞・可知発話・歌唱に失敗し英語のみである\autocite{btd068}。波形拡散は局所忠実度に留まり長域形式・歌詞整合・動機可制御性は確立していない\autocite{btd064}。対照の文章・音響間隙が残り、抽象説明文は文章条件を損ない、比率$1$・$2$は単一民生画像処理装置で学習不能である\autocite{btd067}。拍速度はほぼ同等で$10$秒断片のみである\autocite{btd069}。人間選好検証は$3$・$3.5$世代で$5$・$6$世代ではなく十秒器楽のみであり、新世代模型と相関 exact 値は本稿では確立しない\autocite{btd136}。

\section{映像の系譜：時間的一貫性からメディア基盤へ}\sectionkicker{VIDEO}\label{sec:video}

\begin{multicols}{2}
\raggedcolumns
映像生成の系譜は、画像生成に時間を足した付録ではない。時間的一貫性、同一性保持、音声同期、長域物語構造、継続と編集保存という独立の評価軸を持ち、各軸の主張は条件と版と検査点に束縛される\autocite{btd070,btd071}。本節は、分解型拡散と多重超解像の連鎖から資料効率化、厳選潜在動画と動作接続器、閉鎖頂点と開放混合専門家を経て物体水準編集基準測定へ至る系譜を、忠実度と整合と同一性と同期と長域構造の分離によりたどる\autocite{btd070,btd071,btd073}。各段階の共通の読み方は、対応上限時間と実証された整合を混同しないことである。最大対応時間は実証された整合ではなく、版と検査点の通貨拘束が必須である\autocite{btd070,btd071,btd073,btd074,btd075,btd077,btd124}。

系譜の見取り図は次のとおりである。第一に、二次元畳み込みの膨張と三元分解による分解型動画拡散と、凍結文章符号器を起点とする多重超解像の連鎖が、画像拡散に時間整合を加えて高品位長尺化する出発点になった\autocite{btd070,btd071}。第二に、文章画像の外観整合を凍結し運動のみを無標識動画で学ぶ資料効率化が、対文章動画資料の不足に応答した\autocite{btd073}。第三に、動作接続器と厳選潜在動画が、個人化画像模型の調整なし動作化と厳選資料・段階学習の統一を進めた\autocite{btd075,btd074}。第四に、閉鎖頂点と開放重みの混合専門家配置が能力と workflow の分岐を示し、物体水準編集基準測定が一段階編集から benchmark 測定への転換を定めた\autocite{btd077,btd124,btd138}。版と検査点と母集団の三点拘束を定めておく。開放線は$2.2$で凍結し後継は応用接続のみであり、主張の適用範囲はその版に束縛される\autocite{btd124}。閉鎖頂点の主張は消費本文の頁情報に依拠し、製品終了の lifecycle 事実とともに記録する\autocite{btd077}。報告値は母集団と条件に束縛され、厳選千万部分集合の優位も母集団の偏りを解消しない\autocite{btd075}。連合学習の改善は画像枚数の条件に束縛され、動画資料の規模効果とは区別する\autocite{btd070}。以下ではこの順序で詳述する。

証拠の扱いに関する留保を先に記す。閉鎖頂点の扱いは能力と workflow と lifecycle の証拠に限り、非公開の骨格や学習や符号化の推測は行わない\autocite{btd077}。開放重み系の供給者報告の伸び率は帰属付きで扱い、独立再現なしに順位づけを行わない\autocite{btd124}。基準測定の報告は編集領域内に限り、領域外の優劣主張は確立しない\autocite{btd138}。純粋生成の除去検証の薄さは、領域外への一般化の禁止を意味する\autocite{btd138}。
空間潜在と時間圧縮を区別して読む。画素区画の学習と潜在への挿入と高圧縮符号化では、記憶と計算の包絡が異なる\autocite{btd070,btd071,btd075,btd124}。二次元に時間を足す因子分解とTransformer注意の挿入は設計差として区別し、画像事前分布の転用は有益性と限界を分離し、枠忠実度と運動品質と時間整合は別軸として測る\autocite{btd070,btd071,btd073,btd074,btd075,btd138}。短断片の実証を長域構造や同一性保持の証拠としない\autocite{btd070,btd071,btd073,btd074,btd075}。閉鎖頂点は能力と workflow と lifecycle の証拠に限り、機構の推測は行わない\autocite{btd077}。各遷移は隘路と機構と変化と改善と代償と継承の順序で読む\autocite{btd070,btd071,btd073,btd074,btd075,btd077,btd124,btd138}。

\end{multicols}

\subsection{分解型拡散と多重超解像：時間整合の出発点}
画像拡散に時間整合を加えて高品位長尺化する隘路が出発点である\autocite{btd070}。変化は分解型動画拡散から多重超解像の連鎖である\autocite{btd070,btd071}。第一に、二次元畳み込みを$1 \times 3 \times 3$にし、枠別空間注意と時間注意と相対位置で三元を分解し、画像は時間注意を遮蔽して連合学習し、置換や補間や空間超解像や区画自己回帰延長を再構成誘導で実現した\autocite{btd070}。報告は条件付きで、無条件で距離$295$や指標$57$と真値$60.2$、予測で距離$68.19$や$66.92$、$600$種で$18.6$、画像八枚の連合で$202$から$57.84$へ進んだ\autocite{btd070}。ただし$16$枠区画のみで、長時間は自己回帰や時間超解像の誘導に依存し、資料偏りを反映し、模型は未公開である\autocite{btd070}。第二に、凍結大規模文章符号器に基盤と空間超解像三段と時間超解像三段の計$11.6$Bで$1280 \times 768$の$128$枠毎秒$24$を狙い、速度予測と雑音条件増強で学習し、振動誘導と二段蒸留で$8 \times 8$段階化した\autocite{btd071}。蒸留後は整合$25.03$と$25.19$で約$18$倍速と報告された\autocite{btd071}。

代償として連鎖費用は重く、三次元回転整合は不精密で、偏り危険を伴い未公開である\autocite{btd071}。対応上限時間は実証された整合を意味しない\autocite{btd070,btd071}。局所忠実度と短域整合は区別される\autocite{btd070,btd071}。連合学習の改善は画像枚数の条件に束縛され、動画資料の規模効果とは区別する\autocite{btd070}。蒸留の高速化は誘導の振動操作と二段構成に依存し、一般のサンプリングへの転用は確立していない\autocite{btd071}。未公開模型の主張は本文数値の範囲に留め、再現性の検証は後継の開放系に委ねられた\autocite{btd070,btd071}。対応上限時間は実証された整合を意味しない\autocite{btd070,btd071}。資料効率化の系譜へ継承される\autocite{btd070,btd071,btd073}。

$64$枠延長では再構成誘導が$136.22$や$134.55$を示し、置換の$451.45$や$436.16$との差が報告された\autocite{btd070}。対文章動画の連合では画像なしの$202.28$や$205.42$から画像八枚で$57.84$や$60.72$への変化が条件付きで報告された\autocite{btd070}。多重超解像では基盤拡大の$500$Mから$1.6$Bさらに$5.6$Bが$4096$標本の条件で整合と距離を動かし、$256$と$128$段階の$25.19$や$618$秒に対し蒸留後の$8$と$8$段階が$25.03$や$35$秒を示した\autocite{btd071}。いずれも対応上限時間の記述であり、実証された長域整合ではない\autocite{btd070,btd071}。

\subsection{資料効率化：外観凍結と運動学習の分離}
対文章動画資料の不足が次の隘路になった\autocite{btd073}。変化は文章画像の外観整合を凍結し、運動のみを無標識動画で学ぶ資料効率化である\autocite{btd073}。文章画像多層網に擬似三次元畳み込みと注意を拡張し、毎秒枠数条件を加え、外観事前分布を固定したまま動画集合だけで運動を学習し、事前分布から時間復号、遮蔽補間、時間空間超解像へ共有雑音で進む\autocite{btd073}。報告はゼロショット条件で、$59794$全文脈の距離$13.17$と整合$0.3049$で$23.59$や$0.2631$を上回り、分別条件の指標$33.00$を示した\autocite{btd073}。生成と編集の区別では、純粋生成の主張に留まり、編集保存の定量化は含まない\autocite{btd073}。

代償として、動画でのみ観測される現象の文章化は学習不能で、長時間多場面や物語動画は将来課題とされ、網資料の社会偏りは誇張される\autocite{btd073}。短断片中心であり、一括要点標本は高費用で、運動不足や低速高記憶サンプリングの問題が残る\autocite{btd073}。長時間化には粗い連鎖や動画符号化器や蒸留が必要とされる\autocite{btd073}。対応上限時間は実証された整合を意味しない\autocite{btd073}。外観と運動の分離は、個人化模型の動作化や厳選潜在動画の系譜へ継承される\autocite{btd073,btd074,btd075}。短断片中心の主張は、長域構造の証拠ではない\autocite{btd073}。共有雑音による多段進行は、時間方向の滑らかさと外観保存の両立を狙ったものであり、長時間外挿の保証ではない\autocite{btd073}。社会偏りの誇張は資料選別の課題として記録され、軽減策の定量化は残されていない\autocite{btd073}。対応上限時間は実証された整合を意味しない\autocite{btd073}。

画像対は$2.3$Bの選別済み集合で学び、事前分布を凍結したまま無標識動画だけで運動を学ぶ\autocite{btd073}。推論は事前分布から$16 \times 64 \times 64$の時間復号へ進み、遮蔽補間で$76$枠化し、時間と空間の超解像を共有雑音で進める\autocite{btd073}。全文脈$59794$件のゼロショットで距離$13.17$と整合$0.3049$が$23.59$や$0.2631$と比較され、微調整後は$82.55$と$81.25$が条件付きで報告された\autocite{btd073}。動画でのみ観測される現象の文章化は学習不能であり、長多場面や物語動画は将来課題として残る\autocite{btd073}。

\subsection{因子分解の設計差：二次元＋時間と時空間とTransformer}
分解型拡散は二次元畳み込みを$1 \times 3 \times 3$に膨張させ、枠別空間注意と時間注意と相対位置で三元を分解し、画像は時間注意を遮蔽して連合学習する\autocite{btd070}。多重超解像は基盤に時間注意を置き超解像段階に時間畳み込みを置く分担であり、速度予測と雑音条件増強と振動誘導と二段蒸留を組み合わせる\autocite{btd071}。資料効率化は擬似三次元畳み込みと注意を恒等や零初期値で拡張し毎秒枠数条件を加え、厳選潜在動画は空間初期値に時間畳み込みと注意を挿入し、動作接続器は時間Transformer注意を正弦位置と零初期値残差で挿入する\autocite{btd073,btd075,btd074}。時間畳み込みの除去検証では同一枠化への崩壊が報告された\autocite{btd074}。いずれも短区間の整合を狙った設計差であり、長域への外挿の保証ではない\autocite{btd070,btd071,btd073,btd074,btd075}。

この軸の遷移は、隘路の時間整合の欠如と資料不足から、機構の分解型三元注意と連鎖と擬似三次元拡張により、遮蔽連合から振動誘導と蒸留と共有雑音多段へ進み、$68.19$や$18.6$や約$18$倍速や$13.17$と$0.3049$を条件付きで示し、区画限定を代償として潜在挿入へ継承される\autocite{btd070,btd071,btd073,btd075}。

\subsection{潜在と圧縮の軸：空間潜在と時間圧縮}
分解型拡散は$16$枠区画を$64 \times 64$や$128 \times 128$で学び長時間は誘導に依存し、多重超解像は$1280 \times 768$の$128$枠毎秒$24$を連鎖で狙い画素側の費用が重い\autocite{btd070,btd071}。厳選潜在動画は空間初期値に時間挿入を行う潜在側の統一であり$14 \times 256 \times 384$から$14 \times 320 \times 576$へ進み、動作接続器は五次元膨張潜在で基盤再学習なしに推論時尺度で動作を付与する\autocite{btd075,btd074}。開放重み系は変分符号化を$4 \times 16 \times 16$の$64$倍圧縮に置き区分化で$4 \times 32 \times 32$とし、高低騒音二専門家を信号対雑音比の閾値で切り替える計$27$B有効$14$Bと密な$5$B模型を併置する\autocite{btd124}。本文を超える資料規模や演算量の定量化は含まれない\autocite{btd124}。

民生画像処理装置の単一$24$GB九分未満での五秒$720$P毎秒$24$は配備の記述であり品質順位ではなく、大容量配置は単一$80$GBや複数装置を要する\autocite{btd124}。要点一括標本は高費用であり、蒸留の短縮は振動操作と二段構成に依存する\autocite{btd073,btd075,btd071}。編集側の毎枠$3.07$秒と$1.44$秒や$25.98$秒は領域内条件の費用記述である\autocite{btd138}。開放線は$2.2$で凍結し後継は応用接続のみである\autocite{btd124}。

\begin{center}
\small
\begin{tabularx}{\textwidth}{@{}p{0.20\textwidth}p{0.30\textwidth}X@{}}
\toprule
潜在・圧縮 & 機構・条件 | 包絡・境界 \\
\midrule
画素区画\autocite{btd070} & $16$枠画素、遮蔽連合と再構成誘導 | 長時間は誘導依存、未公開 \\
多重超解像\autocite{btd071} & 凍結符号器、基盤＋空間三段＋時間三段 | 連鎖費用は重い、未公開 \\
外観凍結\autocite{btd073} & $2.3$B凍結、運動のみ無標識学習 | 短断片中心、長多場面は将来課題 \\
潜在挿入\autocite{btd075} & 空間初期値＋時間挿入、三段階 | 短断片中心、低速高記憶サンプリング \\
推論時接続\autocite{btd074} & 五次元膨張、時間Transformer注意 | 約$50$参照、深度実演は乱雑音のみ \\
高圧縮符号化\autocite{btd124} & $4 \times 16 \times 16$と区分化、二専門家 | $24$GB九分未満は配備記述 \\
\bottomrule
\end{tabularx}
\end{center}

\subsection{動作接続器と厳選潜在動画：個人化の動作化と段階学習}
個人化画像模型を調整なしで動作化し、厳選資料と段階学習を統一することが次の隘路になった\autocite{btd075,btd074}。変化は動作接続器と厳選潜在動画である\autocite{btd075,btd074}。第一に、空間初期値に時間畳み込みと注意を挿入し、画像事前学習と低解像動画事前学習と高品質微調整の三段階で学び、撮影 LoRA 変種を備えた\autocite{btd075}。ゼロショット条件で距離$242.02$と$367.23$を上回り、人間序列では画像事前学習が無学習を、厳選千万部分集合が網全体を上回り、微調整後も優位が残った\autocite{btd075}。ただし短断片中心で、低速高記憶サンプリングが残る\autocite{btd075}。第二に、凍結画像模型に時間Transformer注意と領域接続器と運動 LoRA を挿入し、基盤再学習なしで推論時尺度で動作を付与した\autocite{btd074}。平滑性評価は$2.825$と$1.615$を上回り、文章や領域や平滑の三軸で比較された\autocite{btd074}。

代償として、動画資料の品質間隙には接続器が必要で、運動 LoRA は約$50$参照を要し、個人化模型の誤用危険を継承する\autocite{btd074}。深度制御の実演は乱雑音のみである\autocite{btd074}。局所忠実度と短域整合と同一性保持は区別され、分数分超の同一性や長域物語や同期音や継続や編集保存は確立していない\autocite{btd074,btd075}。厳選資料の優位は微調整後も残ることが示されたが、母集団の偏りは解消されない\autocite{btd075}。接続器の参照要求は実用上の制約であり、少数参照化は後継の課題である\autocite{btd074}。開放混合専門家配置の系譜へ接続する\autocite{btd074,btd075,btd124}。

\begin{center}
\small
厳選潜在動画は画像符号化と雑音付与枠の連結と線形増加誘導で単一画像から進み、撮影変種を切り替える\autocite{btd075}。動作接続器は文章や領域や平滑の三軸で比較され、平滑で$2.825$と$1.615$の対比が同一個人化模型の条件で報告された\autocite{btd074}。適合器は推論時尺度で係数を操作し零で除去でき、運動 LoRA は約$50$参照で成立し$5$参照では質感化への劣化が報告された\autocite{btd074}。人間序列の優位は微調整後も残ったが、母集団の偏りの解消ではない\autocite{btd075}。

\begin{tabularx}{\textwidth}{@{}p{0.22\textwidth}p{0.26\textwidth}p{0.24\textwidth}X@{}}
\toprule
段階 & 機構 | 報告（条件付き） | 境界 \\
\midrule
分解型拡散\autocite{btd070} & $1 \times 3 \times 3$膨張、三元分解、遮蔽連合学習、再構成誘導 | 距離$295$・指標$57$、連合$202\to57.84$ | $16$枠区画のみ、長時間は誘導依存、未公開 \\
多重超解像\autocite{btd071} | 凍結文章符号器、基盤＋空間三段＋時間三段、振動誘導＋二段蒸留 | 整合$25.03$・$25.19$、約$18$倍速 | 連鎖費用は重い、回転整合は不精密、未公開 \\
資料効率化\autocite{btd073} | 外観凍結、運動のみ無標識学習、共有雑音の多段進行 | 距離$13.17$・整合$0.3049$（ゼロショット） | 短断片中心、長多場面・物語は将来課題 \\
厳選潜在動画\autocite{btd075} | 時間挿入、三段階学習、撮影 LoRA | 距離$242.02$、厳選千万が網全体を上回る | 短断片中心、低速高記憶サンプリング \\
動作接続器\autocite{btd074} | 時間Transformer注意・領域接続器・運動 LoRA、推論時尺度 | 平滑$2.825$対$1.615$、三軸比較 | 約$50$参照、深度実演は乱雑音のみ \\
\bottomrule
\end{tabularx}
\end{center}

\subsection{閉鎖頂点・開放混合専門家と編集基準測定：基盤への到達}
閉鎖頂点と開放重みの配備が次の分岐になった\autocite{btd077,btd124}。扱いは能力と workflow と lifecycle の証拠に限り、非公開の骨格や学習や符号化の推測は行わない\autocite{btd077}。第一の閉鎖頂点は約一分の文章動画を品質と指示遵守で示し、歴史的転換点となったが、消費本文の頁情報に依拠し、機構の主張は確立しない\autocite{btd077}。製品終了の lifecycle 事実とともに記録し、対応上限時間は実証された整合を意味しないことを明示する\autocite{btd077}。第二の開放重み系は、文章動画と画像動画の高低騒音二専門家で計$27$B有効$14$Bを、SNR 閾値で切替え、高圧縮変分符号化の$4 \times 16 \times 16$に patch 化を重ね、密な統合模型で民生画像処理装置の単一$24$GB九分未満の五秒$720$P毎秒$24$を達成した\autocite{btd124}。供給者報告の$65.6\%$と$83.2\%$の伸びは帰属付きで扱う\autocite{btd124}。開放線は$2.2$で凍結し、後継は応用接続のみである\autocite{btd124}。版と検査点の通貨拘束が必須で、頁の弱点表にある物理や因果や左右や軌跡や自発実体や変形や多人物の失敗は定性的記録に留める\autocite{btd124}。配備経済の軸は別章に委ねる\autocite{btd124}。

物体水準の微細編集を標準化して測ることが最後の隘路になった\autocite{btd138}。変化は一段階編集から benchmark 測定への転換である\autocite{btd138}。先行の一段階編集は疎時空間注意の膨張模型で枠間整合$92.40$を示したが、重なり混合や漂流やちらつきの除去検証に留まった\autocite{btd138}。新しい benchmark は百本の動画と$420$組の指示対で色彩や材質や置換を剛体と非剛体に分け、$15$指標で測った\autocite{btd138}。報告は編集領域内に限る\autocite{btd138}。開放編集系は距離$94.61$と類似$82.55$と尖頭$25.57$で忠実度$89.43$、毎枠$3.07$秒と$1.44$秒や$25.98$秒、精度$46.97$と$43.84$や$48.42$とされた\autocite{btd138}。純粋生成の除去検証は薄く、領域外の優劣主張は確立しない\autocite{btd138}。

継承すべき評価原則は明確である。局所忠実度と短域整合と同一性保持と音声同期と長域構造と継続と編集保存は分離し、単一良標本を長域能力の証拠としてはならない\autocite{btd138}。最大対応時間は実証された整合ではなく、版と条件と母集団の拘束なしの順位づけは行わない\autocite{btd138}。純粋生成の除去検証の薄さは、領域外への一般化の禁止を意味する\autocite{btd138}。編集領域内の報告は、生成品質の主張とは区別して読む\autocite{btd138}。一段階編集の短区間主張は、基準測定の基礎として継承される\autocite{btd138}。これが映像系譜の評価境界である\autocite{btd138}。

\begin{center}
\small
閉鎖頂点は約一分の文章動画を品質と指示遵守で示し歴史的転換点として記録されるが、頁は数値や条件付き得点を与えず定性的主張と日付付き lifecycle 事実のみであり、機構の推測は行わない\autocite{btd077}。対応上限時間の約一分は実証された長域整合を意味しない\autocite{btd077}。音声画像姿勢の$14$Bや動画化と置換の$14$Bは能力記述として読み、同期品質の順位づけは行わない\autocite{btd124}。供給者報告の$65.6\%$と$83.2\%$の伸びは帰属付きで扱い、独立再現なしに順位づけを行わない\autocite{btd124}。

\begin{tabularx}{\textwidth}{@{}p{0.20\textwidth}p{0.28\textwidth}X@{}}
\toprule
評価軸 & 意味 | 未確立事項 \\
\midrule
局所忠実度\autocite{btd070,btd071} & 短区間の見た目と整合 | 長域への外挿は不可 \\
短域整合\autocite{btd070,btd073} | 枠間・時間方向の滑らかさ | $64$枠超の同一性・物体永続は未確立 \\
同一性保持\autocite{btd074,btd075} | 人物・物体の一貫 | 分数分超は未確立 \\
音声同期\autocite{btd074,btd075} | 音と映像の対応 | 未確立 \\
長域構造\autocite{btd073} | 物語・多場面の構成 | 将来課題 \\
継続・補完\autocite{btd074,btd075} | $5.3$秒超の延長 | 未定量化 \\
編集保存\autocite{btd138} | 改変時の既存内容の保持 | 編集領域内のみ、生成への一般化は不可 \\
\bottomrule
\end{tabularx}
\end{center}

\subsection{短断片の壁と長域・同一性・同期の未確立領域}
学習区画は$16$枠や$14$枠や$25$枠の短断片中心であり、$128$枠毎秒$24$の約$5.3$秒や$76$枠補間や$64$枠延長が対応上限時間の記述である\autocite{btd070,btd071,btd073,btd075}。自己回帰延長と時間超解像の誘導、振動誘導と二段蒸留、粗い連鎖や蒸留の組合せは短区間の延長であり、長域物語や分数分超の同一性の証拠にはならない\autocite{btd070,btd071,btd073}。$64$枠超の同一性や物体永続、$5.3$秒超の継続、編集保存は確立していない\autocite{btd070,btd071,btd074,btd075}。音声同期は数値を与えられず、音声画像姿勢の記述は条件の存在の記述であり同期品質の定量化ではない\autocite{btd077,btd124}。単一良標本を長域能力の証拠としない\autocite{btd138}。

\subsection{参照・条件づけ・編集と基準測定の配置}
文章動画の連合条件づけ、振動誘導、毎秒枠数条件と微小条件、画像符号化と雑音付与枠の連結と線形増加誘導、領域接続器と運動 LoRA と推論時尺度、文章拡張の hosted や局所利用が条件の配置である\autocite{btd070,btd071,btd073,btd075,btd074,btd124}。基準測定は百本と$420$組で色彩や材質や置換を剛体と非剛体に分け$15$指標で測り、実写$74$本の八枠ごと抽出と合成$26$本の$35$から$126$枠を覆いと視覚言語判定で測る\autocite{btd138}。訓練や反転なしの適合の条件で距離$94.61$と類似$82.55$と尖頭$25.57$で忠実度$89.43$、毎枠$3.07$秒と$1.44$秒や$25.98$秒、精度$46.97$と$43.84$や$48.42$が報告された\autocite{btd138}。報告は編集領域内に限り、純粋生成の検証の薄さは領域外への一般化の禁止を意味する\autocite{btd138}。

\paragraph{外観と運動の分離の意義}
資料効率化の外観凍結は、文章画像の整合を保ったまま運動だけを学ぶ分離であり、対文章動画資料の不足への応答である\autocite{btd073}。この分離は個人化模型の動作化や厳選潜在動画へ継承され、基盤再学習なしで推論時尺度で動作を付与する接続器の前提になった\autocite{btd073,btd074,btd075}。ただし分離は万能ではない。動画でのみ観測される現象の文章化は学習不能であり、長時間多場面や物語動画は将来課題として残る\autocite{btd073}。厳選千万部分集合が網全体を上回る報告は、資料選別の効果の記述であり、母集団の偏りの解消ではない\autocite{btd075}。微調整後も優位が残ることと偏りが残ることは、両立する記述として読む\autocite{btd075}。運動 LoRA の約$50$参照の要求は、実用上の制約であり、少数参照化は後継の課題である\autocite{btd074}。

\paragraph{単一良標本の扱い}
単一良標本を長域能力の証拠としてはならない\autocite{btd138}。深度制御の実演は乱雑音のみであり、一般条件の証拠ではない\autocite{btd074}。閉鎖頂点の頁情報は定性的主張と日付付き lifecycle 事実に留まり、機構の主張は確立しない\autocite{btd077}。頁の弱点表にある物理・因果・左右・軌跡・自発実体・変形・多人物の失敗は、定性的記録として読む\autocite{btd124}。純粋生成の除去検証の薄さは、領域外への一般化の禁止を意味する\autocite{btd138}。

\paragraph{配備と評価の分離}
民生画像処理装置の単一$24$GB九分未満での五秒$720$P毎秒$24$は配備の記述であり、品質順位ではない\autocite{btd124}。供給者報告の$65.6\%$と$83.2\%$の伸びは帰属付きで扱い、独立再現なしに順位づけを行わない\autocite{btd124}。配備経済の軸は別章に委ねる\autocite{btd124}。要点一括標本の高費用や低速高記憶サンプリングは、運用条件の記述として読む\autocite{btd073,btd075}。版と条件と母集団の拘束なしの順位づけは行わない\autocite{btd138}。

\paragraph{長時間化の処方と限界}
長時間化の処方は複数あるが、いずれも実証された整合の保証ではない\autocite{btd070,btd071,btd073}。区画自己回帰延長と時間超解像の誘導、振動誘導と二段蒸留、粗い連鎖や動画符号化器や蒸留の組合せは、短区間の整合を延長する手続きであり、長域物語構造や分数分超の同一性の証拠にはならない\autocite{btd070,btd071,btd073}。要点一括標本は高費用で運動不足を伴い、低速高記憶サンプリングの問題が残る\autocite{btd073,btd075}。共有雑音による多段進行は時間方向の滑らかさと外観保存の両立を狙ったものであり、長時間外挿の保証ではない\autocite{btd073}。対応上限時間の記述を実証された整合として読むことは、本節の読み方に反する\autocite{btd070,btd071,btd073,btd074,btd075,btd077,btd124,btd138}。

\paragraph{遷移の読み方：隘路から継承まで}
第一の遷移は時間整合の欠如と資料不足を隘路とし、分解型三元注意と連鎖と外観凍結を機構として、遮蔽連合から振動誘導と蒸留と共有雑音多段へ変化し、条件付きの距離や整合で改善を示し、区画限定と連鎖費用を代償として潜在挿入へ継承される\autocite{btd070,btd071,btd073}。第二の遷移は調整なし動作化と段階学習の統一と固定費用の配備を隘路とし、時間Transformer注意と三段階学習と二専門家と高圧縮符号化を機構として、平滑三軸や厳選千万の序列や帰属付き伸び率で改善を示し、参照要求と短断片中心を代償として配備可能な動画系へ継承される\autocite{btd074,btd075,btd124}。第三は閉鎖頂点を workflow の証拠として継承し機構は確立しない\autocite{btd077}。第四は編集の標準化を隘路とし覆いと視覚言語判定を機構として領域内$15$指標で改善を示し検証の薄さを代償として領域内の対比に継承される\autocite{btd138}。第五は軸分離の要請として継承する\autocite{btd138}。

\subsection*{読解上の境界}
最大対応時間は実証された整合ではなく、版と検査点の通貨拘束が必須である\autocite{btd070,btd071,btd073,btd074,btd075,btd077,btd124,btd138}。分解型拡散は$16$枠区画のみを学習し、長時間は自己回帰・時間超解像の誘導に依存し、資料偏りを反映し未公開である\autocite{btd070}。短域運動忠実度のみが示され、$64$枠超の同一性・物体永続、長域物語構造、音声同期、編集保存は確立していない\autocite{btd070}。連鎖費用は重く回転整合は不精密で、$14$M動画文章＋$60$M画像文章＋LAION-$400$Mの学習に偏り危険を伴い未公開である\autocite{btd071}。分数分超の同一性・長域物語・音声同期・$5.3$秒超の継続・編集保存は定量化されていない\autocite{btd074,btd075}。動画でのみ観測される現象の文章化学習は不能で、長多場面・物語動画は将来課題であり、網資料の社会偏りは誇張される\autocite{btd073}。要点一括標本は高費用で運動不足や低速高記憶サンプリングが残り、長時間化には粗い連鎖や動画符号化器や蒸留が必要である\autocite{btd073}。動画資料の品質間隙には接続器が必要で、運動 LoRA は約$50$参照を要し、個人化模型の誤用危険を継承し、深度制御実演は乱雑音のみである\autocite{btd074}。閉鎖頂点の頁は数値・条件スコアを与えず定性的主張と日付付き lifecycle 事実のみであり、弱点表の物理・因果・左右・軌跡・自発実体・変形・多人物の失敗は定性的記録に留める\autocite{btd077}。供給者報告の伸び率は帰属付きで扱い、独立再現なしに順位づけを行わない。開放線は$2.2$で凍結している\autocite{btd124}。基準測定は編集領域のみで純粋生成の検証は薄く、生成一般への優劣主張は確立しない\autocite{btd138}。

対応上限時間の記述と実証された整合の区別が本節の軸であり、版と検査点と母集団の拘束なしに順位づけは行わない\autocite{btd070,btd071,btd073,btd074,btd075,btd077,btd124,btd138}。

\section{長時間・同期・編集：標本の鮮明さを超える時間軸}
\label{sec:temporal}
\sectionkicker{TEMPORAL}

\begin{multicols}{2}
\raggedcolumns
一枚の画像の鮮明さと、数十秒にわたって破綻なく続く映像の正しさは、別の軸の良さである。本節が扱うのは、その別軸をどう測り、どう読み誤らないかという問題である\autocite{btd131,btd135}。別軸の立て方自体が条件付きである。物理的妥当性は二値判定の粗さ付き、口唇同期は話者正面と評価器版固定付き、対話持続は区切り手順と判定模型依存付きである。三者の別軸は互いに異なり、一つの時間軸に畳めない。畳まずに並置することが、本節の構成の理由である。対象は四つの証拠群である。物理的妥当性を問う映像評価の手順、口唇同期の訓練信号と評価指標の分岐、唇読み特徴に基づく同期専門家の置換、そして音声対話の多ターン評価である\autocite{btd131,btd129,btd130,btd135}。いずれにも共通するのは、単発の標本品質を上げた技術が、そのまま長時間の整合を保証しないという事実である。生成器が一枚をきれいに描けることと、物体の振る舞いが数秒後も筋が通っていることは、別の検証を要する。同期が合っていることと、物語の構造が保たれることも同様である\autocite{btd131,btd130}。編集の観点からも同じ分離が要る。改変後の保存目標、局所性、参照制約、失敗様式は、生成の忠実度とは別の技術問題であり、本節の時間軸はその一部を担う。編集の系譜の本体は別節に譲り、本節では時間的一貫性と同期と対話の持続という範囲で扱う。範囲の限定は弱さではなく、主張の適用域の宣言である。本節では、報告条件と評価器の版と母集団を束ねたまま引用し、条件をまたぐ横並びの順位付けを行わない。長時間の主張は長時間の手順で確かめられた範囲でのみ受け取る。それが本節の読みの作法である\autocite{btd131,btd135}。

\subsection{物理的妥当性を二値で問う手順とその天井}
映像生成の物理的妥当性を問う出発点として、VideoPhyと呼ばれる評価手順は、千件の候補記述から人手検証を経て六八八件の評価用記述を整えた経緯を持つ。三段階の選別では、まず生成器による記述候補を集め、次に人手で検証し、さらに図学を専門とする評価者が易と難に分割した\autocite{btd131}。千件から六八八件への絞り込みは、記述の質を揃えるための手順であり、生成器の能力を測るための手順ではない。選別の段階と評価の段階を混ぜて、選別の厳しさを生成の難しさとして語らない。易と難の分割も専門評価者の判断であり、物理法則の分類ではないことを添えて引用する。対象は固体同士が二八九件、固体と流体が二九一件、流体同士が一〇八件に三分され、百三十八種の動作と平均八・五語の短い指示で構成される\autocite{btd131}。三分の読み方も条件付きである。固体同士と固体流体と流体同士の三分は、相互作用の難所を分けて見るための設計であり、生成の優劣を示すための区分ではない。百三十八種の動作という数、平均八・五語という短さも条件の一部であり、長い指示や複雑な筋への一般化は本文で確立していない。試験用三百四十四件という規模、一条件一本三名多数決という手順も値と一体で引用する\autocite{btd131}。十二種の生成器から一万千三百三十本の映像を得て三万六千五百件の注釈を集め、試験用三百四十四件では一条件につき一本を三名が多数決で判定した。二値判定の一致率は物理的適合で七五％、常識で七〇％であり、報告条件での値は同時達成で三九・六％、適合六三・三と常識五三の積み上げの難しさを示す。自動判定に調整したVideoCon-Physicsは適合でROC-AUC八二、常識で七三を示し、Geminiの七三と五八やGPT-4-Visionの五三との差が報告条件で示された\autocite{btd131}。この差は自動判定器同士の比較であり、人手多数決そのものの置換を意味しない。調整の土台は三万六千五百件の注釈であり、試験用三百四十四件の多数決判定との照合を経ている。ただし判定器の学習分布と試験分布の隔たり、生成器十二種の顔ぶれの固定、二〇二四年という時点の固定は外せない。後継の計測や模擬器に根ざす力学検証は本文の範囲外である\autocite{btd131}。常識判断はより主観的であり、後継の計測や模擬器に根ざす力学検証は本文の範囲外である。したがって長時間の整合性を一枚の鮮明さと同一視せず、天井条件付きで読む必要がある。同時達成三九・六％という値は、単発の見栄えが良くても相互作用の持続が破綻し得ることを示す。以降の各項でも、評価器の版と母集団と手順の三点組を欠かさず添え、開いた問いを閉じない姿勢を保つ。

\subsection{凍結専門家で同期を罰する構成}
口唇同期の訓練信号と評価指標の分岐を理解するため、Wav2Lipは専門家識別器を凍結して生成器を罰する構成を採った点が重要である。具体的には、色彩と残差と余弦類似度に二値交差エントロピーを組み合わせた五フレーム窓のSyncNetを専門家とし、LRS2のみで学習した生成器に対して同期損失とL1損失を重み〇・〇三、視覚GAN損失を重み〇・〇七で加えた\autocite{btd129}。重みの読み方も条件付きである。〇・〇三と〇・〇七という配分は報告条件での一例であり、配分を変えた一般化は本文で確立していない。LRS2のみという学習条件、専門家凍結という設計も値と一体であり、専門家を変えた一般化は確立していない。生成器の学習集合と評価器の学習集合が異なるという非対称は、手順の前提として残す\autocite{btd129}。報告条件ではLRS2とLRWとLRS3の無作為音声組み合わせで唇同期距離が六・五一二と六・三八六と六・六五二、信頼度が七・四九〇と七・七八九と七・八八七となり、実映像の七・〇一二前後の距離帯と近接した。ReSyncEDの吹替と無作為と合成音声の選好では、GAN付き構成が六〇・二％と六四・五％と五一・二％を獲得し、無同期に対する合算選好が九〇％に達した\autocite{btd129}。ここで注意すべきは、距離と信頼度の双方が同一の公開SyncNetという一つの評価器に由来する点である。評価器を固定したまま生成器側を変える比較は成立するが、評価器自体が異なる比較は成立しない。学習に用いたLRS2のみという条件、五フレーム窓という条件、損失重み〇・〇三と〇・〇七という条件は、数値と一体で引用する\autocite{btd129}。しかし指標に用いた公開SyncNetは学習データが異なり、SyncNet自体の不安定性を継承する。対象は正面向きの話者映像に限られ、一般の音事象の時間整合や多話者・横顔・遮蔽への拡張は本文では確立していない。したがって唇同期距離の数値を条件越えの真実として横並びにせず、撮影条件と音声条件を束ねたまま読む姿勢が求められる。同期の良さが物語構造や物理整合を保証しない点を明示して読む。この専門家固定の発想は次項の唇読み特徴への置換に継承されるが、両者の値は互換性がなく条件をまたぐ横並びは成立しない。
\end{multicols}

\begin{center}
\small
\begin{tabularx}{\textwidth}{@{}p{0.22\textwidth}p{0.36\textwidth}X@{}}
\toprule
項目 & 報告条件の値 & 読みの留保 \\
\midrule
評価記述の構成 & 六八八件（固体同士二八九・固体流体二九一・流体同士一〇八） & 候補千件からの人手選別を経た範囲 \\
動作と指示 & 百三十八種・平均八・五語 & 短い指示の範囲での検証 \\
収集規模 & 十二生成器・一万千三百三十本・三万六千五百注釈 & 二〇二四年時点の模型群 \\
判定の一致率 & 適合七五％・常識七〇％（三名多数決） & 二値化は画質と運動を畳み込む \\
報告条件の同時達成 & 三九・六％（適合六三・三・常識五三） & 単発の鮮明さと別軸の値 \\
自動判定の調整 & VideoCon-Physicsは適合ROC-AUC八二・常識七三 & 人手多数決の代理に留まる \\
\bottomrule
\end{tabularx}
\end{center}
\autocite{btd131}

\begin{multicols}{2}
\raggedcolumns
\subsection{唇読み特徴への置換と指標の非互換性}
SyncNet損失の置換という変化を担うのが、唇読み特徴に基づくAV-HuBERT専門家の導入である\autocite{btd130}。置換の意味を正確に読むため、前項との対比を整理する。前項の専門家は色彩と残差と余弦類似度に二値交差エントロピーを組み合わせた五フレーム窓のSyncNetであり、本項の専門家はTransformer最終層七六八次元特徴に基づく教師なし音視覚余弦類似度と二値交差エントロピーである。窓の広さは共通するが、特徴の出所も学習データも損失の配分も異なる。共通する五フレームという窓だけを取り出して、両者を一本の改善系列として語らない。置換は標準の交代ではなく、条件付きの差である\autocite{btd129,btd130}。手法はTransformer最終層の七六八次元特徴を用い、生成区間に対する教師なし音視覚余弦類似度と二値交差エントロピーを中心に、GAN損失と知覚損失とL1損失を重み一〇と一と〇・五で組み合わせた。顔切り出しは九六画素、五フレーム窓、メルは十六キロヘルツで十六×八〇の条件である\autocite{btd130}。報告条件ではLRS2でSSIM〇・九四七、PSNR三一・二七三、FID四・五一、landmark距離一・一八八、唇同期信頼度七・九五八と距離六・三〇一、新規のAVSu〇・五〇八とAVSm〇・九三九とAVSv〇・八七九が示された。HDTFではSSIM〇・九三三、PSNR三〇・五七九、FID一六・七六であり、十本十名の五段階利用者調査では三・九二と四・〇二と三・九五を示し、Wav2Lipの二・九一前後との差が報告条件で示された\autocite{btd130}。集合を変えた読み方も条件付きである。LRS2という集合とHDTFという集合は別の顔ぶれと撮影条件を持ち、集合をまたぐ順序付けは行わない。FID四・五一と一六・七六という値は、各集合での観察であり、集合一般の性質ではない。十本十名という規模は効果の方向を示すに留まり、効果の大きさの確定ではない\autocite{btd130}。教師なし変種が相対的に良いという切断結果も含まれる\autocite{btd130}。切断の読み方も条件付きである。重み一〇と一と〇・五という配分、九六画素の切り出し、五フレーム窓、十六キロヘルツ十六×八〇のメルという条件の組のもとでの差であり、配分や切り出しを変えた一般化は本文で確立していない。七六八次元という特徴量の出所、最終層という取り出し位置も条件の一部である\autocite{btd130}。ただし話者正面の発話映像に限定され、新指標の独立した採用は未検証で、利用者調査も小規模である。距離指標の置換が標準として定着したとは言えず、条件付きの改善として継承関係を記述する。前項の公開SyncNetによる距離と本項の唇読み特徴による距離は、切り出しも窓も学習データも異なるため互換性がなく、両者を一本の物差しとして並べることはできない。同一切り出しと同一窓の組での条件付き差として記録し、話者正面という天井を添えて継承する。

\subsection{多ターン対話の区切り手順と逓減の記録}
音声対話の多ターン評価へ視点を移すと、MTR-DuplexBenchは区切り方自体を手順化した点に特徴がある。時刻付き書き起こしと音声区間検出と中規模書き起こし模型から大規模模型による裁定を六走査で重ね、三〇％重なり多数決の中央値と併合で発話ターンを定め、助手窓には正解履歴を与えつつ次利用者を消音する。四次元は合成会話と自然対話と指示と安全であり、成功率と模型評点と拒否率に秒単位の遅延を添える。構成は十ターン合成が二百件、百二十秒自然対話が二百件で平均五・八八ターン、指示が三百件、安全が五百二十件である\autocite{btd135}。四次元の読み方も条件付きである。合成会話と自然対話と指示と安全の四者は別の手順と別の母集団を持ち、四者を一つの順序に載せない。十ターン合成二百件という規模、百二十秒自然対話二百件という規模、指示三百件と安全五百二十件という規模の四点を外して、対話評価一般の姿として語らない。成功率と模型評点と拒否率に秒単位の遅延を添えるという設計も、値の解釈の前提である\autocite{btd135}。報告条件では滑らかな会話成功が第一ターン七三・〇〇％から十ターン通算五七・四〇％へ低下し、割込は七二・五から五四・二、pauseは九三・五から八四・八、背景雑音は五三・〇から二五・七へ低下した。遅延は第一ターン約〇・六四秒から逓増し、カスケードでは九〜十二秒、指示は六八・〇から四一・九％で対照の八六・五と九二・六を下回り、安全は九〇〜九一％で対照の九九・七％を下回った\autocite{btd135}。ただし二〇二五年十一月時点の手順で採用は未検証であり、模型判定と区切りへの依存が残る。区切りの誤りは成功率と遅延の双方に波及し、判定模型の版固定がなければ再現性が揺らぐ\autocite{btd135}。六走査の重ね合わせ、三〇％重なり多数決の中央値と併合という手順の細部は、再現のための条件の一部であり、省略して数値だけを引用しない。助手窓に正解履歴を与えつつ次利用者を消音するという設計も、成功率の解釈に影響する。第一ターンの値を長時間能力に拡張しない記述が求められる。自然対話の平均五・八八ターンという規模感とともに引用する\autocite{btd135}。

\subsection{単発の主張を長時間の証拠に格上げしない}
時間軸の評価は忠実度とは別軸であるという整理が、以上から導かれる\autocite{btd131,btd129,btd130,btd135}。別軸の意味を正確に読むため、各証拠の天井を再確認する。物理は二〇二四年時点の模型群と二値判定の粗さ、同期は話者正面と評価器の版固定、対話は二〇二五年十一月時点の手順と模型判定と区切りへの依存である。三者の天井は互いに異なり、一つの天井に畳めない。天井を外した数値の横並びは、技術史の因果を崩す。VideoPhyの同時達成三九・六％という値は、単発の見栄えが良くても相互作用の持続が破綻し得ることを示し、対応秒数の表示が実証された持続性を意味しないという映像系の境界と接続する\autocite{btd131}。本項ではこの gating を手順化する。第一に、対応秒数の表示や一段階の速さや単一値の分布点は、長時間の証拠に格上げしない。第二に、分解された次元と人手勝率の照合を経ずに単一値で順位付けしない。第三に、評価器の版と母集団と区切り手順の三点を欠かさず添える。第四に、開いた問いを閉じない。gating を外すと技術史の因果が崩れる。各章の継承記述は表現と目的とサンプリングを分けた上で改善と代償を一組にして残す\autocite{btd131,btd135}。Wav2LipとAV-HuBERTの距離と信頼度の差は、同一撮影集合と同一公開SyncNetの組でのみ意味を持ち、異なる唇読み特徴や異なる顔切り出し条件をまたぐ比較は成立しない\autocite{btd129,btd130}。DuplexBenchの第一ターンと通算の落差は、ターンを重ねるほど割込と背景と指示追従が劣化する事実を示し、一問一答の成功を長時間対話の能力と読み替えてはならない\autocite{btd135}。以上を束ねると、時間軸の評価は忠実度とは別軸であるという整理が導かれるという本節の結論は、五つの小節の積み上げである\autocite{btd131,btd129,btd130,btd135}。物理の手順と天井、同期の訓練と評価の分岐、専門家の置換と非互換性、対話の区切りと逓減、そして格上げ防止の gating の五点は、いずれも条件付きの記述であり、条件を外した一般化ではない。五点を一つの順序に載せず、並置のまま残す。並置を保つことが、単発の鮮明さと持続の正しさを分ける読解の実践である。物理得点は妥当性の天井、同期得点は話者正面の天井、対話得点は区切りと判定模型の天井とともに引用し、単発の主張を長時間の証拠に格上げしない。とりわけ物理と同期と対話は評価器の版と母集団と区切り手順の三点組を欠かさず、版と条件を固定して引用する。
\end{multicols}

\begin{center}
\small
\begin{tabularx}{\textwidth}{@{}p{0.24\textwidth}p{0.34\textwidth}X@{}}
\toprule
次元（同一手順内） & 第一ターンの値 & 十ターン通算の値 \\
\midrule
滑らかな会話の成功 & 七三・〇〇％ & 五七・四〇％ \\
割込への対応 & 七二・五％ & 五四・二％ \\
pauseへの対応 & 九三・五％ & 八四・八％ \\
背景雑音への対応 & 五三・〇％ & 二五・七％ \\
指示追従 & 六八・〇％ & 四一・九％（対照八六・五／九二・六） \\
安全性 & 九〇〜九一％ & 対照九九・七％を下回る \\
遅延 & 約〇・六四秒 & 逓増（カスケード九〜十二秒） \\
\bottomrule
\end{tabularx}
\end{center}
\autocite{btd135}

\begin{claimboundary}[本節の読解上の境界]
唇同期得点は条件をまたぐ真実ではなく、物理得点は妥当性の天井付きで読む。対象は二〇二四年時点の模型群であり、二値判定は画質と運動の良否を一つの可否に畳み込む。後継の計測や模擬器に根ざす力学検証は本文の範囲外である。公開SyncNetは学習データが異なり不安定性を継承し、対象は正面向きの話者映像に限られる。一般の音事象の時間整合や多話者・横顔・遮蔽への拡張は確立していない。新指標の独立した採用は未検証で、利用者調査は十本十名の小規模に留まり、距離指標の置換が標準として定着したとは言えない。二〇二五年十一月時点の対話手順の採用は未検証であり、模型判定と区切りへの依存が残り、区間横断の採用と審査者なし安定性は確立していない。編集の保存目標と局所性の本体は別節の範囲であり、本節は時間的一貫性と同期と対話の持続に限る。
\end{claimboundary}

\section{実行と配備：サンプリング・遅延・記憶・公開性}
\label{sec:runtime}
\sectionkicker{RUNTIME}

\begin{multicols}{2}
\raggedcolumns
速さの数値は、機材と設定とサンプリングの予算を外して読むことができない。本節が扱うのは、千段階級のサンプリングを一段階級へ短縮する系譜と、その短縮を実測で裏打ちする手順である。対象は六つの証拠群である。漸進的蒸留による半減反復、潜在整合模型とそのLoRA化、敵対的蒸留、携帯級への構造再設計、そしてH100上の独立計測である\autocite{btd078,btd079,btd080,btd081,btd082,btd137}。六者の並びは手順の短縮から構造の圧縮へ、そして実測の天井へという読みの順序であり、到達の順序ではない。順序を到達に読み替えないことが、本節の構成の理由である。いずれの主張も、ハードウェアと精度と画素数と段階数と学習予算を束ねた条件付きの値であり、束ねを外した横並びの速さ比べは行わない\autocite{btd078,btd081,btd082,btd137}。束ねの五点は省略できない。機材を変えれば律速が変わり、精度を変えれば保持と速度が変わり、画素数を変えれば計算と記憶が変わり、段階数を変えれば反復費用が変わり、学習予算を変えれば蒸留の到達が変わる。五点のいずれかを欠く引用は、適用域の不明な主張になる。構造を変えずに手順だけを短縮する系統と、構造自体を圧縮する系統は別の代償を持つ。前者は一段階と二段階の劣化を残し、後者は探索費用と再学習要件を残す\autocite{btd078,btd082}。代償の所在が異なるため、両者を一つの速さ比べに載せない。手順の短縮は学習予算と誘導条件の組で読み、構造の圧縮は探索費用と再学習要件の組で読む。独立計測は提供元主張の天井を押さえる役割を担い、単一実験室・非査読・H100専用という出自とともに引用する。それが本節の読みの作法である\autocite{btd078,btd137}。

\subsection{半減反復で段階数を畳む漸進的蒸留}
千段階級のサンプリングを半減の反復で短縮する起点として、漸進的蒸留は二段のDDIM教師更新を一段の生徒更新に写像する手順を採った。安定化のため画像と雑音と速度の各パラメータ化に打切りSNR重みと余弦スケジュールを組み合わせ、構造は教師と同一のまま一段階ずつNをN／二へ圧縮する\autocite{btd078}。報告条件ではCIFAR-10で四段階FID三・〇を示し、DDIMの百段階四・一六や十段階一三・三六との差が同一学習目的のままサンプリング側の変更として示された。八段階二・五七、二段階四・五一、一段階九・一二という階段状の劣化も同時に記録され、一段階と二段階の欠落が明確である。ImageNet-64とLSUN-128では四〜八段階まで劣化が小さく保たれる一方、それ以下の急落がDDIM側で明確になる\autocite{btd078}。集合を変えた読み方も条件付きである。CIFAR-10という集合、ImageNet-64という集合、LSUN-128という集合の三者は別の難所を持ち、集合をまたぐ順序付けは行わない。四〜八段階まで劣化が小さいという振舞いは、各集合での観察であり、集合一般の性質ではない。段階数と集合と目的関数を固定した差としてのみ引用し、機材をまたぐ速さ比べに拡張しない。一段階あたり五万更新、一段階と二段階向けに十万更新という予算が示され、総量は元の学習以下に収まる\autocite{btd078}。予算の読み方も条件付きである。更新回数の総量が元の学習以下に収まることは、蒸留が追加学習として軽いことを意味するが、サンプリングの費用が消えたことを意味しない。配備後のサンプリング費用と学習時の蒸留費用は別の欄であり、混ぜて一つの軽さとして語らない。構造が教師と同一のままという条件は、重みの保持費用が変わらないことも意味する。軽くなった欄はサンプリングの反復回数であり、変わっていない欄は構造の規模である\autocite{btd078}。ただし一段階と二段階の劣化、構造圧縮の欠如、蒸留後の確率的サンプリングが蒸留DDIMと非蒸留の中間に位置する点が残る。遅延やVRAMや量子化の束ねは本文になく、条件付きの縮約として継承される。この半減反復の発想は後段の整合蒸留や敵対蒸留に継承されるが、構造圧縮を伴わない点が携帯再設計との分岐になる。段階数と集合と目的関数を固定した差としてのみ引用し、機材をまたぐ速さ比べに拡張しない。

\subsection{潜在整合蒸留と低負担のLoRA化}
潜在空間での整合性蒸留とそのLoRA化は、同一の流れを異なる実装負担で配備する変化である。潜在整合模型は潜在整合関数と雑音予測パラメータ化を用い、誘導付き一段階蒸留で拡張確率フロー常微分方程式を解き、DDIMとDPM-Solverの刻みを間引き係数二〇で進め、一段階から四段階の推論を可能にする。学習はLAION-Aestheticsの千二百万件五百十二画素と六十五万件七百六十八画素で四千段階、約三十二A100時間という予算が示された\autocite{btd079}。予算の読み方も条件付きである。千二百万件と六十五万件という資料規模、五百十二画素と七百六十八画素という画素数、四千段階という学習量の三点を外して、整合蒸留一般の費用として語らない。約三十二A100時間という値は当該資料と当該画素数の組での観察であり、資料や画素数を変えた一般化は本文で確立していない。間引き係数二〇という刻みの進め方、一段階から四段階という推論範囲も条件の一部である\autocite{btd079}。報告条件では五百十二画素で誘導八の四段階FID一一・一〇を示し、DDIM二二・三八やDPM++一八・四三や誘導蒸留一五・一二との差が同一誘導条件で示され、CLIPは二八・六九で二五・八九〜二七・二五を上回った。一段階では三五・三六で一八三前後の対照との差は残るが隔たりも残る\autocite{btd079}。四段階実用と一段階研究の線引きは、この隔たりの記録である。誘導八という条件、五百十二画素という条件、LAION-Aestheticsという資料の条件を外して、整合蒸留一般の速さや良さとして語らない。約三十二A100時間という予算も条件の一部であり、異なる資料や画素数への転用は本文で確立していない\autocite{btd079}。LCM-LoRAはSD-V1.5で六七・五M、1.3Bで一〇五M、SDXLで一九七Mのみを学習し、加速ベクトルと様式ベクトルを〇・八と一・〇等で無学習結合して四段階誘導を差し込む。ただし報告は定性的でFID表はLCM側に委ね、結合重みは発見的で模型毎の一般化は未検証であり、遅延とVRAMの束ねも本文にない\autocite{btd080}。一段階の隔たりと誘導や刻みへの鋭敏さは、四段階実用と一段階研究の線引きを示す。結合重みの発見的手順や模型毎の再調整の要否は未検証であり、加速と様式の無学習結合を一般則と呼ばない\autocite{btd080}。〇・八と一・〇という値は報告条件での一例であり、他の模型や様式への転用は本文で確立していない。学習対象が六七・五Mから一九七Mという小規模部品に限られる点も条件の一部であり、本体全体の再学習の軽さを意味しない。定性的報告と定量的FID表の分担、すなわち質の記述はLoRA側、数の記述はLCM側という分担を崩して引用しない\autocite{btd080}。潜在整合の短縮を携帯や敵対の短縮と同一視せず、学習予算と誘導条件を束ねて記録する\autocite{btd079,btd080}。
\end{multicols}

\begin{center}
\small
\begin{tabularx}{\textwidth}{@{}p{0.30\textwidth}p{0.30\textwidth}X@{}}
\toprule
サンプリングの条件（CIFAR-10・同一目的関数） & FID & 読みの留保 \\
\midrule
漸進蒸留・八段階 & 二・五七 & 四〜八段階は劣化の小さい域 \\
漸進蒸留・四段階 & 三・〇 & DDIM百段階四・一六との差は手順側の変更 \\
漸進蒸留・二段階 & 四・五一 & ここから劣化が明確になる \\
漸進蒸留・一段階 & 九・一二 & 欠落が残る段階 \\
DDIM・十段階 & 一三・三六 & 短縮なしの対照 \\
\bottomrule
\end{tabularx}
\end{center}
\autocite{btd078}

\begin{multicols}{2}
\raggedcolumns
\subsection{写実性を別系統の損失で補う敵対的蒸留}
敵対的蒸留は一段階級の写実性を別系統の損失で補う変化であり、SD2.1とSDXLの生徒を四時刻と零終端SNRで学ぶ。凍結DINOv2のViT-S頭に文章と画像の条件付けとR1正則化を組み合わせたhinge敵対損失に、画素空間の最終NFSD変種である得点蒸留損失を加え、推論時は無誘導で反復精緻化する\autocite{btd081}。損失の読み方も条件付きである。凍結という条件、ViT-Sという規模、R1正則化という安定化、四時刻と零終端SNRという学習条件の四点を外して、敵対蒸留一般の姿として語らない。画素空間の最終NFSD変種という得点蒸留の指定、無誘導で反復精緻化という推論条件も値と一体であり、条件を変えた一般化は本文で確立していない。SD2.1とSDXLという生徒の指定も条件の一部である\autocite{btd081}。報告条件ではCOCO五千件五百十二画素でADD-M一段階FID一九・七とCLIP〇・三二六を示し、UFOGenの二二・五と〇・三一一やInstaFlow-0.9Bの二三・四と〇・三〇四や漸進蒸留四段階の二六・四と〇・三〇〇との差が同一画素数と同一集合で示された。人手Eloでは一段階がLCM-XL四段階を上回り、四段階がSDXL-Base五十段階を上回ると報告された。時間はA100混合精度五百十二画素で一段階〇・〇九秒、DPM二十五段階〇・八八秒という機材と精度と画素数を束ねた値である\autocite{btd081}。ただし教師より多様性がやや低く、写実と過平滑の交換があり、評価は五百十二画素に限られる。記録locatorの転記欠陥は正規二三一一・一七〇四二側の本文で消費されており、その旨を添えて引用する。一段階の人手選好と多様性の低下は同一切断の裏表であり、写実の向上を多様性保持と読み替えない\autocite{btd081}。人手Eloの読み方も条件付きである。一段階がLCM-XL四段階を上回り、四段階がSDXL-Base五十段階を上回るという報告は、人手選好という一つの軸での差であり、多様性や整合や細部忠実の軸での差ではない。軸を変えれば順序は変わり得る。無誘導で反復精緻化するという推論条件、四時刻と零終端SNRという学習条件も数値と一体で引用する\autocite{btd081}。評価画素数と集合規模と機材精度の三点を外したFID比べは行わず、同一五百十二画素同一集合の差として継承する。

\subsection{構造自体を圧縮する携帯級の再設計}
携帯級への再設計は、構造自体を圧縮して一段階を実用域に寄せる変化である。UViTは隘路Transformerを千二十四溝に集約し、六十四画素と三十二画素および外側十六画素の一部自己注意を落としつつ七十七トークンへの交差注意を残し、共有鍵値で五％、GELUからSwishへの置換、ReLU注意の一万微調整、前向き層の四対三化で一〇％、分離畳み込み、十一層化で三八六Mと百八十二GFLOPsに圧縮した。復号側はf8-c8のVAEに蒸留復号器を添える。学習は一億五千万対でUFOGen全微調整と拡散損失の併用が相対的に良いとされた\autocite{btd082}。報告条件ではCOCO三万件一段階FID一一・六七とCLIP〇・三二〇を示し、UFOGenの一二・七八と〇・三一七やSnapFusion八段階の一三・五と〇・三〇八やSD五十段階の九・六二と〇・三〇四との差が同一集合で示され、五十段階では八・六五と〇・三二五、人手選好は一段階二六・六五で五十段階二七・三〇に迫った。iPhone-15-Pro実測は一段階が四＋九二＋一四二＝二三八ミリ秒で、長期スケジュールの六四六と一七四〇と七四二九との差が機種固定で示され、〇・二秒五百十二画素生成が可能とされた\autocite{btd082}。実測の読み方も条件付きである。四と九二と一四二という内訳、六四六と一七四〇と七四二九という対照の二点を外して、携帯一般の速さとして語らない。iPhone-15-Proという機種、五百十二画素という画素数、一段階という段階数の三点を固定した値であり、三点のいずれかを変えた予測は本文で確立していない。人手選好一段階二六・六五と五十段階二七・三〇という近接も、同一切り分けでの観察である\autocite{btd082}。ただし探索に五百十二TPU十五日を要し、ReLU注意は微調整必須で再学習なき転移に隔たりが残り、VRAM頂点と量子化遅延は本文で確立していない\autocite{btd082}。探索費用の読み方も条件付きである。五百十二TPU十五日という費用は一度きりの投資であり、配備後のサンプリング費用ではない。投資と継続費用を混ぜて一つの重さとして語らない。再学習なき転移に隔たりが残るという点は、圧縮構造が特定の資料と目的関数に結びついていることを意味し、圧縮一般の軽さとして語らない。三八六Mと百八十二GFLOPsという値は十一層化と分離畳み込みと注意削減の組の結果であり、層構成を変えれば変わる\autocite{btd082}。構造圧縮と蒸留の組合せが一段階実用を支える一方、探索費用と再学習要件は配備の代償として残る。機種固定の実測を他機種の予測に拡張せず、画素数とスケジュールと層構成を束ねた値として引用する。

\subsection{機材の天井を押さえる独立計測}
独立計測で機材の天井を押さえるのが、H100上のFLUX.2-klein-4Bの計測である。手順はDiffusersのFluxPipelineをbf16とflow照合スケジュールでH100 SXM八十GBに固定し、五解像度と三段階数でCLIP整列とLPIPSと属性束縛とバッチ一〜十六と並列一と二と四と八と費用を測った\autocite{btd137}。報告条件では五百十二画素四段階〇・一九秒で毎秒五・一四枚、七百六十八画素〇・三三秒、千二百八十×七百二十で〇・五〇秒、千二十四画素四段階〇・五七秒で毎秒一・七七枚、八段階一・〇〇秒、十二段階一・四三秒が示された。VRAMは五百十二画素十六GBから千二十四画素十七・八GB、バッチ八頂点で毎秒一・九一枚、並列は一で〇・四三、二で〇・八二の九七％、四で一・四八の八七％、八で二・五三の七四％、CLIP平均〇・三三五で段階数にほぼ平坦、二段階〇・三七五から十二段階〇・三六七、単価は提示雲価格で約〇・〇〇〇四ドル、SDXL比五〜十倍とされた\autocite{btd137}。費用の読み方も条件付きである。提示雲価格という出所、約〇・〇〇〇四ドルという単価、SDXL比五〜十倍という比率の三点を外して、一段階級一般の費用として語らない。比率の相手であるSDXL側の条件、雲価格の時点も値と一体であり、時点や相手を変えた一般化は本文で確立していない。費用の主張は提供元の提示に基づくものであり、独立検証のない数値は事実として引用しない\autocite{btd137}。ただし単一実験室の非査読でH100専用であり、消費者GPU再現は未了である\autocite{btd137}。出自の限定は数値と一体で引用する。H100 SXM八十GBという機材、bf16という精度、DiffusersのFluxPipelineという管線、flow照合スケジュールという設定の四点を外して、一段階級一般の速さとして語らない。並列一で〇・四三、二で〇・八二の九七％、四で一・四八の八七％、八で二・五三の七四％という逓減は、当該機材と当該管線の組での振舞いであり、他機材への予測ではない\autocite{btd137}。解像度と段階数と精度と機材を外した速さ比べは行わず、束ね条件ごとの読み替えに留める。段階数を増やしても整列がほぼ平坦という振舞いは、速度と品質の交換が解像度側に寄る事実を示す。平坦の読み方も条件付きである。二段階〇・三七五から十二段階〇・三六七という平坦は、CLIP整列という一つの軸での振舞いであり、LPIPSや属性束縛の軸での振舞いではない。軸を変えれば段階数の効き方は変わり得る\autocite{btd137}。携帯実測や蒸留報告との比較は機材も精度も異なるため行わず、束ね条件ごとの天井として並置する。以上を束ねると、短縮の系譜は手順の短縮と構造の圧縮と実測の天井の三層で読むべきことが分かる\autocite{btd078,btd079,btd080,btd081,btd082,btd137}。一段階の値を多段階の値と直接横並びにせず、予算と条件の組で読み替える。手順の短縮は反復回数を削り、構造の圧縮は保持と計算を削り、実測の天井は主張の適用範囲を区切る。三層を混ぜて一つの速さとして語らないことが、本節の結びの作法である。結びの三層を再確認する。手順の短縮は漸進と整合と敵対の三系統の並置であり、三系統の順序付けではない。構造の圧縮は携帯級の一例であり、圧縮一般の確定ではない。実測の天井はH100とbf16と管線固定の一例であり、機材一般の確定ではない。三点を並置のまま残し、一段階の到達を多段階の到達に読み替えない\autocite{btd078,btd079,btd080,btd081,btd082,btd137}。
\end{multicols}

\begin{center}
\small
\begin{tabularx}{\textwidth}{@{}p{0.26\textwidth}p{0.22\textwidth}p{0.20\textwidth}X@{}}
\toprule
条件（H100・bf16・管線固定） & 遅延 & 処理能力 & 備考 \\
\midrule
五百十二画素・四段階 & 〇・一九秒 & 毎秒五・一四枚 & VRAM十六GB \\
七百六十八画素・四段階 & 〇・三三秒 & ― & 同一管線の実測 \\
千二十四画素・四段階 & 〇・五七秒 & 毎秒一・七七枚 & VRAM十七・八GB \\
千二十四画素・八段階 & 一・〇〇秒 & ― & 段階数逓増 \\
千二十四画素・十二段階 & 一・四三秒 & ― & 整列はほぼ平坦 \\
千二百八十×七百二十・四段階 & 〇・五〇秒 & ― & 同一管線の実測 \\
\bottomrule
\end{tabularx}
\end{center}
\autocite{btd137}

\begin{claimboundary}[本節の読解上の境界]
実行時間の数値は機材と設定とサンプリング予算に縛られ、単独実験室の限定を明示する。一〜二段階の劣化、構造圧縮の欠如、蒸留後の確率的サンプリングの中間性は残る。遅延・RTF・VRAM・機材・量子化・オフロードの束ねは漸進蒸留の本文になく、一段階の隔たりと解法・間引き・誘導への鋭敏さは残り、利用者集合向けの調整なし汎用部品はない。結合重みは発見的で模型毎の一般化は未検証であり、明示的なVRAM頂点とint8量子化遅延は携帯再設計の本文で確立していない。計測は単一提供元実験室の非査読でH100専用であり、消費者GPU再現と提供元主張の同等性は当該構成を越えて確立していない。記録locator二三一一・一七〇四九は転記欠陥であり、正規二三一一・一七〇四二側の本文で消費した。
\end{claimboundary}

\section{評価の方法論：指標は交換可能ではない}
\label{sec:evaluation}
\sectionkicker{EVALUATION}

\begin{multicols}{2}
\raggedcolumns
分布の近さと構成の正しさと人手の選好は、互いに交換できない物差しである。本節が扱うのは、評価指標の使い分けとその妥当性の境界である。対象は十の証拠群である。分布類似のFIDとIS、構成診断のGenEvalとT2I-CompBench、人手選好のPick-a-Pic、音響距離のFAD、読書音声のLibriSpeech、話者照合のGE2E、そして映像多次元のVBenchとVBench-2.0である\autocite{btd083,btd084,btd085,btd086,btd087,btd088,btd089,btd090,btd092,btd093}。本節では条件の異なる数値を横断比較して順位表化しない\autocite{btd085,btd086,btd092}。順位表化の禁止は本節の編集方針の一部である。分布の近さと構成の正しさと人手の選好は別の軸であり、軸を変えれば順序は変わり得る。軸を固定しない比較は、比較の体をなさない。提供元による評価の数値は独立再現まで隔離し、事実として引用しない。本文の消費水準も明示する。FIDに関する主張は要旨段階の消費に留まることを明示し、LibriSpeechに関する蔵書事実は断片段階の確認に留まることを明示する。指標の値を構成の良さに読み替えず、標本数と埋め込みと前処理の三点を添えて記録する。それが本節の読みの作法である\autocite{btd083,btd089}。

\subsection{分布類似の起点を本文境界付きで押さえる}
分布類似の起点を本文境界付きで押さえるため、FIDは要旨段階、ISは本文段階として区別する。FIDはInception符号化層の埋め込みに多変量正規を当てた生成集合と大規模実集合のFréchet距離であり、二時刻更新則とともにCelebAとCIFAR-10とSVHNとLSUN寝室で試されたと要旨は述べる。FAD側の第二節の記述でも同趣旨が補強される。ただし消費本文は要旨に限られる。ここで明示する。本節のFIDに関する記述は要旨段階の消費に基づくものであり、正確な埋め込み層と標本数と前処理は未検証で、分布水準のプロンプト盲検という性格が残り、FID値や必要標本数や埋め込み依存は確立していない\autocite{btd083}。要旨段階の消費とは、arXiv要旨の確認とFAD側第二節の記述による補強の範囲を指し、全文HTMLは404でar5iv変換は失敗し、PDFとTeXのみで本文消費に至っていない経緯を含む。要旨に書かれた試行集合の名指しは、試行の存在を示すが、試行の条件を示さない。埋め込み層の指定、前処理の手順、標本数の明示がなければ、値の再現はできない\autocite{btd083}。要旨段階の主張を本文で裏付けられた言語では記述しない。ISはImageNet学習済みInceptionによるexp(E[KL])であり、条件付きの低エントロピーが物体らしさ、周辺の高エントロピーが多様性を意味し、約五万標本を要すると本文は求める。報告条件ではCIFAR-10五万標本で実一一・二四±〇・一二、自手法八・〇九±〇・〇七、切断で三・八七±〇・〇三、人手真偽七八・七％で上位一％選別時七一・四％が示された\autocite{btd084}。ただし独立学習模型の目安に限られ、直接最適化は敵対例を生み、ImageNet物体バイアスが残る\autocite{btd084}。目安の読み方も条件付きである。約五万標本という要件、CIFAR-10という集合、ImageNet学習済みという埋め込みの三点を外して、IS一般の良さとして語らない。記録locatorの転記欠陥は正規一六〇六・〇三四九八側本文で消費した旨を添え、順位表化しない。転記欠陥の注記は出典の取り違えを防ぐための必須要素であり、条件束ねの一部として残す\autocite{btd084}。要旨段階と本文段階と断片段階を混ぜない引用が求められる。分布指標の値を構成の良さに読み替えず、標本数と埋め込みと前処理の三点を添えて記録する\autocite{btd083,btd084}。

\subsection{全体点から分離診断への移動}
全体点では捉えられない構成失敗を診断するため、GenEvalとT2I-CompBenchは検出器と問答を分離する方向に進んだ。GenEvalは六課題×五百五十三定型プロンプトでCOCO八十物体とBerlin-Kay十一色を扱い、Mask2Formerを信頼度〇・三、計数で〇・九に定め、存在と個数と重心位置を測り、色彩は切り出しへのCLIP ViT-L/14で測り、二値正誤を課題平均する\autocite{btd085}。報告条件では人手一致八三％、注釈者間八八％、全員一致部分九一％でCLIPScoreの八〇％と八七％との差が同一手順で示され、IF-XL全体〇・六一とSD-v2.1〇・五〇、位置一五％、束縛三五％の壁が示された\autocite{btd085}。一致率の読み方も条件付きである。人手一致八三％という値は同一手順での観察であり、手順を変えた一般化ではない。注釈者間八八％という値、全員一致部分九一％という値は、判定の安定性の範囲を示すが、判定の正しさを示さない。位置一五％と束縛三五％という壁は、検出器と問答の天井を映すものであり、模型の到達の確定ではない\autocite{btd085}。ただしCOCO語彙に閉じ、写真学習検出器は孔と重なりと図案に弱く、美観と写実は捨象される\autocite{btd085}。語彙の閉じ方の読み方も条件付きである。COCO八十物体とBerlin-Kay十一色という範囲は、定型プロンプトの設計と一体であり、範囲外の物体や色彩への一般化は本文で確立していない。信頼度〇・三と計数〇・九という閾値、切出しへのCLIP ViT-L/14という色彩手順も条件の一部であり、閾値や手順を変えた一般化は確立していない\autocite{btd085}。T2I-CompBenchは八千プロンプトを色彩と形状と質感の束縛、二次元と三次元と非空間の関係、一〜八の計数、複合に分け、分離BLIP問答の肯定確率積、空間と計数はUniDetに深さとIoU閾値、複合はCLIPと問答とUniDetの三合一とMLLM〇〜百の思考連鎖で測る。ただし第五節と第六節の模型表は未消費で順位は確立せず、検出器と問答の上限と幻覚、物理不可能なプロンプトを意図的に含む点が残る\autocite{btd086}。未消費の扱いも条件付きである。模型毎の得点順序、MLLM対専用指標の勝者は確立せず、推測して補わない。物理不可能な組成を意図的に含むという設計は、診断の難所を示すためのものであり、生成の失敗を示すためのものではない。検出器と問答の上限、MLLMの不安定さと幻覚も手順の天井として残す\autocite{btd086}。位置と束縛の壁は検出器と問答の天井を映し、模型の優劣ではなく課題の難所を示す。第五節と第六節の未消費を理由に順位を推測せず、本文の手順水準で引用する。全体点から分離診断への移動を評価の前進と位置づけ、検出器依存という代償と一組で残す\autocite{btd085,btd086}。

\subsection{人手選好を公開資産化する手順と母集団の天井}
人手選好を公開資産化する変化として、Pick-a-Picは二択と引分を沈黙の選好として集めた。背骨にSD2.1とDreamlikeとSDXLを置き誘導尺度を振り、五十八万三千七百四十七学習例と三万七千五百二十三プロンプトと四千三百七十五利用者をプロンプト非重複分割で集め、PickScoreはCLIP-Hを四千段階微調整し、InstructGPT式KL選好目的と逆プロンプト頻度重みで学ぶ\autocite{btd087}。学習手順の読み方も条件付きである。四千段階という学習量、CLIP-Hという土台、KL選好目的という目的関数、逆プロンプト頻度重みという重み付けの四点を外して、選好模型一般の姿として語らない。背骨三種と誘導尺度の振りという収集条件も値と一体であり、収集条件を変えた一般化は確立していない。選好命中七〇・五％という値は同一分割での観察であり、分割を変えた一般化ではない\autocite{btd087}。報告条件では選好命中七〇・五％を示し、専門家六八・〇％やCLIP-H六〇・八％や無作為五六・八％との差が同一分割で示され、MS-COCO勝率相関は〇・九一七でFIDの−〇・九〇〇との対比、Elo相関は〇・七九〇±〇・〇五四でHPSの〇・六七〇との対比、順位勝率は七一・三〜八五・一％が示された\autocite{btd087}。ただし自己選択の交流利用者であり、NSFW残渣が残り、超人という表現は文脈盲検の注釈者に対する比較であって依頼者本人に対する比較ではない\autocite{btd087}。母集団の読み方も条件付きである。五十八万三千七百四十七学習例と三万七千五百二十三プロンプトと四千三百七十五利用者という規模は、収集の厚みを示すが、母集団の偏りを解消しない。プロンプト非重複分割という設計、逆プロンプト頻度重みという学習手順も条件の一部であり、分割や重みを変えた一般化は確立していない\autocite{btd087}。背骨と誘導尺度を外した一般化は確立しない。公開選好の資産化は分布指標の限界を補うが、母集団の偏りという新たな天井を伴う。依頼者本人の満足と注釈者の勝率を混同せず、背骨と誘導と分割を固定した差として引用する。分布と構成と選好の三者は互いに代替せず、条件と母集団を束ねた並置に留める。
\end{multicols}

\begin{center}
\small
\begin{tabularx}{\textwidth}{@{}p{0.20\textwidth}p{0.38\textwidth}X@{}}
\toprule
診断枠組み & 手順の要点 & 境界 \\
\midrule
GenEval & 六課題×五百五十三定型・検出器と切出CLIPで二値平均 & COCO八十語彙に閉じる・美観と写実は捨象 \\
T2I-CompBench & 八千プロンプト・問答とUniDetと三合一とMLLM採点 & 第五・六節の模型表は未消費・順位は確立しない \\
Pick-a-Pic & 二択と引分の大規模収集・選好目的で微調整 & 自己選択利用者・NSFW残渣・依頼者比較ではない \\
FAD & VGGish百二十八次元・一秒窓の正規間距離 & 位相盲検・長期構造の欠落・音楽強調検証のみ \\
\bottomrule
\end{tabularx}
\end{center}
\autocite{btd085,btd086,btd087,btd088}

\begin{multicols}{2}
\raggedcolumns
\subsection{耳と話し言葉の土台を切り分ける}
耳と話し言葉の土台は、知覚距離と公開読書音声と話者照合損失を切り分けて整えられた。FADはVGGish百二十八次元を一秒窓〇・五秒刻みで切り、評価集合と清浄背景集合の正規間Fréchet距離とし、Magnatagatune五百四十時間背景と六十時間評価で十三種歪みの強度走査で検証された\autocite{btd088}。報告条件では人手相関〇・五二を示し、SDR〇・三九や余弦−〇・一五や振幅L2−〇・〇一との差が六万九千三百対二十名三百本五秒の手順で示され、清浄〇・二、約三百本二十五分から利用可とされた。ただし対数メル入力は位相盲検で一秒窓は長期構造を落とし、高域と低域通過で外れが残り、音楽強調検証のみで音声合成への転移は未確立である\autocite{btd088}。窓長の読み方も条件付きである。一秒窓〇・五秒刻みという設計、百二十八次元という特徴量、十三種歪みの強度走査という検証手順の三点を外して、FAD一般の良さとして語らない。約三百本二十五分から利用可という実用条件も、清浄背景集合との対比という前提付きである\autocite{btd088}。LibriSpeechは約千時間十六キロヘルツのLibriVox読書英語を分割整列し、Kaldi手順と文章と言語模型をopenSLR-12で公開した事実が確認される。ただしここで明示する。本節のLibriSpeechに関する記述は断片段階の確認に基づくものであり、分割と整列と基準誤り率は本文未確立で読書音声に限られる。蔵書の事実を基準性能に格上げしない\autocite{btd089}。断片段階の確認とは、openSLRとICASSP引用断片による蔵書事実の確認を指し、ICASSPのPDFはバイナリで解析不能のため本文消費に至っていない経緯を含む。記録locator一五〇七・〇八二一一は無関係の数学論文を指す転記欠陥であり、正規IEEE権威側で消費した旨を添える。約千時間十六キロヘルツという規模の表示は、蔵書の目録的事実であり、品質の証拠ではない\autocite{btd089}。GE2Eは六十四話者×十発話でL2正規d-vectorをLSTMと線形で四十次元対数メルの二十五ミリ秒十ミリ秒から作り、問合せを除く重心との尺度付余弦でsoftmaxないし難例対比と多読重みで学び、百六十フレーム五〇％重なりで推論する。等誤り率は条件依存で三・五五から三・一〇、複読で二・六七から二・三八、条件独立で四・〇六前後から三・五五と報告条件で示されたが、学習集合は社内で類似度は自然さではなく聴取試験もない\autocite{btd090}。類似度の読み方も条件付きである。六十四話者×十発話という学習条件、百五十Mと一・二M発話という社内集合の規模、百六十フレーム五〇％重なりという推論条件の三点を外して、話者類似一般の良さとして語らない。尺度付余弦という指標、問合せを除く重心という照合手順も条件の一部である。開放データ再現と話者音声合成類似度としての妥当性は確立していない\autocite{btd090}。この三者は互いに代替せず、音楽強調の距離を読書音声の品質に転用せず、社内類似度を自然さに読み替えない。埋め込み依存と母集団と窓長の三点を添えた引用が求められ、条件固定の差としてのみ比較する\autocite{btd088,btd089,btd090}。

\subsection{映像の多次元化と単一値の拒否}
映像の多次元化は、単一値への圧縮を拒む設計として継承された。VBenchは十六の分離次元に各約百プロンプトと八内容分類を置き、画質はDINOとCLIPの跨フレーム整合、運動適応ちらつき、補間事前による滑らかさ、RAFT動態、LAION美観、MUSIQ撮像で測り、条件整合はGRiTの物体と複数と色彩と規則空間、UMT動作、Tag2Text場面、CLIPとViCLIP様式で測り、次元毎に人手勝率照合を添える\autocite{btd092}。報告条件では四模型で主体整合八六〜九二％に対し空間一八〜三七％、動態四二〜九〇％という偏りが示され、経験最小最大とWebVid平均が添えられた。ただし背骨偏りを継承し、静止映像が整合次元を欺き、文章映像は文章画像、とりわけSDXL級の複数物体と空間構成に後れ、次元間重みは設計上ない\autocite{btd092}。偏りの読み方も条件付きである。四模型という対象、約百プロンプトという規模、八内容分類という範囲の三点を外して、映像評価一般の順序として語らない。経験最小最大とWebVid平均という基線は、値の位置を示すためのものであり、順位を示すためのものではない。主体整合八六〜九二％に対し空間一八〜三七％という偏りは、難所の所在を示す\autocite{btd092}。VBench-2.0は五群十八次元各約七十プロンプトで、複雑意味はLLaVA-Video-7B要約とQwen2.5-7B判定、多問問答、YOLO-World片のViT異常、ArcFace同一、CoTracker撮影法、SIFTとFLANNとRANSACとRAFT幾何、LoRA調整Qwen2.5-VL-3B個体保持、Sora以外均一プロンプト精製で測る\autocite{btd093}。計測手段の読み方も条件付きである。要約と判定に用いる模型の版、幾何検証の手順、個体保持の調整方法の三点を外して、多次元評価一般の姿として語らない。版スナップショットと手順固定を欠く再利用は行わず、最小最大と平均の基線とともに引用する\autocite{btd093}。複雑筋一〇〜一一％、動空間約二〇％、運動順序一五〜二九％の低さと解剖六〇〜八九％、二百八十四時間九五％バーの品質管理が報告条件で示されたが、幻覚は緩和のみで版スナップショットが要り、長期安定と分級要約誤りは本文指摘のまま残る\autocite{btd093}。低さの読み方も条件付きである。五群十八次元各約七十プロンプトという設計、Sora以外均一プロンプト精製という条件、LLaVA-Video-7B要約とQwen2.5-7B判定という模型依存の三点を外して、映像生成一般の到達として語らない。前処理と冗長設問による緩和、版スナップショットの要求も手順の天井として残す\autocite{btd093}。低さを示す値は失敗の宣告ではなく、物語や順序や空間の難所の所在を示す。版スナップショットと手順固定を欠く再利用は行わず、最小最大と平均の基線とともに引用する。単一値への圧縮を拒む設計自体を評価の成果と位置づけ、重み付けの不在を欠落ではなく設計として残す\autocite{btd092,btd093}。

以上を束ねると、評価の方法論は物差しの使い分けの史である。分布類似は生成集合と実集合の距離を測るが、構成の正しさを保証しない\autocite{btd083,btd084}。分離診断は課題の難所を指し示すが、検出器と問答の天井を映す\autocite{btd085,btd086}。人手選好は分布指標の限界を補うが、母集団の偏りを伴う\autocite{btd087}。音響距離と読書音声と話者照合は互いに代替しない\autocite{btd088,btd089,btd090}。映像多次元は単一値への圧縮を拒む\autocite{btd092,btd093}。条件と母集団と版を束ね、順位表化せずに並置する。それが指標は交換可能ではないという表題の意味である。要旨段階と本文段階と断片段階を混ぜず、提供元主張を事実に格上げせず、未消費の表を推測で補わない。三つの禁を守ることが、本節の結びの作法である。結びの三禁を再確認する。要旨段階の主張を本文の言語で語らないこと、提供元主張を事実に格上げしないこと、未消費の表を推測で補わないことである。三禁のいずれかを破れば、指標の値が一人歩きし、条件の異なる数値の横断比較が順位表として流通する。流通した順位表は技術史の因果を歪める。禁を守り、条件と母集団と版を束ねた並置に留める\autocite{btd083,btd085,btd086,btd089}。
\end{multicols}

\begin{center}
\small
\begin{tabularx}{\textwidth}{@{}p{0.20\textwidth}p{0.38\textwidth}X@{}}
\toprule
映像評価枠組み & 設計の要点 & 境界 \\
\midrule
VBench & 十六分離次元・各約百プロンプト・次元毎に人手勝率照合 & 背骨偏りを継承・静止映像が整合次元を欺く・次元間重みなし \\
VBench-2.0 & 五群十八次元・各約七十プロンプト・版固定の多手法計測 & 幻覚は緩和のみ・版スナップショット必須・長期安定は残る \\
\bottomrule
\end{tabularx}
\end{center}
\autocite{btd092,btd093}

\begin{claimboundary}[本節の読解上の境界]
条件をまたぐ順位表化は行わず、提供元評価は独立再現まで隔離する。FIDに関する記述は要旨段階の消費に留まり、正確な埋め込み層と標本数と前処理は未検証で、FID値と必要標本数と埋め込み依存は確立していない。ISは独立学習模型の目安に限られ、直接最適化は敵対例を生み、ImageNet物体バイアスが残る。GenEvalはCOCO語彙に閉じ、写真学習検出器は孔と重なりと図案に弱く、美観と写実を捨象する。T2I-CompBenchの第五節と第六節の模型表は未消費で、模型毎の順位とMLLM対専用指標の勝者は確立しない。Pick-a-Picは自己選択利用者でNSFW残渣が残り、超人とは文脈盲検注釈者に対する比較である。FADは位相盲検で一秒窓は長期構造を落とし、音声合成への転移は未確立である。LibriSpeechに関する記述は断片段階の確認に留まり、分割と整列と基準誤り率は本文未確立で読書音声に限られる。GE2Eは社内集合で類似度は自然さではなく、開放データ再現と話者音声合成類似度としての妥当性は確立していない。映像評価は背骨偏りを継承し、次元間重みは設計上なく、幻覚は緩和のみで版スナップショットを要する。
\end{claimboundary}

\section{収束問題：統合モデルか連携する専門家群か}
\label{sec:convergence}
\sectionkicker{CONVERGENCE}

\begin{multicols}{2}
\raggedcolumns
生成の各領域で育った模型群は、一つの統合模型に収束するのか、それとも連携する専門家群のまま残るのか。本節はこの問いに結論を出さない\autocite{btd094,btd097}。問いの立て方自体が条件付きである。統合の側には多課題の被覆と往復の畳み込みと資料側の離散化と拡散の合成という四つの立場があり、連携の側には様式別の保持と凍結模型の再利用という実践がある。四対一の構図ではなく、四つの条件付き一例と頁限定の二主張の並置である。扱うのは六つの証拠群である。多課題統合のUnified-IO、発話音声と文章の往復を畳むAudioPaLM、資料側離散化でany-to-any会話を狙うAnyGPT、拡散を合成するCoDi、そして来歴機構のC2PAとSynthIDである\autocite{btd094,btd095,btd096,btd097,btd098,btd099}。いずれも統合への一歩を示すが、その射程は条件付きであり、代償と未確立を伴う\autocite{btd094,btd095,btd096,btd097}。一歩の意味を正確に読むため、各立場の代償を再確認する。多課題統合はVQ-GAN律速と軌道限定、発話音声と文章の往復は合成目標の多さと短尺劣化、資料側離散化は高損失と天井と専用評価の不在、合成拡散は見出し欠落の継承と時間修正の部分性である。四者の代償は互いに異なり、一つの代償に畳めない。畳まずに一組で残すことが、本節の構成の理由である。知覚と言語模型の来歴はTS-003に属し、本特集TS-002の生成条件史には数えない\autocite{btd094}。境界の置き方は本特集の編集方針の一部である。TS-002が扱うのは、表現と予測と生成目的と条件づけと制御と編集と時間構造とサンプリングと実行と評価の選択の進化であり、知覚と言語の理解の史はTS-003の領域である。境界をまたぐ拡張は行わず、境界上の事実は境界上の事実として残す。音声と音楽の生成や音声条件の任意モダリティ間変換が本文で確立していない箇所は、その旨を明示する。音声と音楽の生成や音声条件の任意モダリティ間変換が本文で確立していない箇所は、その旨を明示する。来歴機構については、C2PAは主催頁の範囲、SynthIDは製品頁の範囲でのみ引用し、仕様本体や独立評価の言語では記述しない。収束の主張は問いとして立てたまま残す。それが本節の読みの作法である\autocite{btd094,btd098}。

\subsection{課題別頭部を持たない多課題統合の射程}
分業から統合への第一歩として、Unified-IOは課題別頭部を持たない一つの系列変換で七種以上の視覚言語課題を扱った。生成関連では文章画像と穴埋めと領域画像を含み、入力出力を離散トークンに均質化する。文章はSentencePiece指示、濃密地図と画像はVQ-GAN符号、箱と点は千区分位置トークンとし、純粋T5符号復号器に二次元片と位置符号で与え、スパン雑音一五％と画像雑音七五％の事前学習の後に八十超の資料を温度混合で多課題学習し、評価用微調整を行わない\autocite{btd094}。報告条件ではGRIT試験平均六四・三でXLがGPV-2三二・〇との差を示し、七種完全対応の例とされ、NYUv2誤差〇・三八五はVQVAE律速、ImageNet七九・一、COCO CIDEr一二六・八、SmallからXLの単調改善で飽和なしとされた\autocite{btd094}。ただしVQ-GANが濃密出力精度を抑え、純粋Transformerは金字塔と損失事前を欠き、Unrestricted軌道に限られる\autocite{btd094}。律速の読み方も条件付きである。NYUv2誤差〇・三八五という値はVQVAE律速という天井付きであり、符号器を変えれば変わる。GRIT試験平均六四・三という値はUnrestricted軌道という条件付きであり、軌道を変えた一般化は本文で確立していない。SmallからXLの単調改善という傾向も、飽和なしという観察付きであり、外挿の保証ではない\autocite{btd094}。音声と音楽の生成や音声条件の任意モダリティ間変換は本文で確立せず、知覚と言語模型の来歴はTS-003に属して本特集の生成条件史には数えない。ここで境界を明示する。本節がTS-002として扱うのは画像と文章の生成課題の範囲であり、音声を条件とする任意モダリティ間変換や知覚史の評価はTS-003の領域として分け、拡張して数えない。収束の問いは開いたまま、条件付きの統合例として位置づける。この統合は多課題の被覆を示すが、生成の忠実度や長時間性や音声条件を保証しない。VQ-GAN律速と軌道限定の二点を添えて引用し、知覚史の来歴に拡張しない。

\subsection{話し言葉と文章の往復を一つの復号器に畳む}
話し言葉と文章の往復を一つの復号器に畳む変化が、AudioPaLMである。PaLMとPaLM-2のSentencePiece語彙に多言語w2v-BERTないしUSM由来の千二十四音声トークンを二十五Hzで拡張し、全重みを札付き混合で学ぶ。混合にはASRとASTとS2STを一復号思考連鎖に束ねた形を含み、復号はAudioLM第二段と第三段か百分の一高速の並列SoundStormで波形化し、三秒声プロンプトで言語越え声質転移を行う\autocite{btd095}。復号の読み方も条件付きである。AudioLM第二段と第三段という経路、百分の一高速という並列SoundStormの速度、三秒という声プロンプトの長さの三点を外して、音声化一般の姿として語らない。速度の値は品質の値ではなく、声質転移の三秒条件は短尺の保証ではない。復号器の高速化を品質の向上と混同しない\autocite{btd095}。報告条件ではASTでAudioPaLM-2がAudioPaLM比で観測AST二八％増、観測ASR一〇七％増を示し、FLEURSゼロショットの観測ASRで二〇・七BLEUを示してWhisper-Large-v2-1.5Bの一九・六との差が同一観測手順で示され、音声間翻訳はASR経由BLEUと声質保持でカスケードに対する差が報告条件で示された。学習はAdafactor五e-5で数千時間級の混合とされる\autocite{btd095}。ただし発話音声と文章のみで画像と映像と音楽を含まず、翻訳目標は合成が多く、三秒プロンプト転移は短尺で劣化する\autocite{btd095}。射程の読み方も条件付きである。観測AST二八％増と観測ASR一〇七％増という値は、同一観測手順での差であり、手順を変えた一般化ではない。FLEURSゼロショット二〇・七BLEUという値は観測ASRという経路付きであり、直接音声間翻訳の値ではない。数千時間級の混合という学習規模も条件の一部であり、規模を変えた一般化は本文で確立していない\autocite{btd095}。記録locatorの転記欠陥は正規二三〇六・一二九二五側本文で消費した旨を添え、統合の射程を合成目標付きで読む。音声と文章の往復に限る射程を外した拡張は行わず、合成目標の多さという条件を添えて引用する。声質転移の三秒条件を短尺の保証に読み替えず、復号器の高速化を品質の向上と混同しない。分離カスケードからの移動を示す事例として残し、画像や映像への一般化は確立しないと記す。
\end{multicols}

\begin{center}
\small
\begin{tabularx}{\textwidth}{@{}p{0.18\textwidth}p{0.34\textwidth}X@{}}
\toprule
立場 & 設計の要点 & 射程と天井 \\
\midrule
多課題統合 & 離散均質化・課題別頭部なし・温度混合多課題学習 & 画像文章の生成課題・VQ-GAN律速・軌道限定 \\
発話音声と文章の往復 & 語彙拡張・札付混合・一復号思考連鎖 & 発話音声と文章のみ・合成目標多・短尺転移は劣化 \\
資料側離散化 & 構造不変・様式別離散化・文章橋渡し整列 & 多様式損失は単様式を上回る・五秒上限・専用評価なし \\
合成拡散 & 様式別潜在拡散・重み補間融合・線形拡張 & 見出し欠落を継承・結合指標は整合・時間修正は部分的 \\
\bottomrule
\end{tabularx}
\end{center}
\autocite{btd094,btd095,btd096,btd097}

\begin{multicols}{2}
\raggedcolumns
\subsection{構造不変の資料側工夫と合成拡散の対比}
構造も学習も変えず資料側離散化のみでany-to-any会話を狙うのが、AnyGPTの立場である。画像はSEED、発話音声はSpeechTokenizerのRVQ意味第一層とSoundStorm音響、音楽はEncodecをフレーム毎平坦化して離散化し、LLaMA-2-7Bを二兆文章トークンの土台に文章橋渡し整列とAnyInstruct十万八千件で次トークン学習し、復号は拡散とSoundStormとEncodecで行う\autocite{btd096}。離散化の読み方も条件付きである。SEEDという画像符号、RVQ意味第一層という発話音声符号、フレーム毎平坦化という音楽符号の三点を外して、資料側工夫一般の姿として語らない。二兆文章トークンという土台、十万八千件という整列規模も条件の一部であり、規模を変えた一般化は本文で確立していない。復号の割当も様式毎に異なり、割当を変えた一般化は確立していない\autocite{btd096}。報告条件のゼロショット素体では画像説明CIDEr一〇七・五でSEED-LLaMA一二三・六との差、文章画像CLIP〇・六五で〇・六九との差、Libri清浄ASR誤り八・五でWhisper-L2二・七や人五・八との差、VCTK合成誤り八・五と類似〇・七七でVALL-E七・九と〇・七五との差、MusicCaps理解〇・一一で実〇・一六、生成〇・一四でMousai〇・二三との差が同一ゼロショット手順で示され、音楽は五秒上限とされた\autocite{btd096}。ただし多様式損失は単様式を上回り、トークン化が能力天井となり、長系列文脈が対話深度を抑え、専用のany-to-any評価はない\autocite{btd096}。天井の読み方も条件付きである。二兆文章トークンの土台とAnyInstruct十万八千件という学習条件、SEEDとSpeechTokenizerとEncodecという離散化の選択、拡散とSoundStormとEncodecという復号の割当の三点を外して、任意モダリティ間一般の姿として語らない。ゼロショット素体という条件も値と一体であり、調整後の姿への外挿は行わない\autocite{btd096}。記録locatorの転記欠陥は正規二四〇二・一二二二六側本文で消費した旨を添え、素体値の比較は条件固定の差として読む。五秒上限と高損失と天井の三点を添えて引用し、素体のゼロショット値を調整後の姿と混同しない。拡散を合成してany-to-any生成に寄せるのが、CoDiの設計である。画像はSD1.5、映像は擬時間と潜在シフト、音声はメル画像化とAudioLDM-VAEとボコーダ、文章はOPTIMUSとGPT-2という様式別潜在拡散を独立に置き、文章橋渡しプロンプト符号化を重み補間で融合する。結合生成は凍結模型に交差注意と対比整列の環境符号化を加え、文章画像から文章音声へ映像音声へと線形に拡張して未見組合せを可能にする\autocite{btd097}。拡張の読み方も条件付きである。凍結模型という前提、交差注意と対比整列という追加、文章橋渡しプロンプト符号化の重み補間融合という手順の三点を外して、合成一般の姿として語らない。文章画像から文章音声へ映像音声へという順序は線形拡張の設計であり、順序を変えた一般化は本文で確立していない。未見組合せの可能性は設計の主張であり、未見組合せの品質の確定ではない\autocite{btd097}。報告条件ではAudioCaps文章音声でFD二二・九〇を示し、AudioLDM-L-Full二三・三一やIS八・七七対八・一三やKL一・四〇対一・五九やFAD一・八〇対一・九六との差が同一集合で示され、文章画像COCOでFID一一・二六対SD1.4一一・二一、結合類似度は映像音声〇・二五五対独立〇・二四〇等の上乗せ、説明はほぼ同等とされた\autocite{btd097}。ただし文章橋渡しは見出し欠落を継承し、結合指標は符号化余弦の整合であって忠実度ではなく、映像時間整合の修正は部分的である\autocite{btd097}。整合の読み方も条件付きである。映像音声〇・二五五対独立〇・二四〇という上乗せは、符号化余弦という一つの軸での差であり、忠実度の軸での差ではない。AudioCapsという集合、COCOという集合の二点を外して、結合生成一般の姿として語らない。文章画像から文章音声へ映像音声へという線形拡張の順序も条件の一部であり、順序を変えた一般化は本文で確立していない\autocite{btd097}。記録locatorの転記欠陥は正規二三〇五・一一八四六側本文で消費した旨を添え、整合と忠実を分けて読む。結合の上乗せを忠実度の向上と呼ばず、符号化余弦の整合改善として記録する。指数級対資料を回避する工夫として位置づけ、時間整合の部分的修正という天井を明示して継承する。

\subsection{来歴機構の頁限定と開かれた問い}
来歴機構を仕様と運用の水準で切り分けるため、C2PAとSynthIDは主張の粒度を変えて扱う。C2PAは公開来歴標準として内容証明書の履歴を栄養表示になぞらえ、利用者が随時検証できると主催頁は述べ、運営はAdobeやAmazonやBBCやGoogleやMetaやMicrosoftやOpenAIやSonyやTikTokやTruepicを含む多者委員会とされる\autocite{btd098}。運営の読み方も条件付きである。多者委員会という構成は来歴表示の運用体制を示すが、生成物の来歴保証を示さない。委員の顔ぶれは主催頁の記述時点のものであり、参加の変動や仕様の版の固定は頁の範囲では確立していない。運営の厚みを仕様の確定や配備の証拠に読み替えない。ただしここで明示する。本節のC2PAに関する記述は主催頁の範囲の消費に基づくものであり、技術仕様の宣言と署名と改竄証跡と信頼一覧、音声と音楽の束縛の詳細は未消費である\autocite{btd098}。主催頁の消費とは、表層の説明と仕組みへの案内と仕様採用への誘導と運営委員会の一覧の確認を指し、仕様本体の宣言と主張とハッシュと署名と束縛、音声固有の扱いは未消費である。配備証拠は提供元と仕様源泉に委ねられ、生成関連の保証や署名や改竄証跡や音声束縛や検出頑健は確立していない\autocite{btd098}。主催頁の確認は来歴表示の運用姿勢を示すが、生成物の来歴保証を示さない。SynthIDは作成時に画素とフレームと波形と言語模型次トークン確率skewに透かしを埋め、切抜とフィルタとフレーム率と圧縮、雑音とMP3と速度の常用編集に頑健と頁は主張し、検出器とGemini持込検査で確かめ、報道向け早期試験とされる\autocite{btd099}。ただし製品頁のみで独立頑健評価がなく、音声はLyriaとNotebookLMに紐付き、文章透かしの品質影響は未定量化で、検出率と曲線と容量と音楽固有評価は頁にない\autocite{btd099}。頁限定の読み方も条件付きである。切抜とフィルタとフレーム率と圧縮、雑音とMP3と速度という常用編集の列挙は、頁主張の範囲を示すが、独立評価の範囲を示さない。検出器とGemini持込検査という確認手順も頁の記述であり、第三者の再現ではない。報道向け早期試験という位置づけも頁の記述であり、配備の証拠ではない\autocite{btd099}。頁主張の頑健さを独立評価と呼ばない。音声束縛や検出率や容量の未確立を明示し、配備証拠の不足を残す\autocite{btd098,btd099}。\subsection{問いを閉じない結びの作法}
単一統合か連携専門家群かの結論は出さず、問いは開いたまま残す\autocite{btd094,btd095,btd096,btd097,btd098,btd099}。本項ではこの開き方を手順化する。第一に、多課題の被覆を示す一例を、生成の忠実度や長時間性や音声条件の保証に格上げしない。第二に、発話音声と文章の往復に限る射程を、画像や映像への一般化に拡張しない。第三に、資料側工夫と合成拡散の対比を、整合と忠実の別で整理し、結合の上乗せを忠実度の向上と呼ばない。第四に、来歴機構の頁主張を、独立評価や配備証拠に格上げしない。第五に、TS-002の生成条件史に属する事実とTS-003に属する知覚史を分けたままにする。生成の来歴保証と閲覧時の検証可能性を混同せず、TS-002の生成条件史に属する事実とTS-003に属する知覚史を分けたまま、問いを開いて残す\autocite{btd094,btd095,btd096,btd097,btd098,btd099}。統合か連携かの二択を閉じず、条件付きの一例と頁限定の主張と未確立の範囲を並置する。それが収束問題に対する本特集の姿勢であり、本節の結びの作法である。結びの並置を再確認する。多課題統合は画像文章の生成課題の条件付き一例、発話音声と文章の往復は合成目標付きの条件付き一例、資料側離散化は五秒上限付きの条件付き一例、合成拡散は整合と忠実の分離付きの条件付き一例、来歴二者は頁限定の主張である。五点を一つの結論に畳まず、開いたまま残す\autocite{btd094,btd095,btd096,btd097,btd098,btd099}。
\end{multicols}

\begin{center}
\small
\begin{tabularx}{\textwidth}{@{}p{0.20\textwidth}p{0.36\textwidth}X@{}}
\toprule
来歴機構 & 頁の主張の範囲 & 未確立の範囲 \\
\midrule
C2PA（主催頁のみ） & 履歴の栄養表示・随時検証・多者委員会運営 & 宣言と署名と改竄証跡と信頼一覧・音声音楽束縛・配備証拠 \\
SynthID（製品頁のみ） & 画素・フレーム・波形・確率skew埋込・常用編集に頑健 & 検出率と曲線と容量・音楽固有評価・品質影響の定量化 \\
\bottomrule
\end{tabularx}
\end{center}
\autocite{btd098,btd099}

\begin{claimboundary}[本節の読解上の境界]
知覚と言語模型の来歴は除外し、収束の主張は問いとして立てたまま残す。Unified-IOはVQ-GAN律速で軌道限定であり、音声と音楽の生成や音声条件の任意モダリティ間変換は本文で確立しない。AudioPaLMは発話音声と文章のみで翻訳目標は合成が多く、短尺プロンプト転移は劣化する。AnyGPTは多様式損失が単様式を上回り、五秒上限と長系列文脈の制約があり、専用のany-to-any評価はない。CoDiの結合指標は符号化余弦の整合であって忠実度ではなく、時間整合の修正は部分的である。C2PAに関する記述は主催頁の範囲に留まり、仕様本体の宣言と署名と改竄証跡と音声音楽束縛は未消費で、生成関連の来歴保証は確立していない。SynthIDは製品頁のみで独立頑健評価がなく、検出率と誤検出率と容量と音楽固有の有効性は確立していない。記録locatorの転記欠陥は正規側本文で消費した旨を添える。
\end{claimboundary}

\section{2025--2026年の到達点：能力・workflow・lifecycleの証拠}
\sectionkicker{CAPSTONES}
\label{sec:capstones}

\begin{multicols}{2}
\raggedcolumns
本節は2025--2026年の現行群を、能力・workflow・可用性・lifecycleの証拠としてのみ扱う。構成・学習・符号化・評価手法への推論は行わない。製品挙動からの機構推測は禁じられる。性能の順位付けは行わない。数値は条件付きの記録であり、条件の異なる数値を横断比較しない。各記述は版と日付にバインドする。絶対日付を開示しない頁は日付非開示として明示する。ベンダー由来の速度・品質・規模の表現は、ベンダー主張として帰属付きで引用し、独立測定を待つ主張として隔離する。本節は前節までの機構史を置き換えるものではなく、現場で何が使え、何が手順化され、何が提供終了・凍結・移管されたかを、閉じた製品表面の記録として残すものである。

本節の契約を最初に固定する。第一に、能力・workflow・lifecycleのみを語り、構成・学習・符号化・評価手法の推論には進まない。第二に、ベンダー由来の表示はベンダー主張として帰属付けし、独立測定までは運用判断の根拠にしない。第三に、版と日付のバインドを欠いた引用は行わず、日付非開示の頁はその旨を明示する。第四に、条件の異なる数値を並べて順位を作らない。第五に、抄録限定・lifecycle限定のPARTIALの範囲を越えない。抄録限定の典拠は抄録主張の範囲であり、lifecycle限定の典拠はlifecycle文脈の範囲である\autocite{btd076}\autocite{btd106}。この五点を守ることで、capstoneは機構史の権威を損なわず、現在位置の補助線として機能する。

\subsection{画像：生成と反復編集のworkflow表面}
FLUX.2 [klein](Black Forest Labs、2026年1月)は、文章からの画像生成と単一・複数参照による編集を一つの提供形態に統合した能力事例である\autocite{btd100}。0.5秒を切る推論や40億規模で約13GB VRAMでの動作、5倍規模の模型に匹敵あるいは上回るとの表示、他系とのElo対遅延・VRAMのPareto図表は、いずれもベンダー主張として引用し、独立測定のない順位の根拠にはしない。対になるヘルプ記事「What is FLUX.2」(日付非開示の動的文書)は、最大8--10枚の複数参照、最大4MP(例2048$\times$2048)までの生成、改善された文字描画といった製品表面を記録する\autocite{btd101}。文書は表面の権威にとどまり、構成・学習・符号化への推論には使わない。版と日付のバインドは再利用時の再確認を条件とする。

画像の設計・反復編集workflowでは、Seedream 5.0 Pro(ByteDance Seed、2026年7月)、ChatGPT Images 2.5(OpenAI、2026年9月)、Nano Banana 2(Google、2026年2月)が異なる能力表面を示す\autocite{btd102}\autocite{btd103}\autocite{btd104}。Seedreamは濃密な図表や可視化の生成、点・投げ縄・矩形・素描による指定、16進色や質感の編集、層分離を文書化するが、評価はベンダー主張の段階にとどまり独立裏づけを待つ。ChatGPT Imagesは反復編集(Sketch・雛形・画面上注記・共有)と、Flare既定・Sunburst密着のAPI模型を備え、従来比最大5割の低遅延をうたうが、これもベンダー主張として帰属付けし、開示限定のため機構への推論は行わない。Nano Bananaは世界知識や検索接地を伴う生成編集を掲げるが、生成器内の推論の主張は模型票の権威を要するため本節では能力表面の記録にとどめる。Gemini 3.1 Flash Image模型票(2026年2月)は多峰入出力と人文横並べ・自動評価の表を備える一方、幻覚や小密文字や左右混乱などの制約を票内に列挙し、票外への横断比較を禁じる\autocite{btd105}。票内条件のベンダー評価であり、票外模型との優劣確定には使わない。いずれも能力workflowの証拠であり、知覚史の側へは踏み出さない。

\subsection{音声・対話：実時間利用と版バインドの機能集合}
音声の実時間・多能化では、GPT-Live(OpenAI、2026年7月)、API音声模型群(2026年5月)、Gemini 3.8 Audio模型票(2026年9月)、Seed Audio 1.0(ByteDance Seed、2026年7月)が能力事例となる\autocite{btd107}\autocite{btd108}\autocite{btd109}\autocite{btd110}。GPT-Liveは同時聴話・割込み処理・相槌・生翻訳を伴う全二重対話を文書化するが、実時間率の数値はAPI文書へのバインドを要し本節では確定しない。API表面は実時間系・翻訳系・ストリーミングの版バインドされた機能集合であり、前置指示・並列道具・多段階推論・翻訳・streamingの各機能は版ごとに記録される。価格・遅延の数値は版従属として再利用時の再確認を要する。Gemini音声票は生音声の入出力とLive・思考拡張・音声合成の提供形態、文脈量と出力上限を記すが、票範囲の権威にとどまり独立評価を欠く。Seed Audioは発話・効果音・環境音の場面単位の統合生成と100ms台のtiming制御を掲げ、A/B選好や9割超の可用率や多言語の自然さ得点をベンダー図表で示すが、場面整合の主張はベンダー主張として引用し独立評価待ちとする。遅延・自然さ・類似度・韻律の混同を避けて読む。方言や一部言語の例外を含む限定付き主張であり、条件外への一般化は行わない。

\subsection{音楽：制作手順としてのworkflow記録}
音楽では閉鎖模型のなかで開示の厚いLyria 3.5模型票(Google DeepMind、2026年7月)が軸となる\autocite{btd111}。票は歌詞付き音楽の生成と流派・気分・楽器・歌声の制御、音楽専門家を交えた人文・自動評価による先行版からの改善を記すが、重みは閉鎖であり、構造主張は票の手法にバインドされる範囲でのみ読む。機構全体への推論は行わない。対になるLyria hub(日付非開示の製品系譜)は、60秒断片から3分楽曲までのduration制御や音声言語流派音響の制御など可用性表面を記すにとどまり、機構の主張は票へ委譲する\autocite{btd112}。hubの日付非開示は動的文書の性質であり、再利用時は版バインドの再確認を要する。Stable Audio 3(2026年5月)は可変長生成と続成・補完、少数段階化、局所実行を文書化するが、機能表面は版バインドであり公開範囲は許諾に従属する\autocite{btd113}。三件とも能力workflowの証拠であり、符号化や学習の全体像への推論は行わない。得点の横並び比較は条件不整合のため禁じられる。

商用歌曲workflowでは、SunoとElevenLabs Musicが機構非開示のまま制作表面の証拠となる。Suno release notes(頁自体は単一日付を持たない動的索引であり、日付非開示の頁本体に項目別日付がバインドする)は、作成(文章・音声・画像・映像)、歌曲編集・置換・刈込・Stem・Extend・Remix・Cover・Persona、Studio DAWという制作手順の全体を記録する\autocite{btd114}。Suno v6(2026年9月)はv6・v6-wild・v6-miniの三層と自然文区間編集・複数源mashup・試料分離・拍workflowを告知するが、主張はベンダー主張として引用し独立検証を欠く\autocite{btd115}。ElevenLabs Music文書(日付非開示の能力文書)は、器楽・歌唱・多言語の生成、30秒音響参照、短時間のFinetune、長尺作曲と補完を記すが、評価手法は薄い文書段階にとどまる\autocite{btd116}。三件とも能力workflowの権威のみを持ち、構成・学習・評価の権威にはならない。長尺構造の指標は未確立であり、受容側の記録と接続する。

\subsection{映像：joint化と短区間・長距離の分離}
映像のjoint化では、抄録段階のMovie Genと2026年の現行群を、開示の厚みの差を明示して並べる。Movie Gen(Meta、2024年10月)は抄録のみ消費のPARTIALであり、鋳造・1080p・同期音・指示編集・個人化という抄録主張のみが権威となる\autocite{btd076}。本節では抄録水準を超える記述は行わない。30B・13B・7万時間・16秒級の数値はベンダー主張であり、手法・評価節はscript遮蔽で未消費のため一切の機構推論を禁じる。Seedance(製品表面、版バインド)は30画像・10映像・10音響の参照や照明制御、時刻指定・合成・視点の編集を備え、命名応用での展開とAPI予告を含む製品表面を持つが、物理の妥当性の余地はベンダー自身が認める限定である\autocite{btd117}。Seedance 2.5(2026年7月)は音声映像のjoint生成、30秒単走と多回合延長による長尺化、多数の画像映像音響参照、時刻・合成・視点の編集を文書化するが、持続整合の主張は独立検証待ちのベンダー主張として引用する\autocite{btd118}。Veo(日付非開示の頁)は対話・効果音・環境音を伴う生成と参照素材・場面延長・首尾frameを備え、千余件規模のbenchやVBench系のベンダー性能を掲げるが再現待ちの帰属付き主張である\autocite{btd119}。Veoの参照素材・様式・人物、場面延長、首尾frame、外側描画は制作制御の表面であり、対話・効果音・環境音の統合は短区間の忠実度と長距離の保持を分けて読む。開示断片以外の機構言及は禁じられる。最大対応durationは実証された整合と混同しない。

\subsection{lifecycle：提供終了・廃止・凍結の事実}
lifecycle転換は製品方向の事実として残し、構成上の必然性の証拠にはしない。Imagen頁はlifecycleのみのPARTIALであり、Nano Banana hubへの転送後の姿のみが権威となる\autocite{btd106}。本節ではImagenをlifecycle文脈でのみ用い、能力の詳細はこの所在からは確立しない。2026年の廃止は専用生成器から多能hubへの製品移動の事実であり、能力・構成・学習の証拠にはしない。Gen-4.5(日付非開示の頁)は複雑な連続指示と撮影振付の解釈を掲げるが、追加入力は予告段階であり、料金・秒歩は版バインドのworkflow記録である\autocite{btd121}。頁日付と機能の版バインドを混同しない。Ray3(日付非開示の頁)は文章・画像起点と映像間変換(人物参照・首尾keyframe含む)、衣装・環境・光・配置のModifyを備えるが、頁内bannerと記事の日付は区別して版バインドする\autocite{btd122}。2026年内のsubversion driftのためRay 3.2(2026年6月:16 keyframe・8顔追跡・HDR/EXR)と束ねて版バインドする。Sora 2(2025年9月)は2026年4月26日の提供終了を記す歴史頁であり、終了案内(日付非開示の頁)の内容日と輸出・削除対応、API廃止表(頁自体は日付非開示で、行日2026年9月24日の削除・代替なしを記録)を、頁日付と混同せずに記録する\autocite{btd123}\autocite{btd126}\autocite{btd127}\autocite{btd128}。Sora 2は物理整合の多shot・世界状態保持・様式・音景を掲げる旗艦頁だが、提供終了後は歴史・lifecycleの役割に限る。終了案内の輸出・削除対応、API表の予告・終了・代替なしは行単位のlifecycle境界である。開放凍結・API専念・廃止の方向は製品の事実であり、統合への必然性の証拠にはしない。統合か連携かは未決の問いとして残る。

\subsection{横断の読み方：capstoneを大きく見せない}
本節の27件は、前節までの機構史に対する追加の権威ではなく、利用・手順・提供の側面からの補助線である\autocite{btd100}\autocite{btd107}\autocite{btd111}\autocite{btd076}。画像では参照編集と反復注記が、手順の単位として定着したことが要点である。Seedream系の指定手法(点・投げ縄・矩形・素描)やChatGPT Images系の注記・共有の反復、FLUX.2系の多参照の上限、Nano Banana系の検索接地の提供形態は、いずれも作業の進め方の記録であり、生成の仕組みの記録ではない\autocite{btd101}\autocite{btd102}\autocite{btd103}\autocite{btd104}。文字描画の改善や層分離の表示も、製品表面の変化としてのみ読む。ベンダー主張として引用された遅延・VRAM・Eloの図表は、条件付きの表示であり、現場の再現測定がなければ運用判断の根拠にしない。

可用性の読み方を補う。FLUX.2系の文書化された構成での動作表示、Seedream系の文書化された指定・編集の範囲、ChatGPT Images系の反復注記とAPI模型の版、Nano Banana系の多表面での提供と検証回数の言及、Gemini票の票内制約の列挙は、いずれも提供条件の記録である\autocite{btd100}\autocite{btd101}\autocite{btd102}\autocite{btd103}\autocite{btd104}\autocite{btd105}。開放重みと閉鎖APIの実行境界を示す一例としてのみ読み、構成や学習への手がかりにはしない。ヘルプ記事やhubなどの動的文書は日付非開示の性質を持ち、引用のたびに版バインドの再確認を要する。公開・軽量実行という配置は、実行可能性の存在の記録であり、品質の確定ではない。

音声・音楽では、場面単位の生成と制作手順の定型化が要点である。Seed Audio系の場面単位の統合や100ms台のtiming制御、Lyria系の流派・気分・楽器・歌声の制御、Suno系のStem・Extend・Remix・Cover・PersonaとStudio DAW、ElevenLabs系の30秒参照と短時間Finetuneは、制作の単位の記録である\autocite{btd110}\autocite{btd111}\autocite{btd114}\autocite{btd115}\autocite{btd116}。Lyria hubの60秒断片から3分楽曲までのduration制御や画像起点作曲、Stable Audio系の可変長生成と補完・続成・局所実行も、可用性の記録であり、符号化や学習の証拠にはしない\autocite{btd112}\autocite{btd113}。A/B選好や可用率や自然さ得点のベンダー図表は、ベンダー主張として引用し、方言や一部言語の例外などの限定とともに読む。長尺構造や歌詞遵守の達成は未確立であり、受容側の未決laneと接続する。

映像・lifecycleでは、短区間と長距離の分離と、提供の移動の記録が要点である。Seedance系の30秒単走と多回合延長、Veo系の場面延長と首尾frame、FLUX系の上限10秒・720p等の包絡は、版バインドの提供条件として読む\autocite{btd117}\autocite{btd118}\autocite{btd119}。同期音・対話・効果音の統合は、短区間の忠実度と長距離の保持を分けて読む。同一性永続・長距離構造・継続・編集保存は未確立であり、製品表面の上限表示を達成の証拠にしない。Imagenの廃止、Sora 2の提供終了と輸出・削除対応、API廃止表の行単位の削除・代替なし、Ray系のbannerと記事の日付の区別とsubversion driftへの対応は、製品方向の事実としてのみ読む\autocite{btd106}\autocite{btd121}\autocite{btd122}\autocite{btd123}\autocite{btd126}\autocite{btd127}\autocite{btd128}。日付非開示の頁は頁日付と項目・行日を混同せず、再利用時は版バインドの再確認を条件とする。capstoneは機構史を大きく見せるための章ではなく、座標の現在位置を示す補助線である。

\begin{themeoverview}[27件の配置見取り図]
画像6件は生成・反復編集の手順、音声4件は実時間・版バインドの機能集合、音楽6件は制作手順の定型化、映像4件はjoint化の提供条件、lifecycle 7件は提供終了・廃止・凍結の行単位の境界である。いずれも能力・workflow・lifecycleの記録であり、構成・学習・符号化の記録ではない。
\end{themeoverview}
\autocite{btd100}\autocite{btd107}\autocite{btd111}\autocite{btd076}\autocite{btd106}
\end{multicols}

\begin{claimboundary}[本節の読解上の境界]
本節は能力・workflow・lifecycleの記録に限定し、構成・学習・符号化の推論は行わない。ベンダー由来の数値・順位・速度の表示はベンダー主張として帰属付きで引用し、独立測定までは事実化しない。条件の異なる数値の横断比較は行わない。抄録のみ消費のPARTIAL典拠は抄録主張を超える記述を行わない。lifecycleのみのPARTIAL典拠は、能力・構成の証拠にはしない。日付非開示の頁（ヘルプ記事・ハブ・索引本体・能力文書・モデルハブ・ヘルプ・ハブ・終了案内・廃止表本体・API表関連）は頁日付と項目・行日のバインドを区別する。版バインドは再利用時の再確認を条件とする。
\end{claimboundary}

\section{受容とcounter-signal：現場の27件の記録}
\sectionkicker{RECEPTION}
\label{sec:reception}

\begin{multicols}{2}
\raggedcolumns
本節は現場の27件を、一つのboundedな受容記録として読む。27件は27の技術権威ではなく、単一のDiscovery記録に束ねられた1件の受容証拠である\autocite{btd139}。供給するのは受容・配備・失敗の目撃のみであり、構成・日付・許諾・価格・得点・因果の権威にはならない。人気による順位付けは行わない。条件を越えた得点比較は行わない。X由来の技術的後続は一次・独立権威へ差し戻すまで先導扱いであり、台帳の算術を超える主張は禁じられる。

受容台帳r3は27記録・27URL・26勘定(独立22・ベンダー系2・実装保守1・不明1)から成る\autocite{btd139}。FIRST\_HANDは22・1・4、受容は肯定13・混合5・否定5・中立4、内訳は画像11・音声4・音楽3・映像8・機構1(副次札1)である。画像ではFLUX.2 kleinのMLX・Ollama・ComfyUIによる局所・持込実行への関心と遅延摩擦、軽量模型への移行信号、残差の文字誤りが並び、Seedream系の人物整合と参照workflowが好評として記録される。音声4件は全二重の解説、Apple Silicon上実行の関心、F5-TTSを基準とする局所学習主張、雲代替の位置づけであり、整合した実時間率の測定は薄い\autocite{btd139}。いずれも配備・作業・失敗の目撃であり、権威の昇格はない。分類は制作者workflow 9・配備4・bench 3・失敗3・保守1・独立分析3・採用信号2・公式1(副次counter 1)である。

音楽3件は制作者管への導入を示し、短尺の好評が記録される一方、長尺構造の失敗証拠は得られていない\autocite{btd139}。映像8件は並列比較における整合・物理の選好、顔置換での同一性破損、GGUF・ComfyUIにおける記憶・磁碟・streaming制約、口同期と長距離連続の実務課題、強否定の1件から成る。機構1件に副次札を加え、時間空間整合の上限と長距離driftのcounter-signalを保持する。失敗台帳は、長時間の同一性喪失・接合部の同期破綻・非対象域の改変・局所映像の資源制約を確定側に置き、参照誠実度を部分確定、少数段階の質交換を混合と判定する\autocite{btd139}。副次札は二重計上せず、一次分類27の算術を崩さない。比較は条件付きの目撃であり、bench得点の代替にはならない。

\subsection{台帳の読み方：確定側と未決側の分離}
確定側に置くのは、目撃として成立した受容・配備・失敗である。軽量実行への関心と遅延摩擦、人物整合と参照workflowの好評、全二重の採用関心、短尺音楽の好評、並列比較の選好、顔置換の破損、局所実行の資源制約、口同期と連続の実務課題、時間空間整合の上限とdriftの是認がこれに当たる\autocite{btd139}。いずれも技術の仕組みの証拠ではなく、現場での使われ方と壊れ方の証拠である。未決側に置くのは、測定と権威を欠く事項である。実時間率の整合測定、長尺構造と歌詞整合、自動指標と人選好の対応、少数段階の質交換の条件、flow対拡散の純生成比較、消費者 klein再現、独立評価、Wan 2.5以降の権威、Sora機構がこれに当たる\autocite{btd139}。確定側の充実が未決側の解消を意味しないことを厳守し、件数の多さを確からしさに読み替えない。受容記録は技術anchorの配備・実行読解に資するが、権威の昇格はないという一点に尽きる。

\begin{themeoverview}[受容27件の見取り図]
画像11・音声4・音楽3・映像8・機構1(副次札1)を1件のbounded記録として束ねる。確定側は配備・作業・失敗の目撃、未決側はLOW\_SIGNALのまま開いて残す。人気順位の導出と条件越しの比較は禁じられる。
\end{themeoverview}
\autocite{btd139}

\subsection{境界と未決laneの保持}
X境界は厳格に保つ。Xは受容・配備・失敗の証拠のみを供給し、構成・日付・許諾・価格・得点・因果を立てない\autocite{btd139}。X由来の技術的後続は一次・独立権威へ差し戻すまで先導扱いであり、台帳の算術を超える主張は禁じられる。並列比較の選好は条件付きの目撃であり、bench得点の代替にはならない。局所timingやGGUF手順の記録は、配備の存在の記録であり、品質の確定ではない。音声誤りの主張やベンダー人文言及への言及は、先導としてのみ残し、一次資料への差し戻しを条件とする\autocite{btd139}。独立bench・再現の先導(並列比較・局所timing・GGUF手順・音声誤り主張・ベンダー人文言及)が残ることは、不足ではなく bounded な一次通過の当然の帰結である。

LOW\_SIGNAL laneは開いたまま残す。全二重の割込み・重畳測定、多言語・感情複製の劣化、音楽の長尺構造と歌詞整合、自動指標対人選好の矛盾、編集corpus、少数段階ablation、純生成のflow対拡散、消費者 klein再現、独立評価、Wan 2.5以降の権威、Sora機構はいずれも未決であり、膨らませも順位付けもしない。疎laneの欠測は欠陥ではなく、boundedな一次通過の当然の帰結として記録する。技術的後続の必要が生じた場合は、Xの所在ではなく一次資料の所在をauthorityとして記録し直す。受容2頁相当の確定側と本頁の未決側を分けて保持し、後続の起草が境界を踏み越えないよう固定する。27件の算術と台帳の文言が唯一の根拠である\autocite{btd139}。
\end{multicols}

\begin{claimboundary}[本節の読解上の境界]
人気による順位付けは行わない。Xは構成・日付・許諾・価格・得点・因果の権威にならない。27件は27の技術権威ではなく1件のboundedな受容記録である。X由来の技術的後続は一次・独立権威への差し戻しを経るまで仮説扱いとする。LOW\_SIGNAL laneは開いたまま残す。
\end{claimboundary}

\section{結び──座標として読むメディア生成史}
\sectionkicker{SYNTHESIS}
\label{sec:synthesis}

\begin{multicols}{2}
\raggedcolumns
本巻がたどったのは、高次元連続メディア(画素・波形・frame・動き・同期音)を生成可能にする選択の積み重ねである。表現の短縮を起点に、目的関数の分岐、条件づけと案内の分離、制御と参照と同一性、生成と編集の分離が重なり、音声・音楽・映像の各系譜で時間構造と実行経済と評価妥当性が分化した。閉鎖製品は能力・workflow・lifecycleの文脈に留め、受容は配備と失敗の証拠として保存する。この層状の読み方が、本巻の持ち帰りである。

層ごとの要点を振り返る。表現の層では、連続潜在・離散符号帳・RVQ神経符号・意味音響階層・時空間圧縮が、raw系列の費用を tractable な系列へ変えた。圧縮率・token rateの主張は設定バインドであり、意味・音響の分離を混同しないことが要点である\autocite{btd002}。目的関数の層では、自己回帰・敵対・拡散・score・flow・整流が分岐し、骨格(U-Net・Transformer・DiT)と目的と手続きの分離が読みの規律となった。各転換は隘路・機構・改善・代償・継承で読む\autocite{btd023}。条件づけの層では、学習時条件づけ・推論時案内・adapter・文脈参照を区別し、prompt遵守を分布品質や知覚品質と混同しないことが要点である\autocite{btd032}。制御の層では、空間制御・参照・主体性保存を分け、制御信号の忠実度と標本品質を別物として読む\autocite{btd039}。編集の層では、新規生成と保存を伴う改変を分け、保存対象・局所性・参照制約・失敗様式を記録する\autocite{btd049}。

時間の層では、短距離整合と長距離構造と同一性永続と同期を分け、単一標本の主張を長距離能力の根拠にしない\autocite{btd058}。実行の層では、段数削減・構造圧縮・携帯最適化を分け、遅延・VRAM・装置・量子化・offloadのバインドなき数値比較を禁じる\autocite{btd066}。評価の層では、分布指標・構成診断・選好・音声映像多面評価を分け、指標の交換可能性を認めない。版・条件・母集団のバインドが要点である\autocite{btd075}。収束の層では、単一統合と専門家連携を未決の問いとして残し、知覚史の側へ踏み出さない\autocite{btd131}。並行する競合関係(連続潜在と離散符号、拡散とflow、学習時条件と推論時案内、新規生成と編集、短距離と長距離、分布と選好、統合と連携)はいずれも万能ではなく、条件・装置・予算・母集団で損得が変わる\autocite{btd078}\autocite{btd087}\autocite{btd096}。

\subsection{現場への持ち帰り}
手法を選ぶ場面では、対象の名指しが助けになる。何を生成するのか(画像・音声・音楽・映像)、新規か編集か、短区間か長尺か、実時間か非同期か、局所か雲か。対象を名指ししないまま製品表面を選ぶと、効くはずの条件と外れる。現行群の読み方は能力表面に限定し、版と日付にバインドする。文字描画や多参照の上限は製品表面の記録であり、学習や符号化の手がかりにはしない\autocite{btd101}。受容記録は配備・実行の読解に資するが、権威の昇格はない\autocite{btd139}。27件の算術を超える主張は行わない。

\subsection{失敗の診断学}
効かなかったときの診断も、層の言葉で行う。第一に疑うのは、対象の取り違えである。prompt遵守と分布品質、制御忠実度と標本品質、新規生成と編集保存、短距離整合と長距離構造、遅延と処理量を取り違えていないかを確かめる。第二に疑うのは、条件の不一致である。版・装置・予算・母集団・duration・参照条件が、想定と現場でずれていないかを確かめる。第三に疑うのは、深さの超過である。要旨・断片・頁限定の事実を本文読解として扱っていないか、ベンダー主張を独立事実として引用していないか、Xの目撃を技術権威にしていないかを確かめる。三つの疑いのいずれにも当てはまらない場合、残る可能性は分布の変化とlifecycleの移動である。提供終了・凍結・移管の確認時期を決め、一次資料への差し戻しを行う。この往復が、本巻の実践である。

\subsection{立場別の指針}
研究の立場では、層の分離を守る。表現・目的・骨格・手続き、条件と案内、制御と参照、生成と編集、短距離と長距離、指標と選好を混同せず、自分の条件での再適合を怠らない。未解決の系譜には手を出さず、並行関係として記録する。運用の立場では、版バインドとlifecycleを先に固める。提供終了・凍結・移管の行単位の境界を確認し、生きた文書の読み方を守り、受容の失敗台帳(同一性喪失・同期破綻・非対象域改変・資源制約)を点検する。意思決定の立場では、総費用と条件付きで見る。ベンダー主張を事実に格上げせず、未解決を未解決のまま予算に織り込む。三つの立場は、対象の名指しと条件の束ねと深さに応じた言葉という同じ作法を共有する。

\begin{themeoverview}[層の要点対照]
表現は短縮の設計、目的は分岐の選択、条件は遵守と品質の分離、制御は信号と標本の分離、編集は生成と保存の分離、時間は短距離と長距離の分離、実行は数値のバインド、評価は指標の非交換性、収束は未決の問いである。各層の要点は独立した教訓であると同時に、隣の層への問いかけでもある。
\end{themeoverview}
\end{multicols}

\begin{center}
\small
\begin{tabularx}{\textwidth}{@{}p{0.16\textwidth}p{0.34\textwidth}X@{}}
\toprule
場面 & 問うこと & 確かめること \\
\midrule
表現・目的 & 何を短縮し何で生成するか & 設定バインドと骨格・目的・手続きの分離 \\
条件・制御 & 遵守か品質か、信号か標本か & 学習時・推論時・参照の区別 \\
編集・時間 & 新規か保存か、短区間か長尺か & 局所性・参照制約・永続の分離 \\
実行・評価 & バインドなき比較を避ける & 版・装置・予算・母集団の添付 \\
収束・受容 & 未決を閉じない & lifecycleと27件の算術の保持 \\
\bottomrule
\end{tabularx}
\end{center}
\autocite{btd002,btd023,btd032,btd039,btd049,btd058,btd066,btd075,btd131,btd078,btd087,btd096,btd101,btd139}

\begin{claimboundary}[本巻の読解上の境界]
本巻は順位付けも推奨も行わない。要旨・断片・頁限定の事実はその水準でのみ用いる。閉鎖製品は能力・workflow・lifecycleの記録に限定し、構成・学習・符号化の推論は行わない。Xは受容・配備・失敗の証拠に限定し、技術権威にしない。条件の異なる数値の横断比較は行わない。未解決を閉じないことが、本巻の読みの作法である。
\end{claimboundary}


\clearpage
\onecolumn
\printbibliography[title={References}]
\end{document}
